% Étude et représentation graphique d'une fonction — Mathématiques 1ère TI — Cours signé OnBuch+
% Programme officiel MINESEC. Compiler avec : tectonic N07-representation-graphique-fonctions.tex
\newcommand{\DOCMATIERE}{Mathématiques}
\newcommand{\DOCNIVEAU}{1ère TI}
\newcommand{\DOCMODULE}{Relations et opérations fondamentales dans l'ensemble des nombres réels}
\newcommand{\DOCLECON}{N07}
\newcommand{\DOCTITRE}{Étude et représentation graphique d'une fonction}
% OnBuch+ — préambule « cours signés » ENRICHI (compilé avec tectonic / XeLaTeX)
% Système visuel « Pop » d'OnBuch+ : fond crème (jamais blanc), bordures encre
% épaisses, ombres « dures » décalées, titres Archivo Black en capitales.
% Version MATHÉMATIQUES : graphes (pgfplots/tikz), boîtes pédagogiques
% (définition, propriété, méthode, attention, le savais-tu, exercice résolu,
% exercice, TP, bilan), page de garde et signature OnBuch+ ancrée en bas.
%
% Macros attendues dans main.tex AVANT \input{preamble} :
%   \DOCMATIERE  \DOCNIVEAU  \DOCMODULE  \DOCLECON (numéro)  \DOCTITRE
\documentclass[11pt]{article}

\usepackage{fontspec}
\usepackage[french]{babel}
\usepackage[a4paper,margin=2cm,top=3.4cm,bottom=2.8cm,headheight=1.4cm,footskip=1.3cm]{geometry}
\usepackage{amsmath,amssymb,amsfonts}
\usepackage{mathtools} % \xmapsto, \coloneqq... souvent utilisés par les modèles
\usepackage{cancel} % simplifications barrées
% \abs{z} (module, valeur absolue) et \norm{u} (norme), utilisés par les modèles
\providecommand{\abs}{}\let\abs\relax\DeclarePairedDelimiter{\abs}{\lvert}{\rvert}
\providecommand{\norm}{}\let\norm\relax\DeclarePairedDelimiter{\norm}{\lVert}{\rVert}
\usepackage[table]{xcolor} % [table] : fournit aussi \rowcolors (alternance de lignes)
\usepackage{pagecolor}
\usepackage{tikz}
\usetikzlibrary{arrows.meta,positioning,calc,shapes.geometric,shapes.misc,shapes.arrows,decorations.pathmorphing,decorations.pathreplacing,decorations.markings,patterns,shadows,mindmap,trees,fit,backgrounds,matrix,intersections,angles,quotes}
\usepackage{pgfplots}
\usepackage{pgf-pie} % diagrammes circulaires \pie (statistiques)
\pgfplotsset{compat=1.18}
\usepgfplotslibrary{fillbetween,groupplots} % aires entre courbes (calcul intégral)
\usepackage{enumitem}
\usepackage{fancyhdr}
% [most] charge toutes les bibliothèques tcolorbox (listings, minted, raster,
% documentation...) et gonfle inutilement la mémoire TeX sur un document long
% (~300 boîtes) : on ne charge que ce qui est réellement utilisé.
\usepackage[skins,breakable]{tcolorbox}
\usepackage{graphicx}
\usepackage{colortbl}
\usepackage{booktabs}
\usepackage{tabularx}
\usepackage{diagbox} % case d'en-tête diagonale des tableaux à double entrée
% Colonnes extensibles centrée / gauche / droite (C, L, R), souvent utilisées par les modèles
\newcolumntype{C}{>{\centering\arraybackslash}X}
\newcolumntype{L}{>{\raggedright\arraybackslash}X}
\newcolumntype{R}{>{\raggedleft\arraybackslash}X}
\usepackage{array}
\usepackage{multicol}
\usepackage{siunitx}
% parse-numbers=false : accepte aussi \SI{2,5 \times 10^{-3}}{...}
\sisetup{locale=FR,output-decimal-marker={,},group-minimum-digits=4,parse-numbers=false}
\DeclareSIUnit\ohm{\ensuremath{\symup{\Omega}}} % U+2126 absent de la police
\usepackage{qrcode}
% \degree : commande fréquemment utilisée par les modèles pour un angle ou une
% température en dehors de \SI{}{} (non fournie par les paquets chargés).
\providecommand{\degree}{^{\circ}}
% \qed (fin de démonstration), utilisé par les modèles sans amsthm
\providecommand{\qed}{\hfill$\square$}
% \barre : double barre des valeurs interdites dans un tableau de variations (tkz-tab)
\providecommand{\barre}{\ensuremath{\|}}
% environnement proof (amsthm non chargé) utilisé par les modèles
\ifdefined\proof\else\newenvironment{proof}[1][Démonstration]{\par\noindent\textit{#1.}\ }{\qed\par}\fi
\usepackage{lastpage}
\usepackage{hyperref}
\hypersetup{hidelinks}

% ============================================================
% Palette « Pop » OnBuch+ — mêmes hex que la classe P (plus_ui.dart)
% ============================================================
\definecolor{poporange}{HTML}{F59321}
\definecolor{popdark}{HTML}{E07A0C}
\definecolor{popcream}{HTML}{FFF6E8}
\definecolor{popink}{HTML}{1C1714}
\definecolor{poppurple}{HTML}{7A5AE0}
\definecolor{popgreen}{HTML}{2E9E6A}
\definecolor{poppink}{HTML}{E0457B}
\definecolor{popblue}{HTML}{2F7DE1}
\definecolor{popgold}{HTML}{C98A12}
\definecolor{popmuted}{HTML}{8A7F76}
\definecolor{popline}{HTML}{EADFD2}
\definecolor{popgray}{HTML}{8A7F76}
\colorlet{popgrey}{popgray}
% Alias tolérants (le modèle nomme parfois les couleurs différemment)
\colorlet{popwhite}{white}
\colorlet{popblack}{popink}
\colorlet{popred}{poppink}
\colorlet{popyellow}{popgold}
% Teintes claires pour fonds de tableaux / zones
\colorlet{popblueL}{popblue!10!white}
\colorlet{poporangeL}{poporange!14!white}
\colorlet{popgreenL}{popgreen!12!white}
\colorlet{poppurpleL}{poppurple!12!white}
\colorlet{poppinkL}{poppink!12!white}
\colorlet{popgoldL}{popgold!14!white}

\pagecolor{popcream}

% ============================================================
% Typographie
% ============================================================
\defaultfontfeatures{Ligatures=TeX}
\setmainfont{PlusJakartaSans-Medium.ttf}[Path=../../../fonts/,BoldFont=PlusJakartaSans-Bold.ttf]
% Maths en sans-serif (Fira Math) avec chiffres et lettres droites de Plus
% Jakarta, pour que les formules se fondent dans le texte.
\usepackage[math-style=ISO,bold-style=ISO]{unicode-math}
\setmathfont{FiraMath-Regular.otf}[Path=../../../fonts/]
\setmathfont{PlusJakartaSans-Medium.ttf}[Path=../../../fonts/,range={up/num,up/{Latin,latin},\mathup}]
\usepackage{icomma} % virgule décimale sans espace en mode math
% Caractères mathématiques Unicode absents des polices de texte (Plus Jakarta,
% Archivo Black) : rendus par leur équivalent mathématique, en texte comme en
% formule — sinon ils s'affichent en carré vide (ex. « Arithmétique dans ℤ »).
\usepackage{newunicodechar}
\newunicodechar{ℝ}{\ensuremath{\mathbb{R}}}
\newunicodechar{ℤ}{\ensuremath{\mathbb{Z}}}
\newunicodechar{ℂ}{\ensuremath{\mathbb{C}}}
\newunicodechar{ℕ}{\ensuremath{\mathbb{N}}}
\newunicodechar{ℚ}{\ensuremath{\mathbb{Q}}}
\newunicodechar{𝔻}{\ensuremath{\mathbb{D}}}
\newunicodechar{α}{\ensuremath{\alpha}}
\newunicodechar{β}{\ensuremath{\beta}}
\newunicodechar{γ}{\ensuremath{\gamma}}
\newunicodechar{δ}{\ensuremath{\delta}}
\newunicodechar{ε}{\ensuremath{\varepsilon}}
\newunicodechar{θ}{\ensuremath{\theta}}
\newunicodechar{λ}{\ensuremath{\lambda}}
\newunicodechar{μ}{\ensuremath{\mu}}
\newunicodechar{σ}{\ensuremath{\sigma}}
\newunicodechar{τ}{\ensuremath{\tau}}
\newunicodechar{φ}{\ensuremath{\varphi}}
\newunicodechar{ψ}{\ensuremath{\psi}}
\newunicodechar{ω}{\ensuremath{\omega}}
\newunicodechar{Σ}{\ensuremath{\Sigma}}
% majuscules produites par \MakeUppercase dans les titres (α-aminés -> Α)
\newunicodechar{Α}{\ensuremath{\alpha}}
\newunicodechar{Β}{\ensuremath{\beta}}
\newunicodechar{Γ}{\ensuremath{\Gamma}}
\newunicodechar{Δ}{\ensuremath{\Delta}}
\newunicodechar{Ε}{\ensuremath{\varepsilon}}
\newunicodechar{Θ}{\ensuremath{\Theta}}
\newunicodechar{Λ}{\ensuremath{\Lambda}}
\newunicodechar{Μ}{\ensuremath{\mu}}
\newunicodechar{Τ}{\ensuremath{\tau}}
\newunicodechar{Φ}{\ensuremath{\Phi}}
\newunicodechar{Ψ}{\ensuremath{\Psi}}
\newunicodechar{Ω}{\ensuremath{\Omega}}
\newunicodechar{↦}{\ensuremath{\mapsto}}
\newunicodechar{∈}{\ensuremath{\in}}
\newunicodechar{∉}{\ensuremath{\notin}}
\newunicodechar{∧}{\ensuremath{\wedge}}
\newunicodechar{⇔}{\ensuremath{\Leftrightarrow}}
\newunicodechar{⟺}{\ensuremath{\Longleftrightarrow}}
\newunicodechar{⇒}{\ensuremath{\Rightarrow}}
\newunicodechar{⟹}{\ensuremath{\Longrightarrow}}
\newunicodechar{∘}{\ensuremath{\circ}}
\newunicodechar{∪}{\ensuremath{\cup}}
\newunicodechar{∩}{\ensuremath{\cap}}
\newunicodechar{≡}{\ensuremath{\equiv}}
\newunicodechar{⊂}{\ensuremath{\subset}}
\newunicodechar{⊥}{\ensuremath{\perp}}
\newunicodechar{∃}{\ensuremath{\exists}}
\newunicodechar{∀}{\ensuremath{\forall}}
\newunicodechar{∛}{\ensuremath{\sqrt[3]{\,}}}
\newunicodechar{∜}{\ensuremath{\sqrt[4]{\,}}}
\newunicodechar{⃗}{\ensuremath{{}^{\to}}}
\newunicodechar{ⁿ}{\ifmmode^{n}\else\textsuperscript{\ensuremath{n}}\fi}
\newunicodechar{ˣ}{\ifmmode^{x}\else\textsuperscript{\ensuremath{x}}\fi}
\newunicodechar{ᵇ}{\ifmmode^{b}\else\textsuperscript{\ensuremath{b}}\fi}
\newunicodechar{ʳ}{\ifmmode^{r}\else\textsuperscript{\ensuremath{r}}\fi}
\newunicodechar{ᵘ}{\ifmmode^{u}\else\textsuperscript{\ensuremath{u}}\fi}
\newunicodechar{ʸ}{\ifmmode^{y}\else\textsuperscript{\ensuremath{y}}\fi}
\newunicodechar{ᵃ}{\ifmmode^{a}\else\textsuperscript{\ensuremath{a}}\fi}
\newunicodechar{ᵉ}{\ifmmode^{e}\else\textsuperscript{\ensuremath{e}}\fi}
\newunicodechar{ᵏ}{\ifmmode^{k}\else\textsuperscript{\ensuremath{k}}\fi}
\newunicodechar{ᵖ}{\ifmmode^{p}\else\textsuperscript{\ensuremath{p}}\fi}
\newunicodechar{ᴺ}{\ifmmode^{N}\else\textsuperscript{\ensuremath{N}}\fi}
\newunicodechar{ᶜ}{\ifmmode^{c}\else\textsuperscript{\ensuremath{c}}\fi}
\newunicodechar{ʰ}{\ifmmode^{h}\else\textsuperscript{\ensuremath{h}}\fi}
\newunicodechar{ᵠ}{\ifmmode^{q}\else\textsuperscript{\ensuremath{q}}\fi}
\newunicodechar{ₙ}{\ifmmode_{n}\else\textsubscript{\ensuremath{n}}\fi}
\newunicodechar{ₐ}{\ifmmode_{a}\else\textsubscript{\ensuremath{a}}\fi}
\newunicodechar{ᵢ}{\ifmmode_{i}\else\textsubscript{\ensuremath{i}}\fi}
\newunicodechar{ᵣ}{\ifmmode_{r}\else\textsubscript{\ensuremath{r}}\fi}
\newunicodechar{ₖ}{\ifmmode_{k}\else\textsubscript{\ensuremath{k}}\fi}
\newunicodechar{ᵦ}{\ifmmode_{\beta}\else\textsubscript{\ensuremath{\beta}}\fi}
\newunicodechar{ₓ}{\ifmmode_{x}\else\textsubscript{\ensuremath{x}}\fi}
\newfontfamily\popdisplay{ArchivoBlack-Regular.ttf}[Path=../../../fonts/]
\newfontfamily\poplogo{SpaceGrotesk-Bold.ttf}[Path=../../../fonts/]
\newfontfamily\popsemi{PlusJakartaSans-SemiBold.ttf}[Path=../../../fonts/]
\newfontfamily\popxbold{PlusJakartaSans-ExtraBold.ttf}[Path=../../../fonts/]
\newfontfamily\popxboldls{PlusJakartaSans-ExtraBold.ttf}[Path=../../../fonts/,LetterSpace=3.0]

\setlist[enumerate]{leftmargin=1.7em,itemsep=5pt,topsep=4pt}
\setlist[itemize]{leftmargin=1.7em,itemsep=4pt,topsep=4pt,label={\color{poporange}\textbullet}}
\setlength{\parindent}{0pt}
\setlength{\parskip}{5pt}
\renewcommand{\arraystretch}{1.25}

% pgfplots : style Pop par défaut
\pgfplotsset{
  every axis/.append style={axis line style={popink,line width=1pt},tick style={popink},
    grid style={popline,line width=0.5pt},label style={font=\small},tick label style={font=\footnotesize}},
  popcurve/.style={poporange,line width=1.8pt,no markers,smooth},
}
% Même style disponible dans un \draw TikZ simple (hors pgfplots)
\tikzset{popcurve/.style={poporange,line width=1.8pt}}

% ============================================================
% Pill / squiggle
% ============================================================
\newcommand{\poppill}[3]{%
  \tikz[baseline=-0.55ex]\node[draw=popink,line width=1.2pt,rounded corners=2.6pt,%
    fill=#2,inner xsep=6.5pt,inner ysep=3pt]{%
    \popxboldls\fontsize{7}{7}\selectfont\color{#3}\MakeUppercase{#1}};%
}
\newcommand{\popsquiggle}[1]{%
  \tikz[baseline=-0.4ex]{
    \draw[#1,line width=2.6pt,line cap=round]
      plot [smooth] coordinates {(0,0.09) (0.34,0.01) (0.68,0.21) (1.02,0.01) (1.36,0.10)};
  }%
}
% Mot-clé surligné (marqueur orange sous le texte)
\newcommand{\cle}[1]{{\popxbold\color{popdark}#1}}

% ============================================================
% En-tête / pied
% ============================================================
\pagestyle{fancy}
\fancyhf{}
\renewcommand{\headrulewidth}{0pt}
\renewcommand{\footrulewidth}{0pt}
\lhead{\raisebox{-0.15ex}{\poplogo\fontsize{13}{13}\selectfont\color{popink}OnBuch\color{poporange}+}}
\rhead{\poppill{\DOCMATIERE\ \textbullet\ \DOCNIVEAU\ \textbullet\ Leçon \DOCLECON}{white}{popink}}
\lfoot{\popxbold\fontsize{7.6}{7.6}\selectfont\color{popmuted}\MakeUppercase{Cours signé OnBuch+}\ \textbullet\ {\popsemi\DOCTITRE}}
\rfoot{\tikz[baseline=-0.6ex]\node[rounded corners=8.5pt,fill=popink,minimum height=17pt,minimum width=17pt,inner xsep=5pt,inner sep=0]{\popxbold\fontsize{8}{8}\selectfont\color{popcream}\thepage\,/\,\pageref*{LastPage}};}

% ============================================================
% Titres
% ============================================================
\newcommand{\chaptitre}[2]{%
  \begin{tcolorbox}[enhanced,colback=poporange,colframe=popink,boxrule=2.6pt,arc=11pt,
      drop shadow=popink,left=15pt,right=15pt,top=12pt,bottom=11pt,before skip=3pt,after skip=18pt]
    \poppill{#2}{popink}{popcream}\par\vspace{9pt}
    {\raggedright\hyphenpenalty=10000\exhyphenpenalty=10000\popdisplay\fontsize{22}{24}\selectfont\color{popcream}\MakeUppercase{#1}\par}
  \end{tcolorbox}
}
% Section numérotée : pastille orange avec le numéro + titre Archivo
\newcounter{popsec}
\newcommand{\coursec}[1]{%
  \stepcounter{popsec}\setcounter{popsub}{0}%
  \par\vspace{16pt}\noindent
  \begin{minipage}{\linewidth}
  \tikz[baseline=-0.7ex]\node[draw=popink,line width=1.6pt,fill=poporange,rounded corners=4pt,minimum size=22pt,inner sep=0]{\popdisplay\fontsize{12}{12}\selectfont\color{popcream}\thepopsec};%
  \hspace{8pt}{\popdisplay\fontsize{15}{17}\selectfont\color{popink}\MakeUppercase{#1}}\par
  \vspace{4pt}\noindent{\color{poporange}\rule{36pt}{3.6pt}}
  \end{minipage}\par\nopagebreak\vspace{8pt}
}
\newcounter{popsub}
\newcommand{\courssub}[1]{%
  \stepcounter{popsub}%
  \par\vspace{9pt}\noindent{\popxbold\fontsize{11.5}{13}\selectfont\color{popink}{\color{poporange}\thepopsec.\thepopsub}\hspace{6pt}#1}\par\nopagebreak\vspace{4pt}
}

% ============================================================
% Boîtes pédagogiques — même recette : fond blanc, bordure épaisse colorée,
% ombre dure encre, badge plein. Titre optionnel [#1].
% ============================================================
\tcbset{popbase/.style={breakable,enhanced,colback=white,boxrule=2.3pt,arc=9pt,
  shadow={3pt}{-3pt}{0pt}{popink},left=14pt,right=14pt,top=11pt,bottom=11pt,before skip=11pt,after skip=15pt,
  fontupper=\popsemi\fontsize{10.6}{14.5}\selectfont\color{popink},
  fontlower=\popsemi\fontsize{10.6}{14.5}\selectfont\color{popink}}}
% Coche dessinée (le glyphe ✓ n'existe pas dans les polices du document)
\newcommand{\popcheck}[1]{\tikz[baseline=-0.5ex]\draw[#1,line width=1.6pt,line cap=round,line join=round](-0.13,0.0)--(-0.04,-0.09)--(0.13,0.1);}
\newcommand{\popboxhead}[3]{\poppill{#1}{#2}{white}\ifx\relax#3\relax\else\hspace{6pt}{\popxbold\fontsize{10.6}{12}\selectfont\color{popink}#3}\fi\par\vspace{7pt}}

\newtcolorbox{aretenir}[1][]{popbase,colframe=popgreen,before upper={\popboxhead{À retenir}{popgreen}{#1}}}
\newtcolorbox{exemplebox}[1][]{popbase,colframe=poporange,before upper={\popboxhead{Exemple}{poporange}{#1}}}
\newtcolorbox{definition}[1][]{popbase,colframe=popblue,before upper={\popboxhead{Définition}{popblue}{#1}}}
\newtcolorbox{propriete}[1][]{popbase,colframe=poppurple,before upper={\popboxhead{Propriété}{poppurple}{#1}}}
\newtcolorbox{methode}[1][]{popbase,colframe=popgold,before upper={\popboxhead{Méthode}{popgold}{#1}}}
\newtcolorbox{attention}[1][]{popbase,colframe=poppink,before upper={\popboxhead{Attention}{poppink}{#1}}}
\newtcolorbox{experience}[1][]{popbase,colframe=popblue,colback=popblueL,before upper={\popboxhead{Expérience}{popblue}{#1}}}
% « Le savais-tu ? » : carte violette pleine (contexte, culture, Cameroun)
\newtcolorbox{savaistu}[1][]{popbase,colback=poppurple,colframe=popink,
  fontupper=\popsemi\fontsize{10.4}{14.2}\selectfont\color{white},
  before upper={\poppill{Le savais-tu\,?}{white}{poppurple}\ifx\relax#1\relax\else\hspace{6pt}{\popxbold\color{white}#1}\fi\par\vspace{7pt}}}
% Exercice résolu : énoncé en haut, \tcblower puis la solution sur fond crème
\newcounter{popexo}\newcounter{popexr}
\newtcolorbox{exoresolu}[1][]{popbase,colframe=poporange,colbacklower=popcream,
  segmentation style={popink,line width=1.2pt,dashed},
  before upper={\stepcounter{popexr}\popboxhead{Exercice résolu \thepopexr}{poporange}{#1}},
  before lower={\poppill{Solution}{popink}{popcream}\par\vspace{6pt}}}
% Exercice (non résolu en place) — la correction est dans la partie Corrigés
\newtcolorbox{exercice}[1][]{popbase,colframe=popink,
  before upper={\stepcounter{popexo}\popboxhead{Exercice \thepopexo}{popink}{#1}}}
% Corrigé
\newtcolorbox{corrige}[1][]{popbase,colframe=popgreen,colback=popgreenL,
  before upper={\popboxhead{Corrigé}{popgreen}{#1}}}
% Objectifs / prérequis / bilan
\newtcolorbox{objectifs}{popbase,colframe=popink,colback=poporangeL,
  before upper={\popboxhead{Objectifs de la leçon}{popink}{}}}
\newtcolorbox{prerequis}{popbase,colframe=popink,colback=popblueL,
  before upper={\popboxhead{Prérequis}{popblue}{}}}
\newtcolorbox{bilan}{popbase,colframe=popink,colback=white,boxrule=2.8pt,
  before upper={\popboxhead{Fiche bilan}{poporange}{L'essentiel à connaître pour l'examen}}}
% Figure encadrée avec légende
\newtcolorbox{popfigure}[1][]{enhanced,breakable=false,colback=white,colframe=popline,boxrule=1.4pt,arc=8pt,
  left=10pt,right=10pt,top=10pt,bottom=8pt,before skip=10pt,after skip=12pt,halign=center,
  lower separated=false,halign lower=center,
  fontlower=\popsemi\fontsize{9}{11.5}\selectfont\color{popmuted},
  #1}
\newcommand{\legende}[1]{\par\vspace{6pt}{\popsemi\fontsize{9}{11.5}\selectfont\color{popmuted}#1}}

% ============================================================
% Signature OnBuch+
% ============================================================
\newcommand{\poplogoinline}[1]{{\poplogo\fontsize{#1}{#1}\selectfont\color{popink}OnBuch\color{poporange}+}}
\newcommand{\signatureonbuch}{%
  \begin{tcolorbox}[enhanced,colback=popink,colframe=popink,boxrule=2.2pt,arc=10pt,
      drop shadow=poporange,left=14pt,right=14pt,top=10pt,bottom=10pt,before skip=0pt,after skip=0pt]
    \tikz[baseline=-0.7ex]\node[circle,fill=popgreen,draw=popcream,line width=1.3pt,minimum size=22pt,inner sep=0]{\popcheck{white}};%
    \hspace{9pt}\begin{minipage}[c]{0.62\linewidth}
      {\popxbold\fontsize{12}{13}\selectfont\color{popcream}Signé\ \textbullet\ {\poplogo OnBuch\color{poporange}+}}\par\vspace{2pt}
      {\raggedright\popsemi\fontsize{8}{10.5}\selectfont\color{popline}\MakeUppercase{Équipe pédagogique OnBuch+}\\\MakeUppercase{Conforme au programme officiel MINESEC}\par}
    \end{minipage}\hfill
    {\poplogo\fontsize{22}{22}\selectfont\color{popcream}OnBuch\color{poporange}+}
  \end{tcolorbox}%
}

% ============================================================
% Page de garde
% #1 titre · #2 matière · #3 niveau · #4 module · #5 n° leçon · #6 durée ·
% #7 description / objectifs (texte ou itemize)
% ============================================================
\newcommand{\pagegarde}[7]{%
  \thispagestyle{empty}
  \newgeometry{a4paper,margin=1.7cm,top=1.6cm,bottom=1.4cm}%
  \begin{tikzpicture}[remember picture,overlay]
    \fill[poporange!18] ([xshift=-3.2cm,yshift=-3.6cm]current page.north east) circle (4.6cm);
    \fill[poppurple!14] ([xshift=2.4cm,yshift=7.6cm]current page.south west) circle (3.8cm);
  \end{tikzpicture}%
  {\poplogo\fontsize{30}{30}\selectfont\color{popink}OnBuch\color{poporange}+}\hfill
  \raisebox{0.6ex}{\poppill{Cours signé \textbullet\ Premium}{popink}{popcream}}\par
  \vspace{1pt}\popsquiggle{poppurple}\par
  \vspace{8pt}
  \begin{tcolorbox}[enhanced,colback=poporange,colframe=popink,boxrule=2.8pt,arc=13pt,
      drop shadow=popink,left=18pt,right=18pt,top=12pt,bottom=13pt,after skip=10pt]
    \poppill{#2 \textbullet\ #3}{popink}{popcream}\hspace{5pt}\poppill{#4}{white}{popink}\par\vspace{8pt}
    {\popdisplay\fontsize{14}{14}\selectfont\color{popink}LEÇON #5}\par\vspace{4pt}
    {\raggedright\hyphenpenalty=10000\exhyphenpenalty=10000\popdisplay\fontsize{25}{27}\selectfont\color{popcream}\MakeUppercase{#1}\par}
    \vspace{8pt}
    {\popxbold\fontsize{9.5}{9.5}\selectfont\color{popink}\MakeUppercase{Durée conseillée : #6}}
  \end{tcolorbox}
  \begin{tcolorbox}[enhanced,colback=white,colframe=popink,boxrule=2.2pt,arc=10pt,
      drop shadow=popink,left=16pt,right=16pt,top=10pt,bottom=10pt,after skip=8pt]
    \poppill{Dans cette leçon}{poppurple}{white}\par\vspace{5pt}
    \popsemi\fontsize{9.3}{12.6}\selectfont\color{popink}#7
  \end{tcolorbox}
  \noindent\begin{minipage}[t]{0.487\linewidth}
    \begin{tcolorbox}[enhanced,colback=popgreen,colframe=popink,boxrule=2.2pt,arc=10pt,
        drop shadow=popink,left=11pt,right=11pt,top=10pt,bottom=10pt,height=3.1cm,valign=center]
      \tcbox[colback=white,colframe=popink,boxrule=1.3pt,arc=5pt,boxsep=0pt,nobeforeafter,
        left=3pt,right=3pt,top=3pt,bottom=3pt,valign=center]{\qrcode[height=1.9cm]{https://play.google.com/store/apps/details?id=com.luvvix.onbuch&hl=fr}}%
      \hspace{8pt}\begin{minipage}[c]{0.5\linewidth}
        {\popdisplay\fontsize{11}{12.5}\selectfont\color{white}\MakeUppercase{Télécharge OnBuch}}\par\vspace{4pt}
        {\raggedright\popsemi\fontsize{8.4}{11}\selectfont\color{white}Gratuit \textbullet\ Google Play\\Tuteur IA, annales, concours, résultats\par}
      \end{minipage}
    \end{tcolorbox}
  \end{minipage}\hfill
  \begin{minipage}[t]{0.487\linewidth}
    \begin{tcolorbox}[enhanced,colback=poppurple,colframe=popink,boxrule=2.2pt,arc=10pt,
        drop shadow=popink,left=11pt,right=11pt,top=10pt,bottom=10pt,height=3.1cm,valign=center]
      \tikz[baseline=-0.7ex]\node[circle,fill=white,draw=popink,line width=1.3pt,minimum size=1.6cm,inner sep=0]{\popdisplay\fontsize{24}{24}\selectfont\color{poppurple}+};%
      \hspace{8pt}\begin{minipage}[c]{0.6\linewidth}
        {\popdisplay\fontsize{11}{12.5}\selectfont\color{white}\MakeUppercase{Passe à OnBuch+}}\par\vspace{4pt}
        {\raggedright\popsemi\fontsize{8.4}{11}\selectfont\color{white}Tous les cours signés, Tuteur IA illimité, fiches PDF — onglet OnBuch+ de l'app\par}
      \end{minipage}
    \end{tcolorbox}
  \end{minipage}
  \vspace{8pt}
  \popcontenu
  \vfill
  \signatureonbuch
  \restoregeometry
  \newpage
}

% Rangée « Ce cours contient » (page de garde)
\newcommand{\poptile}[3]{% #1 couleur #2 gros texte #3 légende
  \begin{tcolorbox}[enhanced,colback=#1,colframe=popink,boxrule=2pt,arc=9pt,shadow={2.5pt}{-2.5pt}{0pt}{popink},
      left=6pt,right=6pt,top=9pt,bottom=9pt,halign=center,valign=center,height=1.75cm,width=\linewidth,nobeforeafter]
    {\popdisplay\fontsize{12.5}{13}\selectfont\color{white}\MakeUppercase{#2}}\par\vspace{3pt}
    {\centering\popsemi\fontsize{7.8}{9.5}\selectfont\color{white}#3\par}
  \end{tcolorbox}}
\newcommand{\popcontenu}{%
  \par\noindent{\popxboldls\fontsize{7.5}{7.5}\selectfont\color{popmuted}CE COURS SIGNÉ CONTIENT}\par\vspace{7pt}
  \noindent\begin{minipage}[t]{0.235\linewidth}\poptile{popblue}{Cours}{complet, illustré, pas à pas}\end{minipage}\hfill
  \begin{minipage}[t]{0.235\linewidth}\poptile{poporange}{Méthodes}{et exercices résolus}\end{minipage}\hfill
  \begin{minipage}[t]{0.235\linewidth}\poptile{poppink}{Exercices}{type examen + corrigés}\end{minipage}\hfill
  \begin{minipage}[t]{0.235\linewidth}\poptile{popgreen}{Fiche bilan}{l'essentiel à retenir}\end{minipage}\par}

% Sommaire
\newtcolorbox{sommaire}{enhanced,colback=white,colframe=popink,boxrule=2.2pt,arc=10pt,
  drop shadow=popink,left=18pt,right=18pt,top=13pt,bottom=13pt,before skip=6pt,after skip=20pt,
  before upper={\poppill{Sommaire}{poppurple}{white}\par\vspace{9pt}\popsemi\fontsize{10.6}{14.5}\selectfont\color{popink}}}

% Signature de fin de cours — poussée tout en bas de la dernière page
\newcommand{\findecours}{%
  \par\vfill\nopagebreak
  \begin{center}\popsquiggle{poporange}\end{center}\vspace{4pt}
  \signatureonbuch
}
\AtBeginDocument{\def\llbracket{\mathopen{[\mkern-3.5mu[}}\def\rrbracket{\mathclose{]\mkern-3.5mu]}}} % intervalles d'entiers (glyphe ⟦ absent de la police)
% Glyphes absents de Fira Math : redéfinis à partir de glyphes disponibles
\AtBeginDocument{%
  \def\setminus{\mathbin{\backslash}}\let\smallsetminus\setminus
  \def\vdots{\mathord{\vbox{\baselineskip3.2pt\lineskiplimit0pt\kern4pt\hbox{.}\hbox{.}\hbox{.}}}}%
  \def\ddots{\mathinner{\mkern1mu\raise6.5pt\hbox{.}\mkern2mu\raise3.6pt\hbox{.}\mkern2mu\raise0.7pt\hbox{.}\mkern1mu}}%
  \def\triangle{\mathord{\scalebox{0.9}{$\symup{\Delta}$}}}%
  \def\nabla{\mathord{\raisebox{\depth}{\scalebox{1}[-1]{$\symup{\Delta}$}}}}%
  \def\checkmark{\popcheck{popdark}}%
  \def\star{\mathbin{\ast}}\def\bigstar{\mathord{\ast}}%
  \def\diamond{\mathbin{\scalebox{0.6}{$\Diamond$}}}%
  \def\bigcup{\mathop{\mathchoice{\raisebox{-0.3em}{\scalebox{1.7}{$\cup$}}}{\scalebox{1.25}{$\cup$}}{\cup}{\cup}}\displaylimits}%
  \def\bigcap{\mathop{\mathchoice{\raisebox{-0.3em}{\scalebox{1.7}{$\cap$}}}{\scalebox{1.25}{$\cap$}}{\cap}{\cap}}\displaylimits}%
  \def\lesseqqgtr{\mathrel{\lessgtr}}\def\bigoplus{\mathop{\oplus}\displaylimits}%
}
\providecommand{\ssi}{si et seulement si} % abréviation fréquente
\AtBeginDocument{\providecommand{\wideparen}{\overparen}} % arc de cercle

%% FIGFIX-BEGIN v5 — correctifs globaux des figures (docs/onbuch-generation/tools/figfix)
\usepackage{adjustbox}\usepackage{etoolbox}\tcbuselibrary{raster}
% toute figure TikZ/pgfplots plus large que la ligne est réduite à la largeur disponible
\BeforeBeginEnvironment{tikzpicture}{\begin{adjustbox}{max width=\linewidth}}
\AfterEndEnvironment{tikzpicture}{\end{adjustbox}}
% légendes pgfplots lisibles par-dessus les courbes
\pgfplotsset{every axis/.append style={legend style={fill=white,fill opacity=0.92,text opacity=1,draw=black!15,inner sep=3pt}}}
% étiquettes d'arêtes (syntaxe edge["..."]) avec fond blanc
\tikzset{every edge quotes/.append style={fill=white,inner sep=1.2pt}}
%% FIGFIX-END
\begin{document}

\pagegarde{Étude et représentation graphique d'une fonction}{Mathématiques}{1ère TI}{Relations et opérations fondamentales dans l'ensemble des nombres réels}{N07}{6 h}{Cette leçon achève l'étude des fonctions en Première : elle apprend à mener une étude complète (domaine, limites, dérivée, variations, asymptotes, symétries) et à construire la représentation graphique des fonctions polynômes de degré au plus 3, homographiques et rationnelles du type (ax²+bx+c)/(dx+e). Elle montre aussi comment exploiter une courbe pour résoudre graphiquement équations et inéquations.
\vspace{4pt}
\begin{multicols}{2}\small
\begin{itemize}
\item À la fin de cette leçon, je sais déterminer par leurs équations les asymptotes parallèles aux axes à la courbe d'une fonction
\item montrer qu'une droite donnée est asymptote oblique à la courbe d'une fonction
\item étudier et représenter graphiquement une fonction homographique en exhibant son centre de symétrie
\item étudier et représenter graphiquement une fonction polynôme de degré inférieur ou égal à 3
\item étudier et représenter graphiquement une fonction du type x ↦ (ax²+bx+c)/(dx+e) avec a et d non nuls
\item démontrer qu'un point est centre de symétrie ou qu'une droite est axe de symétrie d'une courbe
\end{itemize}\end{multicols}}

\begin{sommaire}
\begin{multicols}{2}
\begin{enumerate}
\item Asymptotes parallèles aux axes
\item Asymptotes obliques
\item Éléments de symétrie d'une courbe
\item Étude des fonctions homographiques
\item Étude des fonctions polynômes de degré ≤ 3
\item Fonctions x ↦ (ax²+bx+c)/(dx+e) et exploitation graphique
\item Activité d'intégration
\item Méthodes et exercices résolus
\item Exercices
\item Corrigés des exercices
\item Fiche bilan
\end{enumerate}
\end{multicols}
\end{sommaire}

\begin{objectifs}
\begin{itemize}
\item À la fin de cette leçon, je sais déterminer par leurs équations les asymptotes parallèles aux axes à la courbe d'une fonction
\item montrer qu'une droite donnée est asymptote oblique à la courbe d'une fonction
\item étudier et représenter graphiquement une fonction homographique en exhibant son centre de symétrie
\item étudier et représenter graphiquement une fonction polynôme de degré inférieur ou égal à 3
\item étudier et représenter graphiquement une fonction du type x ↦ (ax²+bx+c)/(dx+e) avec a et d non nuls
\item démontrer qu'un point est centre de symétrie ou qu'une droite est axe de symétrie d'une courbe
\item exploiter la courbe d'une fonction pour résoudre graphiquement des équations et des inéquations
\item dresser un tableau de variations complet intégrant limites et asymptotes
\end{itemize}
\end{objectifs}

\begin{prerequis}
\begin{itemize}
\item Limites de fonctions (en un point, en +∞ et −∞) et opérations sur les limites
\item Dérivation : dérivées des fonctions usuelles, dérivée de u/v, signe de la dérivée et sens de variation
\item Fonctions homographiques vues en Seconde et équations du second degré
\item Équations de droites, parallélisme, lecture graphique de courbes
\item Parité d'une fonction (fonctions paires et impaires)
\end{itemize}
\end{prerequis}

\newpage

\coursec{Asymptotes parallèles aux axes}

Une coopérative de producteurs de cacao de la région du Centre vend sa récolte à un exportateur. Le coût total de production, en milliers de FCFA, est modélisé par
\[ C(x) = \frac{2x^2 + 3x + 1}{x + 1} \]
où $x$ désigne la quantité produite, en tonnes. Le trésorier de la coopérative se pose trois questions très concrètes : \emph{à partir de quelle quantité le coût se stabilise-t-il ? Existe-t-il une quantité « interdite » pour laquelle le modèle n'a plus de sens ? Comment prévoir le coût pour de très grandes productions sans tout calculer ?}

Ces questions semblent différentes, mais elles ont un point commun : elles demandent de comprendre le \cle{comportement} d'une fonction \emph{aux bords} de son domaine — près d'une valeur interdite, ou quand $x$ devient très grand. Or ce comportement se lit admirablement sur la courbe : celle-ci se rapproche parfois indéfiniment d'une droite, sans jamais la toucher (ou presque). Ces droites « guides » s'appellent des \cle{asymptotes}. Dans cette première section, nous étudions les asymptotes \emph{parallèles aux axes} : les asymptotes horizontales et les asymptotes verticales. À la fin de la leçon complète, tu sauras répondre à toutes les questions du trésorier.

\courssub{Rappels sur les limites}

Avant de définir les asymptotes, rappelons les deux types de limites qui les produisent. Ces rappels sont issus de la leçon sur les limites ; ils sont indispensables ici.

\begin{aretenir}[Les deux situations qui donnent une asymptote parallèle aux axes]
\begin{itemize}
\item \textbf{Limite finie à l'infini.} On dit que $f(x)$ tend vers le réel $b$ quand $x$ tend vers $+\infty$, et on écrit $\lim_{x \to +\infty} f(x) = b$, lorsque $f(x)$ devient aussi proche de $b$ qu'on veut, pourvu que $x$ soit assez grand. (Même idée en $-\infty$.)
\item \textbf{Limite infinie en un point.} On dit que $f(x)$ tend vers $+\infty$ quand $x$ tend vers $a$, et on écrit $\lim_{x \to a} f(x) = +\infty$, lorsque $f(x)$ devient aussi grand qu'on veut, pourvu que $x$ soit assez proche de $a$. On distingue souvent la limite à gauche ($x \to a$, $x < a$, notée $x \to a^-$) et la limite à droite ($x \to a^+$).
\end{itemize}
\end{aretenir}

\begin{exemplebox}[Deux limites de référence]
\begin{itemize}
\item $\lim_{x \to +\infty} \dfrac{1}{x} = 0$ : quand $x$ devient immense, $\frac{1}{x}$ devient minuscule. Limite \emph{finie} à l'infini.
\item $\lim_{x \to 0^+} \dfrac{1}{x} = +\infty$ : quand $x$ s'approche de $0$ par valeurs positives, $\frac{1}{x}$ explose. Limite \emph{infinie} en un point.
\end{itemize}
\end{exemplebox}

L'idée intuitive est simple : une \textbf{limite finie à l'infini} signifie que la courbe « s'aplatit » vers une hauteur fixe ; une \textbf{limite infinie en un point} signifie que la courbe « s'envole » verticalement près d'une certaine abscisse. Chacune de ces deux situations dessine une droite invisible que la courbe épouse de plus en plus : une asymptote.

\courssub{Asymptote horizontale}

\begin{definition}[Asymptote horizontale]
Soit $f$ une fonction définie sur un intervalle de la forme $[A\,;\,+\infty[$ (resp. $]-\infty\,;\,A]$) et $(\mathcal{C}_f)$ sa courbe représentative dans un repère.
Si
\[ \lim_{x \to +\infty} f(x) = b \quad \text{(resp. } \lim_{x \to -\infty} f(x) = b\text{)}, \]
avec $b$ un nombre \textbf{réel} (limite finie), alors la droite d'équation $\boxed{y = b}$ est appelée \cle{asymptote horizontale} à $(\mathcal{C}_f)$ en $+\infty$ (resp. en $-\infty$).
\end{definition}

Autrement dit : quand on s'éloigne vers la droite (ou vers la gauche) du graphique, la courbe se rapproche indéfiniment de la droite horizontale $y = b$. La distance verticale entre la courbe et la droite, qui vaut $|f(x) - b|$, tend vers $0$.

\begin{exemplebox}[La fonction inverse]
Pour $f(x) = \dfrac{1}{x}$, on a $\lim_{x \to +\infty} \dfrac{1}{x} = 0$ et $\lim_{x \to -\infty} \dfrac{1}{x} = 0$. La droite $y = 0$ (l'axe des abscisses) est donc asymptote horizontale à l'hyperbole, à la fois en $+\infty$ et en $-\infty$.
\end{exemplebox}

\begin{attention}[Asymptote horizontale : deux conditions indispensables]
Pour conclure à une asymptote horizontale, il faut vérifier \textbf{deux} choses : la limite est prise en $+\infty$ ou en $-\infty$, \emph{et} elle est \textbf{finie}. Si $\lim_{x \to +\infty} f(x) = +\infty$ (comme pour $f(x) = x^2$), il n'y a \textbf{pas} d'asymptote horizontale : la courbe s'envole, elle ne se stabilise pas. Ne confonds jamais « la courbe monte » avec « la courbe a une asymptote ».
\end{attention}

\courssub{Asymptote verticale}

\begin{definition}[Asymptote verticale]
Soit $f$ une fonction et $a$ un réel tel que $f$ soit définie au voisinage de $a$ (à gauche, à droite, ou des deux côtés), mais pas nécessairement en $a$.
Si l'une des limites suivantes est infinie :
\[ \lim_{x \to a^-} f(x) = \pm\infty \qquad \text{ou} \qquad \lim_{x \to a^+} f(x) = \pm\infty, \]
alors la droite d'équation $\boxed{x = a}$ est appelée \cle{asymptote verticale} à la courbe $(\mathcal{C}_f)$.
\end{definition}

Ici, c'est quand $x$ s'approche d'une valeur précise $a$ que la courbe « s'envole » vers le haut ou vers le bas. La droite verticale $x = a$ agit comme un mur que la courbe longe sans jamais franchir.

\begin{attention}[Une valeur interdite ne donne pas toujours une asymptote verticale]
C'est le piège classique du Probatoire ! Une valeur interdite $a$ (qui annule le dénominateur) ne donne une asymptote verticale \textbf{que si la limite en $a$ est infinie}. Contre-exemple : soit
\[ h(x) = \frac{x^2 - 1}{x - 1}, \quad \text{définie sur } \mathbb{R} \setminus \{1\}. \]
La valeur $x = 1$ est interdite. Pourtant, pour $x \neq 1$ :
\[ h(x) = \frac{(x-1)(x+1)}{x-1} = x + 1, \qquad \text{donc } \lim_{x \to 1} h(x) = 2, \]
limite \textbf{finie}. Il n'y a \textbf{pas} d'asymptote verticale en $x = 1$ : la courbe est la droite $y = x+1$ avec simplement un « trou » au point $(1\,;\,2)$. Moralité : \emph{il faut toujours calculer la limite avant de conclure}.
\end{attention}

\begin{popfigure}
\begin{center}
\begin{tikzpicture}[scale=0.72, every node/.style={scale=0.85}]
% Configuration 1 : limite à gauche = +infini
\begin{scope}[shift={(0,0)}]
\draw[->, popmuted] (-0.4,0) -- (3,0);
\draw[->, popmuted] (0,-0.4) -- (0,3);
\draw[dashed, poppink, thick] (2,-0.4) -- (2,3) node[above, poppink] {$x=a$};
\draw[popblue, very thick, domain=0.3:1.75, samples=60] plot (\x, {0.5/(2-\x) + 0.3});
\node[popdark] at (1.5,-0.9) {$\lim_{x \to a^-} f(x) = +\infty$};
\end{scope}
% Configuration 2 : limite à gauche = -infini
\begin{scope}[shift={(4.6,0)}]
\draw[->, popmuted] (-0.4,0) -- (3,0);
\draw[->, popmuted] (0,-3) -- (0,0.6);
\draw[dashed, poppink, thick] (2,-3) -- (2,0.6) node[above, poppink] {$x=a$};
\draw[popblue, very thick, domain=0.3:1.75, samples=60] plot (\x, {-0.5/(2-\x) - 0.3});
\node[popdark] at (1.5,-3.6) {$\lim_{x \to a^-} f(x) = -\infty$};
\end{scope}
% Configuration 3 : limite à droite = +infini
\begin{scope}[shift={(9.2,0)}]
\draw[->, popmuted] (-0.4,0) -- (3,0);
\draw[->, popmuted] (0,-0.4) -- (0,3);
\draw[dashed, poppink, thick] (0.8,-0.4) -- (0.8,3) node[above, poppink] {$x=a$};
\draw[popblue, very thick, domain=1.05:2.7, samples=60] plot (\x, {0.5/(\x-0.8) + 0.3});
\node[popdark] at (1.5,-0.9) {$\lim_{x \to a^+} f(x) = +\infty$};
\end{scope}
% Configuration 4 : limite à droite = -infini
\begin{scope}[shift={(13.8,0)}]
\draw[->, popmuted] (-0.4,0) -- (3,0);
\draw[->, popmuted] (0,-3) -- (0,0.6);
\draw[dashed, poppink, thick] (0.8,-3) -- (0.8,0.6) node[above, poppink] {$x=a$};
\draw[popblue, very thick, domain=1.05:2.7, samples=60] plot (\x, {-0.5/(\x-0.8) - 0.3});
\node[popdark] at (1.5,-3.6) {$\lim_{x \to a^+} f(x) = -\infty$};
\end{scope}
\end{tikzpicture}
\end{center}
\legende{Les quatre configurations d'une limite infinie en un point $a$ : dans chaque cas, la droite $x = a$ (en pointillés) est asymptote verticale à la courbe.}
\end{popfigure}

\courssub{Lien avec le domaine de définition}

Où faut-il chercher les asymptotes verticales ? La réponse tient en une phrase : \emph{aux bornes du domaine de définition}. Pour une \cle{fonction rationnelle} (quotient de deux polynômes), les candidats sont les \textbf{valeurs interdites}, c'est-à-dire les réels qui annulent le dénominateur.

\begin{methode}[Trouver les asymptotes parallèles aux axes]
Pour déterminer toutes les asymptotes parallèles aux axes de la courbe d'une fonction $f$ :
\begin{enumerate}
\item \textbf{Déterminer le domaine de définition} $D_f$. Pour une fonction rationnelle, résoudre « dénominateur $= 0$ » pour trouver les valeurs interdites.
\item \textbf{En chaque valeur interdite $a$} : calculer $\lim_{x \to a^-} f(x)$ et $\lim_{x \to a^+} f(x)$. Si l'une au moins est infinie, la droite $x = a$ est asymptote verticale. Si les deux sont finies et égales, il n'y a pas d'asymptote (simple « trou »).
\item \textbf{En $+\infty$ et en $-\infty$} (si le domaine le permet) : calculer $\lim_{x \to +\infty} f(x)$ et $\lim_{x \to -\infty} f(x)$. Chaque limite \textbf{finie} $b$ donne une asymptote horizontale $y = b$.
\item \textbf{Conclure proprement} en rédigeant : « la droite d'équation $x = a$ est asymptote verticale à $(\mathcal{C}_f)$ », « la droite d'équation $y = b$ est asymptote horizontale à $(\mathcal{C}_f)$ en $+\infty$ ».
\end{enumerate}
\end{methode}

\begin{exoresolu}[Asymptotes de $f(x) = \dfrac{2x+1}{x-3}$]
Montrer que les droites $x = 3$ et $y = 2$ sont asymptotes à la courbe de la fonction $f$ définie par $f(x) = \dfrac{2x+1}{x-3}$.
\tcblower
\textbf{Étape 1 : domaine de définition.} Le dénominateur s'annule pour $x - 3 = 0$, c'est-à-dire $x = 3$. Donc $D_f = \mathbb{R} \setminus \{3\} = ]-\infty\,;\,3[ \cup ]3\,;\,+\infty[$.

\textbf{Étape 2 : étude en la valeur interdite $x = 3$.} Pour $x$ proche de $3$, le numérateur vaut $2 \times 3 + 1 = 7 > 0$.
\begin{itemize}
\item Si $x \to 3^+$, alors $x - 3 \to 0^+$, donc $\lim_{x \to 3^+} f(x) = +\infty$.
\item Si $x \to 3^-$, alors $x - 3 \to 0^-$, donc $\lim_{x \to 3^-} f(x) = -\infty$.
\end{itemize}
Les limites sont infinies : la droite $\boxed{x = 3}$ est \textbf{asymptote verticale} à $(\mathcal{C}_f)$.

\textbf{Étape 3 : étude en $+\infty$ et $-\infty$.} Pour $x \neq 0$, factorisons par le terme dominant :
\[ f(x) = \frac{2x+1}{x-3} = \frac{x\left(2 + \frac{1}{x}\right)}{x\left(1 - \frac{3}{x}\right)} = \frac{2 + \frac{1}{x}}{1 - \frac{3}{x}}. \]
Or $\lim_{x \to \pm\infty} \dfrac{1}{x} = 0$ et $\lim_{x \to \pm\infty} \dfrac{3}{x} = 0$, donc
\[ \lim_{x \to +\infty} f(x) = \frac{2}{1} = 2 \qquad \text{et} \qquad \lim_{x \to -\infty} f(x) = 2. \]
Les limites sont finies et égales à $2$ : la droite $\boxed{y = 2}$ est \textbf{asymptote horizontale} à $(\mathcal{C}_f)$ en $+\infty$ et en $-\infty$.
\end{exoresolu}

\begin{popfigure}
\begin{center}
\begin{tikzpicture}
\begin{axis}[
    width=0.85\linewidth, height=8cm,
    axis lines=middle,
    xlabel={$x$}, ylabel={$y$},
    xmin=-4, xmax=10, ymin=-6, ymax=10,
    xtick={-2,0,2,3,4,6,8}, ytick={-4,-2,0,2,4,6,8},
    grid=both, grid style={popline!40},
    legend style={at={(0.03,0.97)}, anchor=north west, draw=popline}
]
\addplot[popcurve, domain=-4:2.6, samples=150] {(2*x+1)/(x-3)};
\addplot[popcurve, domain=3.4:10, samples=150] {(2*x+1)/(x-3)};
\addplot[dashed, thick, poppink] coordinates {(3,-6) (3,10)};
\addplot[dashed, thick, poppink, domain=-4:10, samples=2] {2};
\legend{$\mathcal{C}_f$, $x=3$, $y=2$}
\end{axis}
\end{tikzpicture}
\end{center}
\legende{Courbe de $f(x) = \dfrac{2x+1}{x-3}$ : la droite $x = 3$ est asymptote verticale, la droite $y = 2$ est asymptote horizontale (en pointillés rouges).}
\end{popfigure}

\courssub{Position relative de la courbe et de son asymptote horizontale}

Savoir que la droite $y = b$ est asymptote ne suffit pas toujours : on veut souvent préciser si la courbe s'en rapproche \emph{par au-dessus} ou \emph{par en dessous}. C'est une question de \textbf{signe}.

\begin{propriete}[Position relative de $(\mathcal{C}_f)$ et de l'asymptote $y = b$]
Supposons que la droite $y = b$ soit asymptote horizontale à $(\mathcal{C}_f)$ en $+\infty$. Alors, sur un intervalle $[A\,;\,+\infty[$ :
\begin{itemize}
\item si $f(x) - b > 0$, la courbe est \textbf{au-dessus} de son asymptote ;
\item si $f(x) - b < 0$, la courbe est \textbf{en dessous} de son asymptote ;
\item si $f(x) - b$ change de signe, la courbe \textbf{coupe} son asymptote (ce qui est tout à fait possible !).
\end{itemize}
\end{propriete}

\begin{exemplebox}[Position de la courbe de $f(x) = \dfrac{2x+1}{x-3}$ par rapport à $y = 2$]
Calculons la différence :
\[ f(x) - 2 = \frac{2x+1}{x-3} - 2 = \frac{2x+1 - 2(x-3)}{x-3} = \frac{7}{x-3}. \]
Le signe de $f(x) - 2$ est celui de $x - 3$ :
\begin{itemize}
\item sur $]3\,;\,+\infty[$, $f(x) - 2 > 0$ : la courbe est \textbf{au-dessus} de l'asymptote $y = 2$ ;
\item sur $]-\infty\,;\,3[$, $f(x) - 2 < 0$ : la courbe est \textbf{en dessous} de l'asymptote.
\end{itemize}
On retrouve exactement ce que montre la figure ci-dessus.
\end{exemplebox}

\begin{attention}[Une asymptote n'est pas une barrière interdite]
Le mot « asymptote » évoque une droite que la courbe ne touche jamais — c'est vrai \emph{localement}, à l'infini. Mais rien n'empêche la courbe de \textbf{couper} son asymptote horizontale pour des valeurs finies de $x$ ! Par exemple, la fonction $x \mapsto \dfrac{\sin x}{x}$ coupe son asymptote $y = 0$ une infinité de fois. L'asymptote décrit un comportement \emph{quand $x$ tend vers l'infini}, pas une interdiction de contact.
\end{attention}

\courssub{Lecture graphique d'asymptotes}

Au Probatoire, on te donnera parfois une courbe déjà tracée et on te demandera d'en \emph{lire} les asymptotes. La démarche est l'inverse de la précédente.

\begin{methode}[Lire une asymptote sur un graphique]
\begin{enumerate}
\item \textbf{Asymptote verticale} : repérer une droite verticale le long de laquelle la courbe « s'envole » vers $+\infty$ ou $-\infty$. Lire son abscisse $a$ sur l'axe des abscisses : l'équation est $x = a$.
\item \textbf{Asymptote horizontale} : repérer une droite horizontale dont la courbe se rapproche quand on suit le graphique vers l'extrême droite ou l'extrême gauche. Lire son ordonnée $b$ sur l'axe des ordonnées : l'équation est $y = b$.
\item \textbf{Conclure en langage de limites} : « la courbe admet l'asymptote $x = a$, donc $\lim_{x \to a} f(x) = \pm\infty$ » ; « la courbe admet l'asymptote $y = b$ en $+\infty$, donc $\lim_{x \to +\infty} f(x) = b$ ».
\end{enumerate}
\end{methode}

\begin{exemplebox}[Fonctions sans asymptote]
Toutes les fonctions n'ont pas d'asymptote parallèle aux axes. La fonction $x \mapsto x^2$ vérifie $\lim_{x \to +\infty} x^2 = +\infty$ : pas d'asymptote horizontale. Définie sur $\mathbb{R}$ tout entier, elle n'a pas de valeur interdite : pas d'asymptote verticale. Sa courbe n'a donc \textbf{aucune} asymptote parallèle aux axes. De même, $g(x) = \dfrac{1}{x^2}$ a bien une asymptote verticale $x = 0$ (car $\lim_{x \to 0} \frac{1}{x^2} = +\infty$) et une asymptote horizontale $y = 0$ (car $\lim_{x \to \pm\infty} \frac{1}{x^2} = 0$), mais la fonction $x \mapsto x^3$ n'en a aucune.
\end{exemplebox}

\begin{exercice}[À toi de jouer]
Soit $g$ la fonction définie par $g(x) = \dfrac{3x - 2}{x + 1}$.
\begin{enumerate}
\item Déterminer le domaine de définition de $g$.
\item Calculer les limites de $g$ aux bornes de son domaine.
\item En déduire les équations des asymptotes à la courbe de $g$, en précisant leur nature.
\item Étudier la position de la courbe par rapport à son asymptote horizontale.
\end{enumerate}
\end{exercice}

\begin{corrige}[Exercice 1]
\textbf{1.} $x + 1 = 0 \Leftrightarrow x = -1$, donc $D_g = \mathbb{R} \setminus \{-1\}$.

\textbf{2.} En la valeur interdite : pour $x$ proche de $-1$, le numérateur vaut $3 \times (-1) - 2 = -5 < 0$. Si $x \to (-1)^+$, $x+1 \to 0^+$, donc $\lim_{x \to (-1)^+} g(x) = -\infty$ ; si $x \to (-1)^-$, $x+1 \to 0^-$, donc $\lim_{x \to (-1)^-} g(x) = +\infty$. À l'infini : $g(x) = \dfrac{3 - \frac{2}{x}}{1 + \frac{1}{x}} \to 3$, donc $\lim_{x \to \pm\infty} g(x) = 3$.

\textbf{3.} La droite $x = -1$ est asymptote verticale ; la droite $y = 3$ est asymptote horizontale en $+\infty$ et en $-\infty$.

\textbf{4.} $g(x) - 3 = \dfrac{3x - 2 - 3(x+1)}{x+1} = \dfrac{-5}{x+1}$. Sur $]-1\,;\,+\infty[$, $g(x) - 3 < 0$ : la courbe est en dessous de l'asymptote. Sur $]-\infty\,;\,-1[$, $g(x) - 3 > 0$ : la courbe est au-dessus.
\end{corrige}

\begin{savaistu}[D'où vient le mot « asymptote » ?]
Le mot \emph{asymptote} vient du grec ancien \emph{asumptôtos}, qui signifie littéralement « \textbf{qui ne se rencontre pas} » (de \emph{a-}, privatif, et \emph{sumpiptein}, tomber ensemble, se rencontrer). Les Grecs de l'Antiquité, notamment Apollonius de Perga (III\textsuperscript{e} siècle av. J.-C.), avaient déjà remarqué que certaines courbes — l'hyperbole en particulier — se rapprochent indéfiniment de droites sans jamais les atteindre. Plus de deux mille ans plus tard, cette même idée aide le trésorier de notre coopérative de cacao à prévoir ses coûts : les mathématiques voyagent à travers les siècles et les continents !
\end{savaistu}

\begin{aretenir}[L'essentiel de la section]
\begin{itemize}
\item $\lim_{x \to \pm\infty} f(x) = b$ (finie) $\Rightarrow$ la droite $y = b$ est \textbf{asymptote horizontale}.
\item $\lim_{x \to a^-} f(x) = \pm\infty$ ou $\lim_{x \to a^+} f(x) = \pm\infty$ $\Rightarrow$ la droite $x = a$ est \textbf{asymptote verticale}.
\item On cherche les asymptotes verticales parmi les \textbf{valeurs interdites}, mais une valeur interdite avec limite finie ne donne \textbf{pas} d'asymptote.
\item La position de la courbe par rapport à l'asymptote $y = b$ se détermine par le \textbf{signe de $f(x) - b$}.
\item Une asymptote décrit un comportement aux bornes : la courbe peut couper son asymptote horizontale ailleurs.
\end{itemize}
\end{aretenir}

Dans la section suivante, nous verrons que certaines courbes — comme celle de la fonction coût de notre coopérative — ne se rapprochent ni d'une horizontale ni d'une verticale, mais d'une droite \emph{inclinée} : ce sont les \textbf{asymptotes obliques}.

\coursec{Asymptotes obliques}

Dans la section précédente, nous avons appris à détecter les asymptotes \emph{parallèles aux axes} : une droite verticale $x = x_0$ quand $f(x)$ explose au voisinage d'une valeur interdite, une droite horizontale $y = \ell$ quand $f(x)$ se stabilise vers un nombre fini en $+\infty$ ou en $-\infty$. Mais toutes les courbes ne se comportent pas ainsi à l'infini. Certaines fonctions « partent » vers l'infini en suivant de très près une droite \emph{inclinée}, ni horizontale ni verticale. C'est précisément l'objet de cette section : donner un sens rigoureux à cette idée, apprendre à trouver une telle droite, et savoir si la courbe la survole ou la longe par en dessous.

\courssub{Motivation : une courbe qui « suit » une droite inclinée}

Considérons la fonction définie par
\[
f(x) = \frac{x^2 + 3x + 1}{x + 1}.
\]
Son ensemble de définition est $\mathbb{R} \setminus \{-1\}$. En $-1$, le dénominateur s'annule sans que le numérateur s'annule (il vaut $1 - 3 + 1 = -1 \neq 0$) : la droite d'équation $x = -1$ est donc une asymptote verticale, comme nous l'avons vu dans la section précédente.

Intéressons-nous maintenant au comportement en $+\infty$. Pour $x$ « grand », le numérateur se comporte comme $x^2$ et le dénominateur comme $x$, donc
\[
\lim_{x \to +\infty} f(x) = \lim_{x \to +\infty} \frac{x^2}{x} = +\infty.
\]
Il n'y a donc \textbf{pas d'asymptote horizontale}. Mais regardons de plus près. Calculons quelques valeurs :

\begin{center}
\begin{tabular}{|c|c|c|c|}
\hline
\rowcolor{popblueL}
$x$ & $f(x)$ & $x + 2$ & $f(x) - (x+2)$ \\
\hline
$10$ & $\dfrac{131}{11} \approx \num{11,909}$ & $12$ & $\approx \num{-0,091}$ \\
\hline
$100$ & $\dfrac{10301}{101} \approx \num{101,990}$ & $102$ & $\approx \num{-0,0099}$ \\
\hline
$1000$ & $\dfrac{1003001}{1001} \approx \num{1001,999}$ & $1002$ & $\approx \num{-0,001}$ \\
\hline
\end{tabular}
\end{center}

Phénomène remarquable : quand $x$ grandit, $f(x)$ devient presque égal à $x + 2$. L'écart $f(x) - (x+2)$ se rapproche de $0$. Graphiquement, cela signifie que la courbe $(\mathcal{C}_f)$ se « colle » à la droite d'équation $y = x + 2$, qui n'est ni horizontale ni verticale : c'est une \cle{asymptote oblique}.

\begin{popfigure}
\begin{center}
\begin{tikzpicture}
\begin{axis}[
  width=0.95\linewidth, height=8.2cm,
  axis lines=middle, xlabel={$x$}, ylabel={$y$},
  xmin=-8, xmax=6, ymin=-8, ymax=8,
  xtick={-7,-6,...,5}, ytick={-7,-6,...,7},
  tick label style={font=\small},
  grid=both, grid style={popline!40},
  legend style={font=\small, at={(0.02,0.98)}, anchor=north west}
]
\addplot[popcurve,domain=-7.5:-1.35,samples=200]{(x^2+3*x+1)/(x+1)};
\addplot[popcurve,forget plot,domain=-0.68:5.5,samples=200]{(x^2+3*x+1)/(x+1)};
\addplot[poporange, dashed, thick, domain=-8:6, samples=2]{x+2};
\draw[popgreen, dashed, thick] (axis cs:-1,-8) -- (axis cs:-1,8);
\legend{$(\mathcal{C}_f)$, $y = x+2$}
\node[popgreen, font=\small, rotate=90, anchor=south] at (axis cs:-1.15,5.5) {$x=-1$};
\end{axis}
\end{tikzpicture}
\end{center}
\legende{Courbe de $f(x) = \dfrac{x^2+3x+1}{x+1}$. En pointillés verts : l'asymptote verticale $x=-1$ ; en pointillés orange : l'asymptote oblique $y = x+2$, que la courbe rejoint en $+\infty$ et en $-\infty$.}
\end{popfigure}

\courssub{Définition de l'asymptote oblique}

L'idée intuitive est simple : une droite $(D)$ est asymptote oblique à une courbe quand la \emph{distance verticale} entre la courbe et la droite tend vers $0$ à l'infini. Or cette distance verticale, pour une abscisse $x$ donnée, vaut $\big| f(x) - (ax+b) \big|$. D'où la définition rigoureuse.

\begin{definition}[Asymptote oblique]
Soit $f$ une fonction définie au voisinage de $+\infty$ (respectivement de $-\infty$) et $(D)$ la droite d'équation $y = ax + b$ avec $a \neq 0$.

On dit que $(D)$ est \cle{asymptote oblique} à la courbe $(\mathcal{C}_f)$ en $+\infty$ (respectivement en $-\infty$) lorsque
\[
\lim_{x \to +\infty} \big[ f(x) - (ax + b) \big] = 0
\qquad \Big(\text{respectivement } \lim_{x \to -\infty} \big[ f(x) - (ax + b) \big] = 0\Big).
\]
\end{definition}

\begin{aretenir}[L'idée à retenir]
Dire que $(D) : y = ax+b$ est asymptote oblique, c'est dire que \textbf{l'écart} entre la courbe et la droite devient aussi petit qu'on veut, pourvu que $x$ soit assez grand. Ce n'est \emph{pas} dire que $f(x)$ tend vers l'infini : c'est dire que $f(x)$ et $ax+b$ tendent vers l'infini \emph{ensemble, en restant collés l'un à l'autre}.
\end{aretenir}

\begin{attention}[Direction asymptotique $\neq$ asymptote oblique]
Piège classique du Probatoire ! Le fait que $\lim\limits_{x \to +\infty} f(x) = +\infty$ ne suffit \textbf{pas} à conclure à l'existence d'une asymptote oblique. Par exemple, pour $g(x) = x^2$, on a bien $\lim\limits_{x \to +\infty} g(x) = +\infty$, mais aucune droite $y = ax+b$ ne vérifie $\lim [g(x) - (ax+b)] = 0$ : la parabole « s'écarte » de toute droite. On dit seulement que $(\mathcal{C}_g)$ admet une \cle{direction asymptotique} parallèle à l'axe des ordonnées. Pour conclure à une asymptote oblique, il faut \emph{impérativement} exhiber $a$ et $b$ tels que l'écart $f(x) - (ax+b)$ tende vers $0$.
\end{attention}

\courssub{Méthode 1 : vérifier directement que l'écart tend vers 0}

Lorsque l'énoncé \emph{donne} la droite candidate, la méthode est directe : on calcule la limite de la différence.

\begin{methode}[Montrer qu'une droite donnée est asymptote oblique]
Pour montrer que la droite $(D) : y = ax + b$ est asymptote oblique à $(\mathcal{C}_f)$ en $+\infty$ :
\begin{enumerate}
\item Former la différence $f(x) - (ax+b)$ et la \textbf{simplifier} (réduire au même dénominateur si $f$ est rationnelle).
\item Calculer $\lim\limits_{x \to +\infty} \big[ f(x) - (ax+b) \big]$.
\item Si cette limite vaut $0$, \textbf{conclure par une phrase} : « la droite $(D)$ d'équation $y = ax+b$ est asymptote oblique à $(\mathcal{C}_f)$ en $+\infty$ ».
\item Recommencer éventuellement l'étude en $-\infty$ : la même droite peut être asymptote des deux côtés, ou non.
\end{enumerate}
\end{methode}

\begin{exemplebox}[Vérification directe avec $f(x) = \dfrac{x^2+3x+1}{x+1}$]
Montrons que $(D) : y = x+2$ est asymptote oblique à $(\mathcal{C}_f)$ en $+\infty$. Pour $x \neq -1$ :
\begin{align*}
f(x) - (x+2) &= \frac{x^2+3x+1}{x+1} - (x+2) = \frac{x^2+3x+1 - (x+2)(x+1)}{x+1} \\
&= \frac{x^2+3x+1 - (x^2+3x+2)}{x+1} = \frac{-1}{x+1}.
\end{align*}
Or $\lim\limits_{x \to +\infty} \dfrac{-1}{x+1} = 0$. Donc $(D) : y = x+2$ est asymptote oblique à $(\mathcal{C}_f)$ en $+\infty$. Le même calcul montre que $\lim\limits_{x \to -\infty} \dfrac{-1}{x+1} = 0$ : $(D)$ est aussi asymptote oblique en $-\infty$.
\end{exemplebox}

\courssub{Méthode 2 : la division euclidienne du numérateur par le dénominateur}

Comment \emph{trouver} la droite $y = ax+b$ quand l'énoncé ne la donne pas ? Pour une fonction rationnelle dont le numérateur est de degré exactement un de plus que le dénominateur, l'outil roi est la \cle{division euclidienne} des polynômes, vue en Seconde.

Effectuons la division de $x^2 + 3x + 1$ par $x + 1$ :
\[
x^2 + 3x + 1 = (x+1)(x+2) - 1.
\]
En divisant les deux membres par $x+1$ (pour $x \neq -1$) :
\[
f(x) = \frac{x^2+3x+1}{x+1} = x + 2 - \frac{1}{x+1}.
\]
Sous cette forme, tout devient transparent : $f(x)$ est la somme de la partie affine $x+2$ et d'un terme $\dfrac{-1}{x+1}$ qui tend vers $0$ à l'infini. L'asymptote oblique est donc $y = x+2$ : c'est le \textbf{quotient} de la division euclidienne.

\begin{propriete}[Asymptote oblique des fonctions $x \mapsto \dfrac{ax^2+bx+c}{dx+e}$ (admise)]
Soit $f$ la fonction définie par $f(x) = \dfrac{ax^2+bx+c}{dx+e}$ avec $a \neq 0$ et $d \neq 0$. La division euclidienne de $ax^2+bx+c$ par $dx+e$ s'écrit
\[
ax^2+bx+c = (dx+e)(a'x + b') + r, \qquad r \in \mathbb{R},
\]
d'où, pour tout $x \neq -\dfrac{e}{d}$ :
\[
f(x) = a'x + b' + \frac{r}{dx+e}.
\]
Comme $\lim\limits_{x \to \pm\infty} \dfrac{r}{dx+e} = 0$, la droite d'équation $y = a'x + b'$ (le \textbf{quotient} de la division) est asymptote oblique à $(\mathcal{C}_f)$ en $+\infty$ et en $-\infty$. \emph{Cette propriété est admise ; en revanche, savoir effectuer la division et conclure est exigible.}
\end{propriete}

\begin{aretenir}[Réflexe à acquérir]
Devant une fonction du type $f(x) = \dfrac{ax^2+bx+c}{dx+e}$ ($a \neq 0$, $d \neq 0$) : il existe \emph{toujours} une asymptote oblique, et on l'obtient en posant la division euclidienne du numérateur par le dénominateur. L'asymptote est $y = \text{quotient}$.
\end{aretenir}

\courssub{Position relative de la courbe et de son asymptote oblique}

Une fois l'asymptote trouvée, une question naturelle se pose : la courbe est-elle \emph{au-dessus} ou \emph{en dessous} de la droite ? La réponse se lit sur le signe de l'écart $f(x) - (ax+b)$.

\begin{propriete}[Position relative]
Soit $(D) : y = ax+b$ asymptote oblique à $(\mathcal{C}_f)$ en $+\infty$ (ou en $-\infty$).
\begin{itemize}
\item Si $f(x) - (ax+b) > 0$ sur un intervalle, alors $(\mathcal{C}_f)$ est \textbf{au-dessus} de $(D)$ sur cet intervalle ;
\item Si $f(x) - (ax+b) < 0$ sur un intervalle, alors $(\mathcal{C}_f)$ est \textbf{en dessous} de $(D)$ sur cet intervalle.
\end{itemize}
\end{propriete}

Reprenons notre exemple : $f(x) - (x+2) = \dfrac{-1}{x+1}$.
\begin{itemize}
\item Si $x > -1$, alors $x+1 > 0$, donc $\dfrac{-1}{x+1} < 0$ : la courbe est \textbf{en dessous} de son asymptote.
\item Si $x < -1$, alors $x+1 < 0$, donc $\dfrac{-1}{x+1} > 0$ : la courbe est \textbf{au-dessus} de son asymptote.
\end{itemize}
C'est exactement ce que l'on observe sur la figure : à droite, la courbe rampe sous la droite orange ; à gauche, elle la survole.

\begin{popfigure}
\begin{center}
\begin{tikzpicture}
\begin{axis}[
  width=0.95\linewidth, height=7.5cm,
  axis lines=middle, xlabel={$x$}, ylabel={$y$},
  xmin=0, xmax=60, ymin=1.5, ymax=62.5,
  xtick={0,10,...,60}, ytick={10,20,...,60},
  tick label style={font=\small},
  grid=both, grid style={popline!40},
  legend style={font=\small, at={(0.98,0.05)}, anchor=south east}
]
\addplot[popcurve,domain=0.5:60,samples=300]{(x^2+3*x+1)/(x+1)};
\addplot[poporange, dashed, thick, domain=0:60, samples=2]{x+2};
\legend{$(\mathcal{C}_f)$, $y = x+2$}
\draw[<->, poppink, thick] (axis cs:40,42) -- (axis cs:40,41.0244) node[midway, right, font=\small, poppink] {écart $= \dfrac{1}{x+1}$};
\end{axis}
\end{tikzpicture}
\end{center}
\legende{Zoom sur les grandes valeurs de $x$ : la courbe (bleue) se confond presque avec l'asymptote (orange, pointillés). L'écart vertical $\dfrac{1}{x+1}$ devient invisible à l'œil nu : il tend vers $0$.}
\end{popfigure}

\courssub{Exercice résolu et complément pour la Terminale}

\begin{exoresolu}[Étude complète d'une asymptote oblique]
Soit $g$ la fonction définie par $g(x) = \dfrac{2x^2 - 3x + 4}{x - 1}$ et $(\mathcal{C}_g)$ sa courbe représentative.
\begin{enumerate}
\item Déterminer l'ensemble de définition de $g$.
\item Montrer que pour tout $x \neq 1$, $g(x) = 2x - 1 + \dfrac{3}{x-1}$.
\item En déduire que $(\mathcal{C}_g)$ admet une asymptote oblique $(\Delta)$ dont on donnera l'équation.
\item Étudier la position relative de $(\mathcal{C}_g)$ et de $(\Delta)$.
\end{enumerate}
\tcblower
\textbf{1.} $g$ est définie lorsque $x - 1 \neq 0$, donc $\mathcal{D}_g = \mathbb{R} \setminus \{1\}$.

\textbf{2.} Effectuons la division euclidienne de $2x^2 - 3x + 4$ par $x - 1$ :
\[
2x^2 - 3x + 4 = (x-1)(2x - 1) + 3.
\]
Vérification : $(x-1)(2x-1) = 2x^2 - x - 2x + 1 = 2x^2 - 3x + 1$, et $2x^2 - 3x + 1 + 3 = 2x^2 - 3x + 4$. Correct. En divisant par $x-1 \neq 0$ :
\[
g(x) = 2x - 1 + \frac{3}{x-1}.
\]

\textbf{3.} On a $g(x) - (2x-1) = \dfrac{3}{x-1}$, et
\[
\lim_{x \to +\infty} \frac{3}{x-1} = 0 \qquad \text{et} \qquad \lim_{x \to -\infty} \frac{3}{x-1} = 0.
\]
Donc la droite $(\Delta) : y = 2x - 1$ est asymptote oblique à $(\mathcal{C}_g)$ en $+\infty$ et en $-\infty$.

\textbf{4.} Le signe de l'écart $g(x) - (2x-1) = \dfrac{3}{x-1}$ est celui de $x-1$ :
\begin{itemize}
\item si $x > 1$, l'écart est positif : $(\mathcal{C}_g)$ est \textbf{au-dessus} de $(\Delta)$ ;
\item si $x < 1$, l'écart est négatif : $(\mathcal{C}_g)$ est \textbf{en dessous} de $(\Delta)$.
\end{itemize}
\end{exoresolu}

\begin{attention}[Erreurs fréquentes à éviter]
\begin{itemize}
\item \textbf{Oublier de conclure.} Calculer la limite de l'écart et trouver $0$ ne suffit pas : il faut écrire la phrase de conclusion « donc la droite d'équation ... est asymptote oblique à la courbe en ... ». C'est elle qui rapporte le point.
\item \textbf{Se tromper dans la division.} Vérifie toujours ton quotient en redéveloppant : $(dx+e)(a'x+b') + r$ doit redonner exactement le numérateur.
\item \textbf{Confondre les rôles.} L'asymptote est le \emph{quotient} $a'x + b'$, jamais le reste $r$.
\item \textbf{Étudier la position sans l'écart.} La position relative se déduit du signe de $f(x) - (ax+b)$, pas d'une impression sur le graphique.
\end{itemize}
\end{attention}

\begin{savaistu}[Pour aller plus loin : trouver $a$ et $b$ sans connaître la droite]
En Terminale, tu apprendras une méthode générale pour deviner l'asymptote oblique d'une fonction quelconque (pas seulement rationnelle). Si $(D) : y = ax+b$ est asymptote oblique en $+\infty$, alors nécessairement
\[
a = \lim_{x \to +\infty} \frac{f(x)}{x} \qquad \text{puis} \qquad b = \lim_{x \to +\infty} \big[ f(x) - ax \big].
\]
Vérifions sur notre exemple : $\dfrac{f(x)}{x} = \dfrac{x^2+3x+1}{x(x+1)}$ tend vers $1$, donc $a = 1$ ; puis $f(x) - x = \dfrac{x^2+3x+1}{x+1} - x = \dfrac{2x+1}{x+1}$ tend vers $2$, donc $b = 2$. On retrouve $y = x + 2$. En Première, cette méthode n'est pas exigible, mais elle éclaire pourquoi la division euclidienne « marche » : le quotient capture exactement la partie affine de $f$.
\end{savaistu}

\begin{savaistu}[Les asymptotes autour de nous]
Les asymptotes obliques apparaissent dès qu'une grandeur croît « presque proportionnellement » à une autre, avec un écart qui s'efface. Par exemple, le coût total d'une ligne de production (usine de brasserie à Douala, unité de transformation de manioc) comprend un coût fixe et un coût variable : le coût \emph{moyen} par unité suit souvent une courbe qui se rapproche d'une droite oblique ou d'une asymptote horizontale. Comprendre ces comportements limites, c'est pouvoir prévoir — et c'est exactement le métier du mathématicien... et du bon gestionnaire.
\end{savaistu}

\coursec{Éléments de symétrie d'une courbe}

Dans la section précédente, nous avons appris à détecter les \cle{asymptotes obliques}, ces droites dont la courbe se rapproche indéfiniment aux grandes valeurs de $x$. Nous savons donc désormais décrire le comportement d'une courbe \emph{aux bords} de son domaine. Mais une courbe possède souvent une autre propriété remarquable, qui concerne son \emph{architecture globale} : la \cle{symétrie}. Observe un papillon, la façade d'une case traditionnelle, ou le logo d'une compagnie : la symétrie est partout autour de nous. En mathématiques, reconnaître un centre ou un axe de symétrie d'une courbe présente un avantage considérable : \textbf{il suffit d'étudier la fonction sur la moitié de son domaine}, puis de compléter la courbe par symétrie. C'est un gain de temps précieux au Probatoire, et c'est l'objet de cette section.

\courssub{Rappels : fonctions paires et fonctions impaires}

Tu as déjà rencontré en classe de Seconde deux situations de symétrie particulières, liées au repère lui-même.

\begin{definition}[Fonction paire, fonction impaire]
Soit $f$ une fonction de domaine $D_f$.
\begin{itemize}
\item $f$ est \cle{paire} si $D_f$ est symétrique par rapport à $0$ (c'est-à-dire : $x \in D_f \Rightarrow -x \in D_f$) et si pour tout $x \in D_f$ : $f(-x) = f(x)$. Sa courbe admet alors \textbf{l'axe des ordonnées} pour axe de symétrie.
\item $f$ est \cle{impaire} si $D_f$ est symétrique par rapport à $0$ et si pour tout $x \in D_f$ : $f(-x) = -f(x)$. Sa courbe admet alors \textbf{l'origine $O$ du repère} pour centre de symétrie.
\end{itemize}
\end{definition}

\begin{exemplebox}[Rappels illustrés]
La fonction $x \mapsto x^2$ est paire : $(-x)^2 = x^2$. Sa parabole est symétrique par rapport à l'axe des ordonnées. La fonction $x \mapsto \dfrac{1}{x}$ est impaire : $\dfrac{1}{-x} = -\dfrac{1}{x}$. Son hyperbole est symétrique par rapport à l'origine.
\end{exemplebox}

Ces deux cas sont des cas \emph{particuliers} : le centre de symétrie est $O(0\,;0)$, l'axe de symétrie est la droite $x = 0$. Mais qu'en est-il lorsque la symétrie se fait par rapport à un point $\Omega(a\,;b)$ quelconque, ou par rapport à une droite verticale $x = a$ quelconque ? C'est toute la question de cette section. L'idée directrice sera la suivante : \textbf{se ramener au cas connu} (parité ou imparité) par un simple changement de repère.

\courssub{Centre de symétrie d'une courbe}

\textbf{Intuition.} Dire qu'un point $\Omega$ est centre de symétrie d'une courbe $(\mathcal{C}_f)$ signifie que la courbe est \emph{globalement invariante} par la symétrie centrale de centre $\Omega$ : chaque fois qu'un point $M$ appartient à la courbe, son symétrique $M'$ par rapport à $\Omega$ appartient aussi à la courbe. Autrement dit, $\Omega$ est le \cle{milieu} de tout segment $[MM']$ joignant deux points symétriques de la courbe. Tourne la feuille d'un demi-tour autour de $\Omega$ : la courbe retombe exactement sur elle-même.

\begin{definition}[Centre de symétrie]
Soit $f$ une fonction de courbe représentative $(\mathcal{C}_f)$ dans un repère $(O\,;\vec{\imath},\vec{\jmath})$. Le point $\Omega(a\,;b)$ est un \cle{centre de symétrie} de $(\mathcal{C}_f)$ si pour tout point $M(x\,;f(x))$ de la courbe, le symétrique $M'$ de $M$ par rapport à $\Omega$ appartient encore à $(\mathcal{C}_f)$.
\end{definition}

\begin{propriete}[Caractérisation du centre de symétrie]
Soit $f$ une fonction de domaine $D_f$ et $\Omega(a\,;b)$ un point. Alors :
\[ \Omega(a\,;b) \text{ est centre de symétrie de } (\mathcal{C}_f) \iff \begin{cases} D_f \text{ est symétrique par rapport à } a,\\ \text{pour tout } h \text{ tel que } a+h \in D_f : \ f(a+h) + f(a-h) = 2b. \end{cases} \]
De manière équivalente : $\dfrac{f(a+h)+f(a-h)}{2} = b$, ou encore $f(a+h) - b = -\bigl(f(a-h) - b\bigr)$.
\end{propriete}

\begin{experience}[Démonstration guidée : du milieu d'un segment à la formule]
Nous allons établir cette caractérisation à partir de la définition du milieu d'un segment. Cette démonstration est formatrice et sa démarche est exigible.
\begin{enumerate}
\item \textbf{Écris les coordonnées des points en jeu.} Soit $M$ un point de la courbe d'abscisse $a+h$ : $M\bigl(a+h\,;\, f(a+h)\bigr)$. Son symétrique par rapport à $\Omega$, s'il est sur la courbe, est le point $M'$ d'abscisse $a-h$ : $M'\bigl(a-h\,;\, f(a-h)\bigr)$.
\item \textbf{Traduis que $\Omega$ est le milieu de $[MM']$.} Par la formule du milieu, $\Omega$ est milieu de $[MM']$ si et seulement si :
\[ \frac{(a+h)+(a-h)}{2} = a \qquad \text{et} \qquad \frac{f(a+h)+f(a-h)}{2} = b. \]
\item \textbf{Vérifie la première condition.} $\dfrac{(a+h)+(a-h)}{2} = \dfrac{2a}{2} = a$ : elle est \emph{toujours} vraie, quel que soit $h$. L'abscisse ne pose donc aucune condition, si ce n'est que $a+h$ et $a-h$ doivent tous deux appartenir à $D_f$ : c'est exactement dire que $D_f$ est symétrique par rapport à $a$.
\item \textbf{Conclus.} Il reste la condition sur les ordonnées : $\dfrac{f(a+h)+f(a-h)}{2} = b$, c'est-à-dire $f(a+h) + f(a-h) = 2b$ pour tout $h$ admissible. $\blacksquare$
\end{enumerate}
\end{experience}

\begin{attention}[Le domaine d'abord, les calculs ensuite !]
Avant tout calcul de $f(a+h) + f(a-h)$, \textbf{vérifie que $D_f$ est symétrique par rapport à $a$}. Si ce n'est pas le cas, $\Omega(a\,;b)$ ne peut pas être centre de symétrie, et tout calcul serait inutile. Exemple : pour $f(x) = \dfrac{x-1}{x-2}$, $D_f = \mathbb{R} \setminus \{2\}$, qui est bien symétrique par rapport à $2$ (car $2+h \in D_f \Leftrightarrow h \neq 0 \Leftrightarrow 2-h \in D_f$). En revanche, la courbe de $x \mapsto \sqrt{x}$ ne peut admettre aucun centre de symétrie d'abscisse $a > 0$, car $D_f = [0\,;+\infty[$ n'est symétrique par rapport à aucun réel $a$.
\end{attention}

\courssub{Axe de symétrie d'une courbe}

\textbf{Intuition.} Une droite verticale d'équation $x = a$ est axe de symétrie de la courbe si celle-ci se superpose à elle-même par pliage le long de cette droite : à deux points de la courbe situés à égale distance de l'axe correspondent la \emph{même ordonnée}.

\begin{definition}[Axe de symétrie parallèle à l'axe des ordonnées]
La droite d'équation $x = a$ est un \cle{axe de symétrie} de $(\mathcal{C}_f)$ si le symétrique de tout point de la courbe par rapport à cette droite appartient encore à la courbe.
\end{definition}

\begin{propriete}[Caractérisation de l'axe de symétrie]
Soit $f$ une fonction de domaine $D_f$. Alors :
\[ \text{La droite } x = a \text{ est axe de symétrie de } (\mathcal{C}_f) \iff \begin{cases} D_f \text{ est symétrique par rapport à } a,\\ \text{pour tout } h \text{ tel que } a+h \in D_f : \ f(a+h) = f(a-h). \end{cases} \]
\end{propriete}

La justification est immédiate : le symétrique du point $M\bigl(a+h\,;f(a+h)\bigr)$ par rapport à la droite $x = a$ est le point de même ordonnée et d'abscisse $a - h$. Il est sur la courbe si et seulement si $f(a-h) = f(a+h)$.

\courssub{Méthode pratique : le changement de repère}

La caractérisation précédente est efficace, mais il existe une méthode équivalente très élégante, qui ramène tout au cas connu des fonctions paires et impaires : \textbf{transporter l'origine du repère au point $\Omega$}.

\begin{methode}[Prouver une symétrie par changement de repère]
Pour montrer que $\Omega(a\,;b)$ est centre de symétrie de $(\mathcal{C}_f)$ :
\begin{enumerate}
\item Pose le changement de coordonnées $X = x - a$ et $Y = y - b$ (c'est-à-dire $x = X + a$, $y = Y + b$). Le point $\Omega$ devient l'origine du nouveau repère $(\Omega\,;\vec{\imath},\vec{\jmath})$.
\item Remplace $x$ par $X + a$ et $y$ par $Y + b$ dans l'équation $y = f(x)$, puis simplifie pour obtenir une équation de la forme $Y = F(X)$.
\item Étudie la \cle{parité} de $F$ :
\begin{itemize}
\item si $F$ est \textbf{impaire} ($F(-X) = -F(X)$), alors $\Omega$ est \textbf{centre de symétrie} ;
\item si $F$ est \textbf{paire} ($F(-X) = F(X)$), alors la droite $x = a$ est \textbf{axe de symétrie}.
\end{itemize}
\item Conclus en rédigeant clairement la propriété obtenue.
\end{enumerate}
Pour un axe de symétrie $x = a$ seul, on pose simplement $X = x - a$ (avec $Y = y$) et l'on montre que la fonction obtenue est paire.
\end{methode}

\begin{exemplebox}[Centre de symétrie de la courbe de $f(x) = \dfrac{x-1}{x-2}$]
Montrons que $\Omega(2\,;1)$ est centre de symétrie de $(\mathcal{C}_f)$, où $D_f = \mathbb{R}\setminus\{2\}$ (symétrique par rapport à $2$ : vérifié plus haut).

\emph{Par la caractérisation.} Pour $h \neq 0$ :
\begin{align*}
f(2+h) + f(2-h) &= \frac{(2+h)-1}{(2+h)-2} + \frac{(2-h)-1}{(2-h)-2} = \frac{1+h}{h} + \frac{1-h}{-h} \\
&= \frac{1+h}{h} + \frac{h-1}{h} = \frac{2h}{h} = 2 = 2 \times 1.
\end{align*}
Donc $f(2+h) + f(2-h) = 2b$ avec $b = 1$ : $\Omega(2\,;1)$ est bien centre de symétrie.

\emph{Par changement de repère.} Posons $X = x - 2$, $Y = y - 1$. Alors $x = X+2$ et :
\[ Y = f(X+2) - 1 = \frac{X+1}{X} - 1 = \frac{X+1-X}{X} = \frac{1}{X}. \]
La fonction $F : X \mapsto \dfrac{1}{X}$ est impaire, donc $\Omega(2\,;1)$ est centre de symétrie de $(\mathcal{C}_f)$. Remarque au passage : on a montré que $f(x) = 1 + \dfrac{1}{x-2}$, ce qui sera précieux pour tracer la courbe.
\end{exemplebox}

\begin{popfigure}
\begin{center}
\begin{tikzpicture}[scale=1.05,>=Stealth]
\draw[popline!60] (-2.2,-2.2) grid[step=1] (6.2,4.2);
\draw[->,popdark,thick] (-2.2,0) -- (6.2,0) node[below left]{$x$};
\draw[->,popdark,thick] (0,-2.2) -- (0,4.2) node[below left]{$y$};
\draw[popmuted,dashed,thick] (2,-2.2) -- (2,4.2);
\draw[popmuted,dashed,thick] (-2.2,1) -- (6.2,1);
\draw[popblue,very thick,domain=-2.1:1.55,samples=120,smooth] plot (\x,{1+1/(\x-2)});
\draw[popblue,very thick,domain=2.45:6.1,samples=120,smooth] plot (\x,{1+1/(\x-2)});
\fill[poporange] (2,1) circle (2.4pt) node[above right,popdark]{$\Omega(2\,;1)$};
\fill[popgreen] (4,1.5) circle (2.2pt) node[above right,popdark]{$M\bigl(2+h\,;f(2+h)\bigr)$};
\fill[poppink] (0,0.5) circle (2.2pt) node[below left,popdark]{$M'\bigl(2-h\,;f(2-h)\bigr)$};
\draw[popgold,thick,dashed] (4,1.5) -- (0,0.5);
\node[popmuted] at (5.4,0.6) {$x=2$};
\node[popmuted] at (-1.6,1.35) {$y=1$};
\end{tikzpicture}
\end{center}
\legende{Courbe de $f(x) = \dfrac{1}{x-2} + 1$ : hyperbole de centre de symétrie $\Omega(2\,;1)$. Les points $M$ et $M'$ sont symétriques par rapport à $\Omega$, milieu de $[MM']$.}
\end{popfigure}

\courssub{Applications fondamentales : hyperbole et parabole}

\begin{propriete}[Centre de symétrie de l'hyperbole homographique]
Soit $f(x) = \dfrac{ax+b}{cx+d}$ avec $c \neq 0$ et $ad - bc \neq 0$. La courbe de $f$ est une hyperbole qui admet pour \cle{centre de symétrie} le point
\[ \Omega\left(-\frac{d}{c}\,;\,\frac{a}{c}\right), \]
intersection de ses deux asymptotes $x = -\dfrac{d}{c}$ et $y = \dfrac{a}{c}$.
\end{propriete}

En effet, la division suivant les puissances décroissantes permet d'écrire $f(x) = \dfrac{a}{c} + \dfrac{k}{x + \frac{d}{c}}$ pour une constante $k$ ; le changement de repère $X = x + \dfrac{d}{c}$, $Y = y - \dfrac{a}{c}$ donne alors $Y = \dfrac{k}{X}$, fonction impaire.

\begin{propriete}[Axe de symétrie de la parabole]
Soit $g(x) = ax^2 + bx + c$ avec $a \neq 0$. La parabole de $g$ admet pour \cle{axe de symétrie} la droite d'équation
\[ x = -\frac{b}{2a}, \]
qui passe par le sommet $S\left(-\dfrac{b}{2a}\,;\, g\left(-\dfrac{b}{2a}\right)\right)$.
\end{propriete}

Vérifions-le : en posant $\alpha = -\dfrac{b}{2a}$, la forme canonique $g(x) = a(x-\alpha)^2 + g(\alpha)$ donne $g(\alpha + h) = a h^2 + g(\alpha)$ et $g(\alpha - h) = a h^2 + g(\alpha)$, donc $g(\alpha+h) = g(\alpha-h)$ pour tout réel $h$.

\begin{exoresolu}[Axe de symétrie de la parabole de $g(x) = x^2 - 4x + 1$]
Montrer que la droite d'équation $x = 2$ est axe de symétrie de la courbe $(\mathcal{C}_g)$, puis préciser le sommet.
\tcblower
Le domaine $D_g = \mathbb{R}$ est symétrique par rapport à $2$. Pour tout réel $h$ :
\begin{align*}
g(2+h) &= (2+h)^2 - 4(2+h) + 1 = 4 + 4h + h^2 - 8 - 4h + 1 = h^2 - 3,\\
g(2-h) &= (2-h)^2 - 4(2-h) + 1 = 4 - 4h + h^2 - 8 + 4h + 1 = h^2 - 3.
\end{align*}
Donc $g(2+h) = g(2-h)$ pour tout $h$ : la droite $x = 2$ est axe de symétrie de $(\mathcal{C}_g)$. Cohérent avec la formule : $-\dfrac{b}{2a} = -\dfrac{-4}{2} = 2$. Le sommet est $S(2\,;\,g(2)) = (2\,;-3)$. On retrouve d'ailleurs la forme canonique $g(x) = (x-2)^2 - 3$.
\end{exoresolu}

\begin{popfigure}
\begin{center}
\begin{tikzpicture}
\begin{axis}[
width=0.85\linewidth, height=7cm,
axis lines=middle, xlabel={$x$}, ylabel={$y$},
xmin=-1.5, xmax=5.5, ymin=-4.5, ymax=5.5,
xtick={-1,0,1,2,3,4,5}, ytick={-3,-2,-1,1,2,3,4,5},
grid=both, grid style={popline!50},
tick label style={font=\small}, label style={font=\small}
]
\addplot[popcurve,domain=-1.1:5.1,samples=200]{x^2 - 4*x + 1};
\addplot[poporange,dashed,thick] coordinates {(2,-4.2) (2,5.2)};
\addplot[mark=*,mark size=1.8pt,popgreen] coordinates {(2,-3)};
\node[popdark,font=\small] at (axis cs:2.9,-3.3) {$S(2\,;-3)$};
\node[poporange,font=\small] at (axis cs:2.35,4.6) {$x=2$};
\addplot[mark=*,mark size=1.5pt,poppink] coordinates {(0,1)};
\addplot[mark=*,mark size=1.5pt,poppink] coordinates {(4,1)};
\node[popdark,font=\small] at (axis cs:-0.15,1.55) {$g(2-h)$};
\node[popdark,font=\small] at (axis cs:4.15,1.55) {$g(2+h)$};
\end{axis}
\end{tikzpicture}
\end{center}
\legende{Parabole de $g(x) = x^2 - 4x + 1$ et son axe de symétrie $x = 2$ (pointillés) : deux points d'abscisses $2-h$ et $2+h$ ont la même image.}
\end{popfigure}

\courssub{Intérêt pratique : réduire le domaine d'étude}

\begin{aretenir}[Pourquoi chercher les symétries ?]
Lorsque la courbe $(\mathcal{C}_f)$ admet un centre de symétrie $\Omega(a\,;b)$ ou un axe de symétrie $x = a$ :
\begin{enumerate}
\item on \textbf{restreint l'étude} (limites, dérivée, variations) à la moitié du domaine, par exemple $D_f \cap [a\,;+\infty[$ ;
\item on trace la partie de courbe correspondante ;
\item on \textbf{complète par symétrie} : symétrie centrale de centre $\Omega$, ou symétrie orthogonale d'axe $x = a$.
\end{enumerate}
On gagne ainsi du temps et de la précision, et l'on évite les incohérences de tracé : une courbe qui ne respecte pas sa symétrie annoncée est une courbe fausse.
\end{aretenir}

\begin{attention}[Pièges de rédaction au Probatoire]
\begin{itemize}
\item Ne confonds pas les deux caractérisations : \textbf{centre} de symétrie $\Leftrightarrow f(a+h) + f(a-h) = 2b$ (on \emph{additionne}) ; \textbf{axe} de symétrie $\Leftrightarrow f(a+h) = f(a-h)$ (on \emph{compare}).
\item L'égalité doit être démontrée \textbf{pour tout $h$ admissible}, pas vérifiée sur une seule valeur : tester $h = 1$ ne prouve rien.
\item Si $D_f$ n'est pas symétrique par rapport à $a$, conclus immédiatement qu'il n'y a ni centre d'abscisse $a$ ni axe $x = a$.
\item Une fonction impaire n'est qu'un cas particulier : centre $O(0\,;0)$. Ne conclus pas « $f$ impaire » quand le centre est $\Omega(a\,;b) \neq O$ ; rédige avec la caractérisation générale.
\end{itemize}
\end{attention}

\begin{savaistu}[La symétrie, un outil d'économie... et d'élégance]
Au Cameroun, les tisserands qui confectionnent le \emph{toghu} du Nord-Ouest exploitent des axes de symétrie pour reproduire leurs motifs par pliage : un seul gabarit suffit pour toute l'étoffe. Le mathématicien fait exactement pareil avec les courbes ! Plus profondément, la grande mathématicienne allemande Emmy Noether a montré au XX\textsuperscript{e} siècle que chaque symétrie d'une loi physique correspond à une quantité conservée (énergie, quantité de mouvement...) : la symétrie est l'un des concepts les plus puissants de toute la science.
\end{savaistu}

\begin{exercice}[Pour t'entraîner]
\begin{enumerate}
\item Soit $f(x) = \dfrac{2x+3}{x-1}$. Montrer que $\Omega(1\,;2)$ est centre de symétrie de $(\mathcal{C}_f)$, d'abord par la caractérisation, puis par changement de repère.
\item Soit $g(x) = -x^2 + 6x - 5$. Déterminer l'axe de symétrie de sa courbe et son sommet.
\item La courbe de $h(x) = x^3 - 3x$ admet-elle un centre de symétrie ? Lequel ?
\end{enumerate}
\end{exercice}

\begin{corrige}[Exercice]
\begin{enumerate}
\item $D_f = \mathbb{R}\setminus\{1\}$, symétrique par rapport à $1$. Pour $h \neq 0$ : $f(1+h) = \dfrac{2h+5}{h}$ et $f(1-h) = \dfrac{5-2h}{-h} = \dfrac{2h-5}{h}$. Somme : $\dfrac{2h+5+2h-5}{h} = \dfrac{4h}{h} = 4 = 2\times 2$. Donc $\Omega(1\,;2)$ est centre de symétrie. Par changement de repère : $X = x-1$, $Y = y-2$ donnent $Y = \dfrac{2(X+1)+3}{X} - 2 = \dfrac{2X+5-2X}{X} = \dfrac{5}{X}$, fonction impaire : confirmation.
\item $-\dfrac{b}{2a} = -\dfrac{6}{-2} = 3$ : axe $x = 3$. Vérification : $g(3+h) = -(9+6h+h^2) + 18 + 6h - 5 = 4 - h^2 = g(3-h)$. Sommet $S(3\,;4)$.
\item $D_h = \mathbb{R}$ et $h(-x) = -x^3 + 3x = -h(x)$ : $h$ est impaire, donc l'origine $O(0\,;0)$ est centre de symétrie de sa courbe.
\end{enumerate}
\end{corrige}

Tu sais désormais reconnaître et prouver les éléments de symétrie d'une courbe, et surtout les exploiter pour étudier une fonction deux fois plus vite. Mettons immédiatement ces outils en action dans la section suivante, consacrée à l'étude complète des \textbf{fonctions homographiques}, dont tu connais déjà le centre de symétrie !

\coursec{Étude des fonctions homographiques}

Dans la section précédente, nous avons appris à reconnaître les éléments de symétrie d'une courbe : un point $\Omega$ peut être centre de symétrie, une droite verticale peut être axe de symétrie. Nous allons maintenant mettre ces outils en action sur une famille de fonctions extrêmement importante au programme de Première : les \cle{fonctions homographiques}. Tu les as déjà rencontrées sans le savoir : la fonction inverse $x \mapsto \dfrac{1}{x}$ en est le prototype. Son étude complète — domaine, limites, asymptotes, variations, symétrie, tracé — est un classique incontournable du Probatoire.

\courssub{Définition et domaine de définition}

\begin{definition}[Fonction homographique]
On appelle \cle{fonction homographique} toute fonction $f$ qui peut s'écrire sous la forme
\[ f(x) = \frac{ax+b}{cx+d} \quad \text{avec } c \neq 0 \ \text{et}\ ad - bc \neq 0, \]
où $a$, $b$, $c$, $d$ sont des réels. Le réel $ad - bc$ est appelé le \cle{déterminant} de la fonction homographique.
\end{definition}

Pourquoi ces deux conditions ? Regardons ce qui se passe si elles sont violées.

\begin{itemize}
\item Si $c = 0$, alors $f(x) = \dfrac{ax+b}{d}$ : c'est une simple fonction \emph{affine}, déjà étudiée en Seconde. Rien de nouveau.
\item Si $ad - bc = 0$, alors $ad = bc$. Comme $c \neq 0$, on a $\dfrac{a}{c} = \dfrac{b}{d}$ si $d \neq 0$ ; notons $k$ cette valeur commune : $a = kc$ et $b = kd$, donc
\[ f(x) = \frac{kc\,x + kd}{cx+d} = \frac{k(cx+d)}{cx+d} = k. \]
(Si $d = 0$, alors $bc = 0$, donc $b = 0$ puisque $c \neq 0$, et $f(x) = \dfrac{ax}{cx} = \dfrac{a}{c}$ : même conclusion.) La fonction est \emph{constante} : son étude est triviale et sa courbe est une droite horizontale privée d'un point.
\end{itemize}

La condition $ad - bc \neq 0$ garantit donc que la fonction homographique est une \emph{vraie} fonction homographique, non dégénérée.

\begin{propriete}[Domaine de définition]
Le domaine de définition de $f(x) = \dfrac{ax+b}{cx+d}$ (avec $c \neq 0$) est
\[ D_f = \mathbb{R} \setminus \left\{ -\frac{d}{c} \right\} = \left] -\infty ; -\frac{d}{c} \right[ \cup \left] -\frac{d}{c} ; +\infty \right[. \]
\end{propriete}

En effet, une fraction n'existe que si son dénominateur est non nul : $cx + d = 0 \iff x = -\dfrac{d}{c}$. Cette valeur $-\dfrac{d}{c}$ est la \cle{valeur interdite} : elle jouera un rôle central dans toute l'étude.

\courssub{La forme canonique : réduire l'étude à la fonction inverse}

Toute la théorie des fonctions homographiques repose sur une idée simple : se ramener à la fonction inverse $x \mapsto \dfrac{1}{x}$, que l'on connaît bien, par un changement de repère. Concrètement, on cherche à écrire
\[ f(x) = \alpha + \frac{\beta}{x - x_0}. \]

\begin{methode}[Obtenir la forme canonique]
Pour écrire $f(x) = \dfrac{ax+b}{cx+d}$ sous la forme $\alpha + \dfrac{\beta}{x - x_0}$ :
\begin{enumerate}
\item Identifier $x_0 = -\dfrac{d}{c}$ (la valeur interdite).
\item Écrire le numérateur en fonction de $(cx+d)$ : on cherche une constante $\alpha$ telle que $ax + b = \dfrac{a}{c}(cx+d) + \text{reste}$, c'est-à-dire que l'on effectue la \emph{division} du numérateur par le dénominateur.
\item Simplifier : $f(x) = \dfrac{a}{c} + \dfrac{b - \frac{ad}{c}}{cx+d} = \dfrac{a}{c} + \dfrac{bc - ad}{c(cx+d)}$.
\end{enumerate}
Comme $c(cx+d) = c^2\left(x + \dfrac{d}{c}\right) = c^2(x - x_0)$, on obtient finalement
\[ f(x) = \alpha + \frac{\beta}{x - x_0} \quad \text{avec } \alpha = \frac{a}{c},\quad x_0 = -\frac{d}{c},\quad \beta = \frac{bc - ad}{c^2}. \]
Comme $ad - bc \neq 0$, on a bien $\beta \neq 0$.
\end{methode}

\begin{exemplebox}[Forme canonique de $f(x) = \dfrac{2x+3}{x-1}$]
On écrit le numérateur en fonction de $(x-1)$ :
\[ 2x + 3 = 2(x-1) + 2 + 3 = 2(x-1) + 5. \]
Donc
\[ f(x) = \frac{2(x-1)+5}{x-1} = 2 + \frac{5}{x-1}. \]
On reconnaît la fonction inverse $\dfrac{5}{x}$ translatée : la courbe de $f$ se déduit de l'hyperbole d'équation $y = \dfrac{5}{x}$ par la translation de vecteur $\vec{i} + 2\vec{j}$.
\end{exemplebox}

\courssub{Limites et asymptotes}

La forme canonique rend les limites presque évidentes. Pour $f(x) = \alpha + \dfrac{\beta}{x - x_0}$ :

\begin{propriete}[Asymptotes d'une fonction homographique]
Soit $f(x) = \dfrac{ax+b}{cx+d}$ avec $c \neq 0$ et $ad - bc \neq 0$. Alors :
\begin{itemize}
\item $\displaystyle\lim_{x \to -\frac{d}{c}^+} f(x) = \pm\infty$ et $\displaystyle\lim_{x \to -\frac{d}{c}^-} f(x) = \mp\infty$ : la droite d'équation $x = -\dfrac{d}{c}$ est \cle{asymptote verticale} à la courbe ;
\item $\displaystyle\lim_{x \to +\infty} f(x) = \lim_{x \to -\infty} f(x) = \dfrac{a}{c}$ : la droite d'équation $y = \dfrac{a}{c}$ est \cle{asymptote horizontale} à la courbe.
\end{itemize}
\end{propriete}

Pour la limite en la valeur interdite, le signe de l'infini dépend du signe de $\beta$ et du sens d'approche : il faut raisonner sur le signe du dénominateur $x - x_0$. Pour la limite en l'infini, on utilise la forme canonique : $\dfrac{\beta}{x - x_0} \to 0$, donc $f(x) \to \alpha = \dfrac{a}{c}$.

\begin{exemplebox}[Limites de $f(x) = \dfrac{2x+3}{x-1} = 2 + \dfrac{5}{x-1}$]
\begin{itemize}
\item Quand $x \to 1^+$ : $x - 1 \to 0^+$, donc $\dfrac{5}{x-1} \to +\infty$ et $f(x) \to +\infty$.
\item Quand $x \to 1^-$ : $x - 1 \to 0^-$, donc $\dfrac{5}{x-1} \to -\infty$ et $f(x) \to -\infty$.
\item Quand $x \to \pm\infty$ : $\dfrac{5}{x-1} \to 0$, donc $f(x) \to 2$.
\end{itemize}
Asymptote verticale : la droite d'équation $x = 1$ ; asymptote horizontale : la droite d'équation $y = 2$.
\end{exemplebox}

\courssub{Dérivée et sens de variation}

\begin{propriete}[Dérivée d'une fonction homographique — démonstration exigible au Probatoire]
Soit $f(x) = \dfrac{ax+b}{cx+d}$ avec $c \neq 0$ et $ad - bc \neq 0$. Alors $f$ est dérivable sur chaque intervalle de son domaine et
\[ f'(x) = \frac{ad - bc}{(cx+d)^2}. \]
Par conséquent, $f$ est \cle{strictement monotone} sur \emph{chacun} des intervalles de son domaine : strictement croissante si $ad - bc > 0$, strictement décroissante si $ad - bc < 0$.
\end{propriete}

\begin{experience}[Démonstration par calcul direct]
Posons $u(x) = ax + b$ et $v(x) = cx + d$. Alors $u'(x) = a$ et $v'(x) = c$, et $v$ ne s'annule pas sur $D_f$. Appliquons la formule de la dérivée d'un quotient $\left(\dfrac{u}{v}\right)' = \dfrac{u'v - uv'}{v^2}$ :
\begin{align*}
f'(x) &= \frac{a(cx+d) - c(ax+b)}{(cx+d)^2} \\
&= \frac{acx + ad - acx - bc}{(cx+d)^2} \\
&= \frac{ad - bc}{(cx+d)^2}.
\end{align*}
Les termes en $x$ disparaissent : le numérateur est une \emph{constante}. Or $(cx+d)^2 > 0$ pour tout $x \in D_f$. Le signe de $f'(x)$ est donc le signe de $ad - bc$, \emph{constant} sur chaque intervalle du domaine. D'où la stricte monotonie annoncée.
\end{experience}

\begin{attention}[Le piège classique du Probatoire]
On ne dit \textbf{jamais} qu'une fonction homographique est « décroissante sur son domaine ». Elle est décroissante (ou croissante) \emph{sur chacun des intervalles} $\left]-\infty ; -\frac{d}{c}\right[$ et $\left]-\frac{d}{c} ; +\infty\right[$ séparément. Contre-exemple flagrant : pour $f(x) = \dfrac{1}{x}$, on a $-1 < 1$ mais $f(-1) = -1 < f(1) = 1$ : si $f$ était décroissante sur tout $\mathbb{R}^*$, ce serait impossible ! Dans le tableau de variations, on place une \emph{double barre} en la valeur interdite et on ne relie jamais les deux branches par une flèche continue.
\end{attention}

Voici le tableau de variations type dans le cas $ad - bc < 0$ (alors $\beta = \dfrac{bc-ad}{c^2} > 0$, donc $f(x) \to +\infty$ quand $x \to x_0^+$ et $f(x) \to -\infty$ quand $x \to x_0^-$), avec $x_0 = -\dfrac{d}{c}$ et $\alpha = \dfrac{a}{c}$ :

\begin{center}
\begin{tabular}{|c|ccc||ccc|}\hline
$x$ & $-\infty$ & & $x_0$ & $x_0$ & & $+\infty$ \\ \hline
$f'(x)$ & & $-$ & & & $-$ & \\ \hline
$f$ & $\alpha$ & $\searrow$ & $-\infty$ & $+\infty$ & $\searrow$ & $\alpha$ \\ \hline
\end{tabular}
\end{center}

Dans le cas $ad - bc > 0$, les flèches montent ($\nearrow$) de part et d'autre de la double barre, avec les limites infinies échangées en $x_0$.

\courssub{Centre de symétrie}

\begin{propriete}[Centre de symétrie de l'hyperbole]
La courbe de $f(x) = \dfrac{ax+b}{cx+d}$ admet pour \cle{centre de symétrie} le point
\[ \Omega\left(-\frac{d}{c} \ ;\ \frac{a}{c}\right), \]
intersection des deux asymptotes. La courbe est une \cle{hyperbole}.
\end{propriete}

\begin{experience}[Preuve par la forme canonique]
Écrivons $f(x) = \alpha + \dfrac{\beta}{x - x_0}$ avec $x_0 = -\dfrac{d}{c}$ et $\alpha = \dfrac{a}{c}$. Pour montrer que $\Omega(x_0 ; \alpha)$ est centre de symétrie, vérifions que pour tout réel $h \neq 0$ :
\[ \frac{f(x_0 + h) + f(x_0 - h)}{2} = \alpha. \]
Or $f(x_0 + h) = \alpha + \dfrac{\beta}{h}$ et $f(x_0 - h) = \alpha - \dfrac{\beta}{h}$. Leur somme vaut $2\alpha$, donc la demi-somme vaut $\alpha$. Le point $\Omega$ est bien centre de symétrie de la courbe.
\end{experience}

\begin{aretenir}[L'essentiel sur les fonctions homographiques]
Pour $f(x) = \dfrac{ax+b}{cx+d}$ avec $c \neq 0$ et $ad - bc \neq 0$ :
\begin{itemize}
\item Domaine : $D_f = \mathbb{R} \setminus \left\{-\dfrac{d}{c}\right\}$ ;
\item Asymptote verticale : $x = -\dfrac{d}{c}$ ; asymptote horizontale : $y = \dfrac{a}{c}$ ;
\item Dérivée : $f'(x) = \dfrac{ad - bc}{(cx+d)^2}$, de signe constant ;
\item Centre de symétrie : $\Omega\left(-\dfrac{d}{c} ; \dfrac{a}{c}\right)$ ;
\item Forme canonique : $f(x) = \dfrac{a}{c} + \dfrac{bc - ad}{c(cx+d)}$.
\end{itemize}
\end{aretenir}

\courssub{Plan d'étude type et exemple complet}

\begin{methode}[Plan d'étude d'une fonction homographique en 6 étapes]
\begin{enumerate}
\item \textbf{Domaine} : résoudre $cx + d = 0$, écrire $D_f = \mathbb{R} \setminus \left\{-\frac{d}{c}\right\}$.
\item \textbf{Limites et asymptotes} : limites en la valeur interdite (à gauche et à droite) $\Rightarrow$ asymptote verticale ; limites en $\pm\infty$ $\Rightarrow$ asymptote horizontale $y = \frac{a}{c}$.
\item \textbf{Dérivée} : calculer $f'(x) = \dfrac{ad - bc}{(cx+d)^2}$ et donner son signe.
\item \textbf{Tableau de variations} : une double barre en la valeur interdite, flèches sur chaque intervalle, limites aux bornes.
\item \textbf{Symétrie} : citer le centre $\Omega\left(-\frac{d}{c} ; \frac{a}{c}\right)$.
\item \textbf{Tracé} : placer les asymptotes en pointillés, le centre $\Omega$, les intersections avec les axes ($f(0)$ et les solutions de $f(x) = 0$), quelques points auxiliaires, puis tracer les deux branches de l'hyperbole.
\end{enumerate}
\end{methode}

\begin{exoresolu}[Étude complète de $f(x) = \dfrac{2x+3}{x-1}$]
Soit $f$ la fonction définie par $f(x) = \dfrac{2x+3}{x-1}$ et $\mathcal{C}_f$ sa courbe dans un repère orthonormé.
\begin{enumerate}
\item Déterminer le domaine de définition de $f$.
\item Calculer les limites aux bornes du domaine et en déduire les asymptotes.
\item Étudier les variations de $f$ et dresser son tableau de variations.
\item Préciser le centre de symétrie de $\mathcal{C}_f$.
\item Déterminer les intersections de $\mathcal{C}_f$ avec les axes, puis tracer $\mathcal{C}_f$.
\end{enumerate}
\tcblower
\textbf{1. Domaine.} $x - 1 = 0 \iff x = 1$. Donc $D_f = \mathbb{R} \setminus \{1\} = \left]-\infty ; 1\right[ \cup \left]1 ; +\infty\right[$.

\textbf{2. Limites et asymptotes.} On utilise la forme canonique $f(x) = 2 + \dfrac{5}{x-1}$.
\begin{itemize}
\item $\displaystyle\lim_{x \to 1^+} f(x) = +\infty$ et $\displaystyle\lim_{x \to 1^-} f(x) = -\infty$ : la droite d'équation $x = 1$ est asymptote verticale.
\item $\displaystyle\lim_{x \to \pm\infty} f(x) = 2$ : la droite d'équation $y = 2$ est asymptote horizontale.
\end{itemize}

\textbf{3. Variations.} Avec $a = 2$, $b = 3$, $c = 1$, $d = -1$ : $ad - bc = 2\times(-1) - 3\times 1 = -5 < 0$. Donc
\[ f'(x) = \frac{-5}{(x-1)^2} < 0 \quad \text{sur chaque intervalle de } D_f. \]
$f$ est strictement décroissante sur $\left]-\infty ; 1\right[$ et sur $\left]1 ; +\infty\right[$.

\begin{center}
\begin{tabular}{|c|ccc||ccc|}\hline
$x$ & $-\infty$ & & $1$ & $1$ & & $+\infty$ \\ \hline
$f'(x)$ & & $-$ & & & $-$ & \\ \hline
$f$ & $2$ & $\searrow$ & $-\infty$ & $+\infty$ & $\searrow$ & $2$ \\ \hline
\end{tabular}
\end{center}

\textbf{4. Centre de symétrie.} $\Omega\left(1 ; 2\right)$, intersection des deux asymptotes.

\textbf{5. Intersections avec les axes.} Avec l'axe des ordonnées : $f(0) = \dfrac{3}{-1} = -3$, point $A(0 ; -3)$. Avec l'axe des abscisses : $f(x) = 0 \iff 2x + 3 = 0 \iff x = -\dfrac{3}{2}$, point $B\left(-\dfrac{3}{2} ; 0\right)$. On place aussi, par exemple, les points d'abscisses $2$ et $-1$ : $f(2) = 7$ et $f(-1) = -\dfrac{1}{2}$. On trace ensuite les deux branches de l'hyperbole (voir figure ci-dessous).
\end{exoresolu}

\begin{popfigure}
\begin{center}
\begin{tikzpicture}
\begin{axis}[
  width=0.85\linewidth,
  axis lines=middle,
  xlabel={$x$}, ylabel={$y$},
  xmin=-4, xmax=6, ymin=-6, ymax=9,
  xtick={-3,-2,-1,1,2,3,4,5},
  ytick={-4,-2,2,4,6,8},
  grid=both, grid style={popline!40},
  tick label style={font=\small}
]
% Asymptotes
\addplot[popmuted, dashed, domain=-4:6] {2};
\draw[popmuted, dashed] (axis cs:1,-6) -- (axis cs:1,9);
% Branches de l'hyperbole
\addplot[popcurve, domain=-4:0.7, samples=150] {(2*x+3)/(x-1)};
\addplot[popcurve, domain=1.35:6, samples=150] {(2*x+3)/(x-1)};
% Centre de symétrie
\addplot[only marks, mark=*, poporange, mark size=2pt] coordinates {(1,2)};
\node[poporange, anchor=south west, font=\small] at (axis cs:1.1,2.2) {$\Omega(1\,;\,2)$};
% Intersections avec les axes
\addplot[only marks, mark=*, popblue, mark size=2pt] coordinates {(0,-3) (-1.5,0)};
\node[popblue, anchor=north west, font=\small] at (axis cs:0.1,-3.2) {$A(0\,;\,-3)$};
\node[popblue, anchor=south east, font=\small] at (axis cs:-1.6,0.3) {$B\left(-\tfrac{3}{2}\,;\,0\right)$};
% Étiquettes asymptotes
\node[popmuted, font=\small, anchor=west] at (axis cs:4.2,2.4) {$y=2$};
\node[popmuted, font=\small, anchor=west, rotate=90] at (axis cs:0.85,7.2) {$x=1$};
\end{axis}
\end{tikzpicture}
\end{center}
\legende{Courbe de $f(x) = \dfrac{2x+3}{x-1}$ : hyperbole d'asymptotes $x = 1$ et $y = 2$, de centre de symétrie $\Omega(1\,;\,2)$.}
\end{popfigure}

\begin{attention}[Rédaction du tracé et erreurs fréquentes]
\begin{itemize}
\item Trace toujours les \emph{asymptotes en pointillés} avant de dessiner la courbe : elles guident le tracé et leur présence est exigée pour la note.
\item La courbe ne \emph{coupe jamais} l'asymptote verticale (valeur interdite). Pour une fonction homographique, elle ne coupe pas non plus l'asymptote horizontale : en effet $f(x) = \dfrac{a}{c}$ équivaut à $c(ax+b) = a(cx+d)$, c'est-à-dire $bc = ad$, ce qui est exclu par l'hypothèse $ad - bc \neq 0$. Attention toutefois : ce fait est propre aux homographiques ; d'autres fonctions peuvent croiser leur asymptote horizontale, ne le généralise pas.
\item N'oublie pas la \emph{double barre} dans le tableau de variations et ne joins jamais les deux branches à travers l'asymptote verticale.
\item Vérifie la cohérence : pour $f(x) = \dfrac{2x+3}{x-1}$, la branche de gauche passe par $A(0\,;\,-3)$ et $B\left(-\tfrac32\,;\,0\right)$ et tend vers $2$ en $-\infty$ ; si ton dessin contredit ces faits, il est faux.
\end{itemize}
\end{attention}

\begin{savaistu}[Les homographiques autour de nous]
Les fonctions homographiques apparaissent dès qu'une grandeur est partagée équitablement : le prix par personne d'un taxi-brousse Yaoundé--Douala dont le coût total est fixe est une fonction homographique du nombre de passagers. En physique, la relation des lentilles convergentes $\dfrac{1}{f'} = \dfrac{1}{p} + \dfrac{1}{p'}$ et bien des lois d'électrocinétique font intervenir des homographiques. Le mathématicien allemand August Ferdinand Möbius a étudié au XIX\textsuperscript{e} siècle leurs généralisations, appelées \emph{transformations de Möbius}, fondamentales en géométrie et en cartographie.
\end{savaistu}

\begin{exercice}[À toi de jouer]
Soit $g(x) = \dfrac{x - 2}{x + 1}$.
\begin{enumerate}
\item Déterminer le domaine de définition de $g$, ses limites aux bornes et ses asymptotes.
\item Montrer que $g(x) = 1 - \dfrac{3}{x+1}$ et en déduire le centre de symétrie de la courbe.
\item Étudier les variations de $g$ et dresser son tableau de variations.
\item Tracer la courbe de $g$ après avoir calculé les intersections avec les axes.
\end{enumerate}
\end{exercice}

\begin{corrige}[Exercice]
\textbf{1.} $x + 1 = 0 \iff x = -1$, donc $D_g = \mathbb{R} \setminus \{-1\}$. En $-1^+$ : $x+1 \to 0^+$ et $x - 2 \to -3 < 0$, donc $g(x) \to -\infty$ ; en $-1^-$ : $x+1 \to 0^-$, donc $g(x) \to +\infty$. Asymptote verticale : la droite d'équation $x = -1$. En $\pm\infty$ : $g(x) \to \dfrac{1}{1} = 1$, asymptote horizontale : la droite d'équation $y = 1$.

\textbf{2.} $1 - \dfrac{3}{x+1} = \dfrac{x+1-3}{x+1} = \dfrac{x-2}{x+1} = g(x)$. La courbe admet donc pour centre de symétrie $\Omega(-1 ; 1)$.

\textbf{3.} $ad - bc = 1\times 1 - (-2)\times 1 = 3 > 0$, donc $g'(x) = \dfrac{3}{(x+1)^2} > 0$ : $g$ est strictement croissante sur $\left]-\infty ; -1\right[$ et sur $\left]-1 ; +\infty\right[$.

\begin{center}
\begin{tabular}{|c|ccc||ccc|}\hline
$x$ & $-\infty$ & & $-1$ & $-1$ & & $+\infty$ \\ \hline
$g'(x)$ & & $+$ & & & $+$ & \\ \hline
$g$ & $1$ & $\nearrow$ & $+\infty$ & $-\infty$ & $\nearrow$ & $1$ \\ \hline
\end{tabular}
\end{center}

\textbf{4.} $g(0) = -2$ et $g(x) = 0 \iff x - 2 = 0 \iff x = 2$. La courbe passe par les points $(0\,;\,-2)$ et $(2\,;\,0)$, avec les asymptotes $x = -1$ et $y = 1$ tracées en pointillés, et elle est symétrique par rapport à $\Omega(-1\,;\,1)$.
\end{corrige}

\coursec{Étude des fonctions polynômes de degré $\leq 3$}

Dans la section précédente, nous avons étudié les fonctions homographiques : leurs courbes sont des hyperboles, dotées d'un centre de symétrie et de deux asymptotes. Nous quittons maintenant le monde des fractions pour celui des \cle{polynômes}, les fonctions les plus simples qui soient : elles sont définies sur $\mathbb{R}$ tout entier, continues et dérivables partout, sans aucune asymptote. Après les fonctions du second degré que tu connais depuis la Seconde, l'objectif de cette section est de maîtriser l'étude complète des fonctions \cle{polynômes de degré $3$}, appelées aussi \cle{fonctions cubiques}. C'est un classique absolu du Probatoire : presque chaque année, un problème propose l'étude d'une telle fonction.

\courssub{Rappels sur les fonctions polynômes de degré $2$}

\begin{definition}[Fonction polynôme de degré $2$]
Une fonction $f$ est un polynôme de degré $2$ s'il existe trois réels $a$, $b$, $c$ avec $a \neq 0$ tels que, pour tout réel $x$ :
\[ f(x) = ax^2 + bx + c. \]
Son ensemble de définition est $\mathbb{R}$.
\end{definition}

Tout polynôme de degré $2$ s'écrit sous sa \cle{forme canonique} :
\[ f(x) = a(x - \alpha)^2 + \beta \quad \text{avec} \quad \alpha = -\frac{b}{2a} \quad \text{et} \quad \beta = f(\alpha) = -\frac{\Delta}{4a}. \]

\begin{propriete}[Courbe d'un trinôme]
La courbe représentative de $f(x) = ax^2 + bx + c$ ($a \neq 0$) est une \cle{parabole} :
\begin{itemize}
\item de sommet $S\left(-\dfrac{b}{2a}\,;\, f\left(-\dfrac{b}{2a}\right)\right)$ ;
\item admettant la droite d'équation $x = -\dfrac{b}{2a}$ comme \cle{axe de symétrie} ;
\item tournée vers le haut si $a > 0$ ($f$ décroissante puis croissante), vers le bas si $a < 0$ ($f$ croissante puis décroissante).
\end{itemize}
\end{propriete}

\begin{exemplebox}[Forme canonique en pratique]
Soit $f(x) = 2x^2 - 8x + 5$. On factorise le coefficient de $x^2$ dans les deux premiers termes :
\[ f(x) = 2\left(x^2 - 4x\right) + 5 = 2\left[(x-2)^2 - 4\right] + 5 = 2(x-2)^2 - 3. \]
La parabole a pour sommet $S(2\,;\,-3)$ et pour axe de symétrie la droite d'équation $x = 2$.
\end{exemplebox}

\courssub{Fonction cubique : définition, dérivée, limites}

\begin{definition}[Fonction polynôme de degré $3$ ou cubique]
On appelle \cle{fonction cubique} toute fonction $f$ définie sur $\mathbb{R}$ par
\[ f(x) = ax^3 + bx^2 + cx + d, \quad \text{avec } a, b, c, d \text{ réels et } a \neq 0. \]
\end{definition}

Une cubique est définie, continue et dérivable sur $\mathbb{R}$ tout entier : il n'y a \textbf{aucune valeur interdite}, aucune asymptote verticale à chercher. Sa dérivée est un polynôme de degré $2$ :
\[ f'(x) = 3ax^2 + 2bx + c. \]
Toute l'étude des variations de $f$ repose donc sur l'étude du signe d'un trinôme, chose que tu sais faire depuis la Seconde grâce au discriminant.

\textbf{Limites aux bornes.} Que vaut $f(x)$ quand $x$ devient très grand ? Prenons l'intuition : dans $x^3 - 3x + 1$, quand $x$ vaut $1000$, le terme $x^3$ vaut un milliard, tandis que $-3x + 1$ vaut à peine $-2999$. Le terme de plus haut degré \og écrase \fg{} tous les autres. Cette intuition se justifie rigoureusement par une factorisation.

\begin{experience}[Pourquoi le terme de plus haut degré impose la limite]
Pour $x \neq 0$, factorisons $f(x) = ax^3 + bx^2 + cx + d$ par son terme dominant $ax^3$ :
\[ f(x) = ax^3\left(1 + \frac{b}{ax} + \frac{c}{ax^2} + \frac{d}{ax^3}\right). \]
Quand $x \to +\infty$, chacun des termes $\frac{b}{ax}$, $\frac{c}{ax^2}$, $\frac{d}{ax^3}$ tend vers $0$, donc la parenthèse tend vers $1$. Ainsi $f(x)$ se comporte comme $ax^3$ : même limite, même signe pour les grandes valeurs de $x$.
\end{experience}

\begin{propriete}[Limite d'un polynôme à l'infini (admise)]
En $+\infty$ et en $-\infty$, tout polynôme a la même limite que son \cle{terme de plus haut degré}. En particulier, pour $f(x) = ax^3 + bx^2 + cx + d$ avec $a > 0$ :
\[ \lim_{x \to -\infty} f(x) = -\infty \quad \text{et} \quad \lim_{x \to +\infty} f(x) = +\infty. \]
Si $a < 0$, les deux limites sont échangées. Une cubique n'a donc \textbf{jamais} d'asymptote horizontale.
\end{propriete}

\begin{attention}[Piège classique]
Ne conclus jamais que $\lim_{x \to +\infty} f(x) = +\infty$ en disant \og car $f$ est un polynôme \fg. La justification attendue est : \emph{un polynôme a en l'infini la limite de son terme de plus haut degré}, ici $ax^3$. Et attention au signe de $a$ : pour $f(x) = -2x^3 + x$, on a $\lim_{x \to +\infty} f(x) = -\infty$.
\end{attention}

\courssub{Discussion des variations selon le discriminant de $f'$}

Le signe de $f'(x) = 3ax^2 + 2bx + c$ dépend de son discriminant, que nous noterons
\[ \Delta' = (2b)^2 - 4 \times 3a \times c = 4b^2 - 12ac = 4(b^2 - 3ac). \]
Trois cas se présentent, qui donnent trois \og allures \fg{} typiques de cubiques.

\textbf{Cas 1 : $\Delta' > 0$.} $f'$ s'annule en deux réels distincts $x_1 < x_2$ et change de signe. La fonction admet \textbf{deux extrema locaux} : si $a > 0$, un maximum local en $x_1$ et un minimum local en $x_2$ ; si $a < 0$, c'est l'inverse. La courbe a une forme de \og S \fg{} couché.

\textbf{Cas 2 : $\Delta' = 0$.} $f'$ s'annule en un unique réel $x_0 = -\dfrac{b}{3a}$ \textbf{sans changer de signe} (trinôme de signe constant, nul en $x_0$). La fonction est strictement monotone sur $\mathbb{R}$, mais sa courbe admet en $x_0$ une \cle{tangente horizontale} : c'est un \cle{point d'inflexion à tangente horizontale}. Exemple : $f(x) = x^3$, dont la tangente en $O(0\,;\,0)$ est l'axe des abscisses, que la courbe traverse.

\textbf{Cas 3 : $\Delta' < 0$.} $f'$ ne s'annule jamais et garde un signe constant (celui de $3a$, soit le signe de $a$). La fonction est \textbf{strictement monotone} sur $\mathbb{R}$ : strictement croissante si $a > 0$, strictement décroissante si $a < 0$.

\begin{popfigure}
\begin{center}
\begin{tikzpicture}[scale=0.85]
% Cas 1 : Delta > 0
\begin{scope}[xshift=0cm]
\draw[->, popmuted] (-2.2,0) -- (2.2,0);
\draw[->, popmuted] (0,-1.8) -- (0,1.8);
\draw[popblue, very thick, domain=-1.9:1.9, samples=100] plot (\x, {0.35*(\x*\x*\x - 2.4*\x)});
\node[popdark, below] at (0,-2) {\textbf{$\Delta' > 0$ : deux extrema}};
\end{scope}
% Cas 2 : Delta = 0
\begin{scope}[xshift=5.2cm]
\draw[->, popmuted] (-2.2,0) -- (2.2,0);
\draw[->, popmuted] (0,-1.8) -- (0,1.8);
\draw[poporange, very thick, domain=-1.55:1.55, samples=100] plot (\x, {0.45*\x*\x*\x});
\node[popdark, below] at (0,-2) {\textbf{$\Delta' = 0$ : inflexion à tangente horizontale}};
\end{scope}
% Cas 3 : Delta < 0
\begin{scope}[xshift=10.4cm]
\draw[->, popmuted] (-2.2,0) -- (2.2,0);
\draw[->, popmuted] (0,-1.8) -- (0,1.8);
\draw[popgreen, very thick, domain=-1.7:1.7, samples=100] plot (\x, {0.3*(\x*\x*\x + 1.6*\x)});
\node[popdark, below] at (0,-2) {\textbf{$\Delta' < 0$ : monotone}};
\end{scope}
\end{tikzpicture}
\end{center}
\legende{Les trois allures possibles de la courbe d'une cubique (ici avec $a>0$), selon le signe du discriminant de $f'$.}
\end{popfigure}

\courssub{Centre de symétrie de la courbe d'une cubique}

Voici un fait remarquable : contrairement à la parabole qui possède un \emph{axe} de symétrie, la courbe d'une cubique possède toujours un \emph{centre} de symétrie. Observe la courbe de $f(x) = x^3$ : elle est symétrique par rapport à l'origine, car $f$ est impaire. Ce phénomène est général.

\begin{propriete}[Centre de symétrie d'une cubique (démonstration guidée)]
La courbe représentative de $f(x) = ax^3 + bx^2 + cx + d$ ($a \neq 0$) admet pour \cle{centre de symétrie} le point
\[ \Omega\left(-\frac{b}{3a}\,;\, f\left(-\frac{b}{3a}\right)\right). \]
C'est le point où la dérivée seconde $f''(x) = 6ax + 2b$ s'annule : on l'appelle \cle{point d'inflexion} de la courbe.
\end{propriete}

\begin{experience}[Démonstration guidée par changement de repère]
Posons $x_0 = -\dfrac{b}{3a}$ et $y_0 = f(x_0)$. Effectuons le changement de repère défini par
\[ x = x_0 + X \quad \text{et} \quad y = y_0 + Y, \]
c'est-à-dire que $\Omega$ devient la nouvelle origine. Développons $f(x_0 + X)$ :
\begin{align*}
f(x_0 + X) &= a(x_0 + X)^3 + b(x_0 + X)^2 + c(x_0 + X) + d \\
&= aX^3 + (3ax_0 + b)X^2 + \left(3ax_0^2 + 2bx_0 + c\right)X + f(x_0).
\end{align*}
Or $3ax_0 + b = 3a \times \left(-\dfrac{b}{3a}\right) + b = 0$ : le terme en $X^2$ disparaît. En notant $k = 3ax_0^2 + 2bx_0 + c = f'(x_0)$, il reste :
\[ Y = f(x_0 + X) - f(x_0) = aX^3 + kX. \]
La fonction $X \mapsto aX^3 + kX$ est \textbf{impaire} (elle ne contient que des puissances impaires de $X$). Sa courbe est donc symétrique par rapport à la nouvelle origine $\Omega$. Conclusion : dans le repère initial, la courbe de $f$ est symétrique par rapport au point $\Omega(x_0\,;\,y_0)$. $\blacksquare$
\end{experience}

\begin{attention}[Centre de symétrie, pas axe de symétrie]
Une erreur fréquente au Probatoire : affirmer que la droite $x = -\dfrac{b}{3a}$ est un axe de symétrie de la cubique. C'est \textbf{faux} : une cubique n'a jamais d'axe de symétrie. Pour prouver que $\Omega(\alpha\,;\,\beta)$ est centre de symétrie, on vérifie que pour tout réel $h$ :
\[ f(\alpha + h) + f(\alpha - h) = 2\beta. \]
\end{attention}

\courssub{Intersections avec les axes}

\begin{itemize}
\item \textbf{Avec l'axe des ordonnées} : immédiat, la courbe coupe $(Oy)$ au point $(0\,;\,d)$ puisque $f(0) = d$.
\item \textbf{Avec l'axe des abscisses} : il faut résoudre l'équation $ax^3 + bx^2 + cx + d = 0$. Il n'existe pas de formule simple au programme ; la stratégie est la recherche d'une \cle{racine évidente} parmi $-2, -1, 0, 1, 2$. Si $\alpha$ est racine, on factorise :
\[ f(x) = (x - \alpha)(ax^2 + b'x + c'), \]
et on termine avec le discriminant du trinôme. Une cubique a donc une, deux ou trois intersections avec l'axe des abscisses.
\end{itemize}

\courssub{Plan d'étude complet d'une fonction cubique}

\begin{methode}[Étudier et représenter $f(x) = ax^3 + bx^2 + cx + d$]
\begin{enumerate}
\item \textbf{Ensemble de définition} : $D_f = \mathbb{R}$.
\item \textbf{Limites} : en $-\infty$ et $+\infty$, $f$ a la limite de $ax^3$.
\item \textbf{Dérivée} : calculer $f'(x) = 3ax^2 + 2bx + c$, puis son discriminant $\Delta'$.
\item \textbf{Signe de $f'$ et variations} : discuter selon le signe de $\Delta'$ (deux extrema, inflexion à tangente horizontale, ou monotonie). Calculer les valeurs des extrema éventuels.
\item \textbf{Tableau de variations} complet, avec les limites.
\item \textbf{Centre de symétrie} : $\Omega\left(-\dfrac{b}{3a}\,;\, f\left(-\dfrac{b}{3a}\right)\right)$ ; éventuellement le justifier par $f(x_0 + h) + f(x_0 - h) = 2y_0$.
\item \textbf{Points clés} : intersection avec $(Oy)$, racines évidentes, tangentes remarquables (en $\Omega$, aux extrema).
\item \textbf{Tracé} : placer $\Omega$, les extrema, les points clés, puis la courbe en respectant le tableau.
\end{enumerate}
\end{methode}

\courssub{Exemple complet : étude de $f(x) = x^3 - 3x + 1$}

\begin{exemplebox}[Étude complète]
Soit $f$ la fonction définie par $f(x) = x^3 - 3x + 1$ et $(\mathcal{C})$ sa courbe dans un repère orthonormé.

\textbf{1. Ensemble de définition.} $f$ est un polynôme : $D_f = \mathbb{R}$.

\textbf{2. Limites.} $f$ a en l'infini la limite de son terme de plus haut degré $x^3$ :
\[ \lim_{x \to -\infty} f(x) = -\infty \quad \text{et} \quad \lim_{x \to +\infty} f(x) = +\infty. \]

\textbf{3. Dérivée et variations.} $f'(x) = 3x^2 - 3 = 3(x^2 - 1) = 3(x-1)(x+1)$.
Discriminant : $\Delta' = 0 - 4 \times 3 \times (-3) = 36 > 0$, racines $x_1 = -1$ et $x_2 = 1$. Comme le coefficient de $x^2$ dans $f'$ est positif, $f'(x) > 0$ sur $]-\infty\,;\,-1[$ et $]1\,;\,+\infty[$, et $f'(x) < 0$ sur $]-1\,;\,1[$.
Valeurs remarquables : $f(-1) = -1 + 3 + 1 = 3$ et $f(1) = 1 - 3 + 1 = -1$.

\textbf{4. Tableau de variations.}
\begin{center}
\begin{tabular}{|c|ccccccc|}
\hline
$x$ & $-\infty$ & & $-1$ & & $1$ & & $+\infty$ \\
\hline
$f'(x)$ & & $+$ & $0$ & $-$ & $0$ & $+$ & \\
\hline
$f$ & $-\infty$ & $\nearrow$ & $3$ & $\searrow$ & $-1$ & $\nearrow$ & $+\infty$ \\
\hline
\end{tabular}
\end{center}

\textbf{5. Centre de symétrie.} $x_0 = -\dfrac{b}{3a} = 0$ et $f(0) = 1$. Le point $\Omega(0\,;\,1)$ est centre de symétrie de $(\mathcal{C})$. Vérification :
\[ f(h) + f(-h) = \left(h^3 - 3h + 1\right) + \left(-h^3 + 3h + 1\right) = 2 = 2 \times 1. \]
La condition $f(0 + h) + f(0 - h) = 2f(0)$ est satisfaite pour tout $h$ : $\Omega$ est bien centre de symétrie. Remarquons la cohérence : les extrema $(-1\,;\,3)$ et $(1\,;\,-1)$ sont symétriques par rapport à $\Omega$ (milieu de leurs coordonnées : $\left(\frac{-1+1}{2}\,;\,\frac{3+(-1)}{2}\right) = (0\,;\,1)$).

\textbf{6. Points clés et tangentes.} $f(0) = 1$ : la courbe coupe $(Oy)$ en $(0\,;\,1)$. La tangente en $\Omega$ a pour coefficient directeur $f'(0) = -3$ : équation $y = -3x + 1$. Les tangentes en $x = -1$ et $x = 1$ sont horizontales. L'équation $f(x) = 0$ n'a pas de racine évidente entière ($f(1) = -1$, $f(-1) = 3$, $f(2) = 3$, $f(-2) = -1$) : on lira les solutions graphiquement.
\end{exemplebox}

\begin{popfigure}
\begin{center}
\begin{tikzpicture}
\begin{axis}[
  width=0.85\linewidth, height=7.5cm,
  axis lines=middle, axis line style={->, popmuted},
  xlabel={$x$}, ylabel={$y$},
  xmin=-2.6, xmax=2.6, ymin=-2.5, ymax=4.5,
  xtick={-2,-1,1,2}, ytick={-1,1,3},
  grid=both, grid style={popline!40, very thin},
  tick label style={font=\small, popmuted},
  label style={popdark}
]
\addplot[popcurve, domain=-2.25:2.25, samples=200]{x^3 - 3*x + 1};
% Tangente en Omega
\addplot[poppink, thick, dashed, domain=-0.9:1.1, samples=2]{-3*x + 1};
% Points remarquables
\addplot[only marks, mark=*, poporange, mark size=2pt] coordinates {(-1,3) (1,-1)};
\addplot[only marks, mark=*, poppurple, mark size=2.5pt] coordinates {(0,1)};
\node[poporange, above left, font=\small] at (axis cs:-1,3) {max $(-1\,;\,3)$};
\node[poporange, below right, font=\small] at (axis cs:1,-1) {min $(1\,;\,-1)$};
\node[poppurple, above right, font=\small] at (axis cs:0.08,1.15) {$\Omega(0\,;\,1)$};
\end{axis}
\end{tikzpicture}
\end{center}
\legende{Courbe de $f(x) = x^3 - 3x + 1$ : maximum local en $(-1\,;\,3)$, minimum local en $(1\,;\,-1)$, centre de symétrie $\Omega(0\,;\,1)$ et tangente en $\Omega$ (en pointillés).}
\end{popfigure}

\begin{exoresolu}[Lecture graphique : résolution de $x^3 - 3x + 1 = 0$]
On donne la courbe $(\mathcal{C})$ de $f(x) = x^3 - 3x + 1$ ci-dessus.
\begin{enumerate}
\item Déterminer graphiquement le nombre de solutions de l'équation $x^3 - 3x + 1 = 0$ et en donner un encadrement à $0{,}5$ près.
\item Résoudre graphiquement l'inéquation $x^3 - 3x + 1 \geq 0$.
\end{enumerate}
\tcblower
\textbf{1.} Les solutions de $f(x) = 0$ sont les abscisses des points d'intersection de $(\mathcal{C})$ avec l'axe des abscisses. La courbe coupe $(Ox)$ en \textbf{trois points} (cohérent avec le tableau : $f$ passe de $-\infty$ à $3$, puis de $3$ à $-1$, puis de $-1$ à $+\infty$, donc s'annule trois fois par continuité). Lecture graphique :
\[ -2 < x_1 < -1{,}5 \; ; \quad 0 < x_2 < 0{,}5 \; ; \quad 1{,}5 < x_3 < 2. \]
(En réalité $x_1 \approx -1{,}88$, $x_2 \approx 0{,}35$, $x_3 \approx 1{,}53$.)

\textbf{2.} L'inéquation $f(x) \geq 0$ revient à chercher où la courbe est \textbf{au-dessus} de l'axe des abscisses (ou le coupe). D'après le graphique :
\[ \mathcal{S} = [x_1\,;\,x_2] \cup [x_3\,;\,+\infty[ \approx [-1{,}88\,;\,0{,}35] \cup [1{,}53\,;\,+\infty[. \]
\textbf{Vérification par un point test} : $f(0) = 1 \geq 0$, et $0$ appartient bien à l'ensemble trouvé.
\end{exoresolu}

\begin{attention}[Les trois pièges du chapitre]
\begin{itemize}
\item \textbf{Oublier que $f'$ est un trinôme} : le signe de $f'(x) = 3ax^2 + 2bx + c$ se détermine avec son discriminant et la règle \og signe de $a$ à l'extérieur des racines \fg, pas au hasard.
\item \textbf{Mal calculer les extrema} : les valeurs $f(x_1)$ et $f(x_2)$ doivent être calculées avec soin ; ce sont elles qui apparaissent dans le tableau de variations, pas $x_1$ et $x_2$.
\item \textbf{Inventer des asymptotes} : une cubique n'a ni asymptote verticale (elle est définie sur $\mathbb{R}$) ni asymptote horizontale (sa limite en l'infini est infinie).
\end{itemize}
\end{attention}

\begin{savaistu}[Tartaglia, Cardan et la querelle du troisième degré]
Au XVI\textsuperscript{e} siècle, résoudre l'équation générale du troisième degré était le plus grand problème ouvert des mathématiques. Le mathématicien italien Niccolò Fontana, dit \emph{Tartaglia} (\og le bègue \fg), découvrit vers 1535 une méthode pour résoudre les équations de la forme $x^3 + px = q$ et gagna grâce à elle des concours publics de mathématiques, très prisés à l'époque. Jérôme Cardan, médecin et mathématicien milanais, lui arracha le secret sous serment de ne jamais le publier... puis le publia en 1545 dans son \emph{Ars Magna}, en citant Tartaglia mais en brisant sa promesse. S'ensuivit une querelle publique d'une violence extraordinaire, faite de pamphlets et de défis mathématiques. Ironie de l'histoire : ces recherches obligèrent les mathématiciens à manipuler des racines carrées de nombres négatifs, ouvrant la voie aux \cle{nombres complexes} que tu étudieras en Terminale. La formule de résolution est aujourd'hui connue sous le nom de \emph{formule de Cardan}.
\end{savaistu}

\begin{exercice}[Pour t'entraîner]
Soit $g$ la fonction définie sur $\mathbb{R}$ par $g(x) = -x^3 + 3x^2 - 2$.
\begin{enumerate}
\item Calculer les limites de $g$ en $-\infty$ et en $+\infty$.
\item Étudier les variations de $g$ et dresser son tableau de variations.
\item Montrer que le point $\Omega(1\,;\,0)$ est centre de symétrie de la courbe de $g$.
\item Déterminer une racine évidente de $g(x) = 0$, factoriser $g(x)$, puis résoudre l'équation $g(x) = 0$.
\end{enumerate}
\end{exercice}

\begin{corrige}[Exercice : étude de $g(x) = -x^3 + 3x^2 - 2$]
\textbf{1.} $g$ a en l'infini la limite de $-x^3$ : $\lim_{x \to -\infty} g(x) = +\infty$ et $\lim_{x \to +\infty} g(x) = -\infty$.

\textbf{2.} $g'(x) = -3x^2 + 6x = 3x(2 - x)$. Racines : $x = 0$ et $x = 2$. Comme le coefficient dominant de $g'$ est négatif, $g'(x) < 0$ sur $]-\infty\,;\,0[$ et $]2\,;\,+\infty[$, $g'(x) > 0$ sur $]0\,;\,2[$. Valeurs : $g(0) = -2$ et $g(2) = -8 + 12 - 2 = 2$.
\begin{center}
\begin{tabular}{|c|ccccccc|}
\hline
$x$ & $-\infty$ & & $0$ & & $2$ & & $+\infty$ \\
\hline
$g'(x)$ & & $-$ & $0$ & $+$ & $0$ & $-$ & \\
\hline
$g$ & $+\infty$ & $\searrow$ & $-2$ & $\nearrow$ & $2$ & $\searrow$ & $-\infty$ \\
\hline
\end{tabular}
\end{center}

\textbf{3.} $x_0 = -\dfrac{b}{3a} = -\dfrac{3}{-3} = 1$ et $g(1) = -1 + 3 - 2 = 0$. Vérifions : pour tout réel $h$,
\begin{align*}
g(1 + h) + g(1 - h) &= \left[-(1+h)^3 + 3(1+h)^2 - 2\right] + \left[-(1-h)^3 + 3(1-h)^2 - 2\right] \\
&= -\left[(1+h)^3 + (1-h)^3\right] + 3\left[(1+h)^2 + (1-h)^2\right] - 4 \\
&= -\left[2 + 6h^2\right] + 3\left[2 + 2h^2\right] - 4 = -2 - 6h^2 + 6 + 6h^2 - 4 = 0 = 2g(1).
\end{align*}
Donc $\Omega(1\,;\,0)$ est bien centre de symétrie de la courbe.

\textbf{4.} Racine évidente : $g(1) = 0$, donc $1$ est racine et $g(x) = (x - 1)(-x^2 + 2x + 2)$. Le trinôme $-x^2 + 2x + 2$ a pour discriminant $\Delta = 4 + 8 = 12$, d'où $x = \dfrac{-2 \pm 2\sqrt{3}}{-2} = 1 \mp \sqrt{3}$. Les solutions sont :
\[ \mathcal{S} = \left\{1 - \sqrt{3}\,;\, 1\,;\, 1 + \sqrt{3}\right\}. \]
On remarque que les trois racines sont symétriques par rapport à $1$, l'abscisse du centre de symétrie : ce n'est pas un hasard !
\end{corrige}

\coursec{Fonctions $x \mapsto \dfrac{ax^2+bx+c}{dx+e}$ et exploitation graphique}

Dans la section précédente, nous avons appris à étudier les fonctions polynômes de degré au plus $3$ : leur domaine est $\mathbb{R}$ tout entier, et leurs courbes ne présentent ni trou ni asymptote. Nous franchissons maintenant une étape décisive : que se passe-t-il lorsqu'un polynôme du second degré est \emph{divisé} par une fonction affine ? Le quotient introduit une valeur interdite, donc une asymptote verticale, et — surprise — une \cle{asymptote oblique} apparaît à l'infini. Ces fonctions, dernière grande famille au programme de Première, synthétisent tout ce que tu as appris : limites, dérivation, variations, symétries. Elles te permettront enfin d'\emph{exploiter} une courbe pour résoudre graphiquement équations et inéquations, compétence reine au Probatoire.

\courssub{Domaine de définition et asymptote verticale}

\begin{definition}[Fonction rationnelle du type étudié]
Soient $a$, $b$, $c$, $d$, $e$ des réels avec $a \neq 0$ et $d \neq 0$. On considère la fonction
\[ f(x) = \frac{ax^2 + bx + c}{dx + e}. \]
Le dénominateur s'annule en $x = -\dfrac{e}{d}$. Le \cle{domaine de définition} de $f$ est donc :
\[ D_f = \mathbb{R} \setminus \left\{ -\frac{e}{d} \right\} = \left] -\infty ; -\frac{e}{d} \right[ \cup \left] -\frac{e}{d} ; +\infty \right[. \]
\end{definition}

Comme pour les fonctions homographiques, lorsque le numérateur ne s'annule pas en $-\dfrac{e}{d}$ (cas général), la limite de $f$ en ce point est infinie, et la droite d'équation $x = -\dfrac{e}{d}$ est \cle{asymptote verticale} à la courbe $\mathcal{C}_f$. Le signe de l'infini ($-\infty$ ou $+\infty$) se détermine par une étude de signes à gauche et à droite de la valeur interdite.

\begin{attention}[Valeur interdite avant tout calcul]
Toute étude de ce type de fonction commence par l'annonce du domaine. Oublier la valeur interdite $-\dfrac{e}{d}$, c'est risquer d'inclure une valeur impossible dans un ensemble de solutions ou de tracer une courbe « continue » qui traverserait l'asymptote verticale : erreur éliminatoire au Probatoire.
\end{attention}

\courssub{Division euclidienne et asymptote oblique}

Voici l'idée directrice. Quand $x$ devient très grand en valeur absolue, le quotient $\dfrac{ax^2+bx+c}{dx+e}$ se comporte comme $\dfrac{ax^2}{dx} = \dfrac{a}{d}x$ : la courbe « suit » une droite de pente $\dfrac{a}{d}$. Pour exhiber cette droite, on effectue la \cle{division euclidienne} du numérateur par le dénominateur.

\begin{propriete}[Forme réduite et asymptote oblique]
Il existe des réels $\alpha$, $\beta$, $\gamma$ (avec $\alpha = \dfrac{a}{d} \neq 0$) tels que, pour tout $x \in D_f$ :
\[ f(x) = \alpha x + \beta + \frac{\gamma}{dx+e}. \]
Comme $\displaystyle\lim_{x \to \pm\infty} \frac{\gamma}{dx+e} = 0$, on a $\displaystyle\lim_{x \to \pm\infty} \big[ f(x) - (\alpha x + \beta) \big] = 0$ : la droite d'équation $y = \alpha x + \beta$ est \cle{asymptote oblique} à $\mathcal{C}_f$ en $-\infty$ et en $+\infty$.
\end{propriete}

\begin{methode}[Obtenir la forme réduite]
Pour écrire $f(x) = \alpha x + \beta + \dfrac{\gamma}{dx+e}$ :
\begin{enumerate}
\item Poser la division euclidienne de $ax^2+bx+c$ par $dx+e$ : quotient $\alpha x + \beta$, reste $\gamma$ (constant car le diviseur est de degré $1$).
\item Vérifier en réduisant au même dénominateur : $(\alpha x + \beta)(dx+e) + \gamma$ doit redonner $ax^2+bx+c$.
\item Conclure : l'asymptote oblique est $y = \alpha x + \beta$. Le signe de $f(x) - (\alpha x + \beta) = \dfrac{\gamma}{dx+e}$ donne en outre la \emph{position relative} de la courbe par rapport à son asymptote.
\end{enumerate}
\end{methode}

\courssub{Dérivée et discussion des variations}

\begin{propriete}[Dérivée]
La fonction $f$ est dérivable sur $D_f$ et, par la formule du quotient $\left(\dfrac{u}{v}\right)' = \dfrac{u'v - uv'}{v^2}$ :
\[ f'(x) = \frac{(2ax+b)(dx+e) - d(ax^2+bx+c)}{(dx+e)^2} = \frac{ad\,x^2 + 2ae\,x + (be - cd)}{(dx+e)^2}. \]
Le numérateur est un \cle{polynôme de degré 2} (car $ad \neq 0$). Le signe de $f'(x)$ est celui de ce numérateur, le dénominateur étant un carré strictement positif sur $D_f$.
\end{propriete}

On calcule donc le discriminant $\Delta$ du numérateur $N(x) = ad\,x^2 + 2ae\,x + (be-cd)$ :
\begin{itemize}
\item si $\Delta \leq 0$ : $f'$ garde un signe constant sur chaque intervalle du domaine ; $f$ est monotone sur chacun d'eux (pas d'extremum) ;
\item si $\Delta > 0$ : $f'$ s'annule en deux réels $x_1 < x_2$ (distincts de $-\dfrac{e}{d}$ en général) ; $f$ admet deux \cle{extrema locaux}, l'un sur chaque branche de la courbe.
\end{itemize}

\begin{attention}[Ne pas confondre les deux discriminants]
Le discriminant qui gouverne les \emph{variations} est celui du numérateur de $f'$, pas celui du numérateur $ax^2+bx+c$ de $f$ (qui, lui, sert à trouver les intersections avec l'axe des abscisses). Note bien la différence dans ta copie.
\end{attention}

\courssub{Centre de symétrie}

\begin{propriete}[Centre de symétrie — admise]
La courbe de $f(x) = \dfrac{ax^2+bx+c}{dx+e}$ admet pour \cle{centre de symétrie} le point $\Omega$, intersection de ses deux asymptotes :
\[ \Omega\left( -\frac{e}{d} \;;\; -\frac{\alpha e}{d} + \beta \right) \quad \text{où } y = \alpha x + \beta \text{ est l'asymptote oblique.} \]
Ce résultat est \emph{admis} ; on le vérifiera sur l'exemple de référence en montrant que $f(\Omega_x - h) + f(\Omega_x + h) = 2\,\Omega_y$ pour tout $h \neq 0$.
\end{propriete}

\courssub{Exemple de référence : étude complète de $f(x) = \dfrac{x^2 - x + 3}{x - 1}$}

\begin{exemplebox}[Étude complète et tracé]
Étudions pas à pas $f(x) = \dfrac{x^2-x+3}{x-1}$.

\textbf{1. Domaine.} $x - 1 = 0 \iff x = 1$, donc $D_f = \mathbb{R} \setminus \{1\}$.

\textbf{2. Limites et asymptote verticale.} En $1$ : le numérateur vaut $1 - 1 + 3 = 3 > 0$. Pour $x < 1$, le dénominateur est négatif, donc $\displaystyle\lim_{\substack{x \to 1 \\ x < 1}} f(x) = -\infty$ ; pour $x > 1$, il est positif, donc $\displaystyle\lim_{\substack{x \to 1 \\ x > 1}} f(x) = +\infty$. La droite $x = 1$ est asymptote verticale. En $\pm\infty$ : $f(x) \sim \dfrac{x^2}{x} = x$, donc $\displaystyle\lim_{x \to -\infty} f(x) = -\infty$ et $\displaystyle\lim_{x \to +\infty} f(x) = +\infty$.

\textbf{3. Division euclidienne.} $x^2 - x + 3 = x(x-1) + 3$, d'où
\[ f(x) = x + \frac{3}{x-1}. \]
La droite d'équation $y = x$ est asymptote oblique en $-\infty$ et en $+\infty$. De plus, $f(x) - x = \dfrac{3}{x-1}$ : pour $x > 1$ la courbe est \emph{au-dessus} de l'asymptote, pour $x < 1$ elle est \emph{en dessous}.

\textbf{4. Dérivée.} À partir de la forme réduite : $f'(x) = 1 - \dfrac{3}{(x-1)^2} = \dfrac{(x-1)^2 - 3}{(x-1)^2}$. Le numérateur s'annule pour $(x-1)^2 = 3$, soit $x = 1 - \sqrt{3}$ ou $x = 1 + \sqrt{3}$. C'est un trinôme en $(x-1)$, positif à l'extérieur des racines, négatif entre elles.

\textbf{5. Extrema.}
\begin{align*}
f(1-\sqrt{3}) &= (1-\sqrt{3}) + \frac{3}{-\sqrt{3}} = 1 - \sqrt{3} - \sqrt{3} = 1 - 2\sqrt{3} \approx -2{,}46 \quad (\text{maximum local})\\
f(1+\sqrt{3}) &= (1+\sqrt{3}) + \frac{3}{\sqrt{3}} = 1 + \sqrt{3} + \sqrt{3} = 1 + 2\sqrt{3} \approx 4{,}46 \quad (\text{minimum local})
\end{align*}
(On a utilisé $\dfrac{3}{\sqrt{3}} = \sqrt{3}$.)

\textbf{6. Tableau de variations.}
\begin{center}
\begin{tabular}{|c|ccccccccc|}\hline
$x$ & $-\infty$ & & $1-\sqrt{3}$ & & $1$ & & $1+\sqrt{3}$ & & $+\infty$ \\ \hline
$f'(x)$ & & $+$ & $0$ & $-$ & & $-$ & $0$ & $+$ & \\ \hline
$f$ & $-\infty$ & $\nearrow$ & $1-2\sqrt{3}$ & $\searrow$ & $-\infty \;\big|\; +\infty$ & $\searrow$ & $1+2\sqrt{3}$ & $\nearrow$ & $+\infty$ \\ \hline
\end{tabular}
\end{center}

\textbf{7. Centre de symétrie (vérification).} Les asymptotes $x=1$ et $y=x$ se coupent en $\Omega(1\,;\,1)$. Vérifions : pour tout $h \neq 0$,
\[ f(1-h) + f(1+h) = \left(1 - h - \frac{3}{h}\right) + \left(1 + h + \frac{3}{h}\right) = 2 = 2 \times 1. \]
$\Omega(1\,;\,1)$ est bien centre de symétrie de $\mathcal{C}_f$.

\textbf{8. Points remarquables.} $f(0) = -3$ (intersection avec l'axe des ordonnées). Le numérateur $x^2 - x + 3$ a pour discriminant $1 - 12 = -11 < 0$ : la courbe ne coupe \emph{jamais} l'axe des abscisses. Les tangentes en $1-\sqrt{3}$ et $1+\sqrt{3}$ sont horizontales.
\end{exemplebox}

\begin{popfigure}
\begin{center}
\begin{tikzpicture}
\begin{axis}[
  width=0.9\linewidth, height=9cm,
  axis lines=middle, xlabel={$x$}, ylabel={$y$},
  xmin=-3.5, xmax=5.5, ymin=-5, ymax=7,
  xtick={-3,-2,-1,1,2,3,4,5}, ytick={-4,-2,2,4,6},
  grid=both, grid style={popline!40}
]
\addplot[popcurve, domain=-3.4:0.82, samples=150] {x + 3/(x-1)};
\addplot[popcurve, domain=1.18:5.4, samples=150] {x + 3/(x-1)};
\addplot[dashed, poporange, thick, domain=-3.4:5.4] {x};
\draw[dashed, poppink, thick] (axis cs:1,-5) -- (axis cs:1,7);
\addplot[only marks, mark=*, popdark, mark size=2pt] coordinates {(-0.732,-2.464) (2.732,4.464) (0,-3)};
\addplot[only marks, mark=square*, popgreen, mark size=2.5pt] coordinates {(1,1)};
\node[popgreen, font=\small] at (axis cs:1.6,0.4) {$\Omega(1\,;\,1)$};
\node[popdark, font=\small, anchor=east] at (axis cs:-0.9,-2.9) {$1-2\sqrt{3}$};
\node[popdark, font=\small, anchor=west] at (axis cs:2.9,4.1) {$1+2\sqrt{3}$};
\node[poporange, font=\small, anchor=north west, rotate=38] at (axis cs:4.1,4.35) {$y=x$};
\node[poppink, font=\small, anchor=west] at (axis cs:1.1,-4.6) {$x=1$};
\end{axis}
\end{tikzpicture}
\end{center}
\legende{Courbe de $f(x)=\dfrac{x^2-x+3}{x-1}$, asymptotes $x=1$ (rose) et $y=x$ (orange), extrema et centre de symétrie $\Omega$.}
\end{popfigure}

\courssub{Exploitation graphique : équations et inéquations}

Une fois la courbe tracée avec soin, elle devient un \emph{instrument de résolution}. Résoudre graphiquement, c'est lire des intersections et des positions relatives — sans aucun calcul algébrique.

\begin{methode}[Résoudre graphiquement $f(x) = k$ ou $f(x) \geq k$]
\begin{enumerate}
\item Tracer la droite horizontale d'équation $y = k$.
\item Pour $f(x) = k$ : lire les \cle{abscisses} des points d'intersection de $\mathcal{C}_f$ avec cette droite. L'ensemble des solutions est la liste de ces abscisses.
\item Pour $f(x) \geq k$ : repérer les intervalles où la courbe est \emph{au-dessus} de la droite (ou la touche), puis lire les abscisses correspondantes sur l'axe des $x$.
\item Exclure systématiquement les valeurs interdites de $D_f$.
\end{enumerate}
On procède de même pour $f(x) = mx + p$ avec la droite $y = mx + p$, et pour $f(x) \leq g(x)$ en comparant les positions de deux courbes.
\end{methode}

\begin{exoresolu}[Lecture graphique sur $\mathcal{C}_f$]
On conserve $f(x) = \dfrac{x^2-x+3}{x-1}$ et sa courbe tracée ci-dessus.

\textbf{1.} Résoudre graphiquement $f(x) = 3$, puis $f(x) \geq 3$.

\textbf{2.} Résoudre graphiquement $f(x) = 5$.

\tcblower

\textbf{1.} La droite $y = 3$ ne rencontre $\mathcal{C}_f$ en \emph{aucun} point : la branche de gauche culmine à $1 - 2\sqrt{3} \approx -2{,}46 < 3$ et la branche de droite descend au plus bas à $1 + 2\sqrt{3} \approx 4{,}46 > 3$. Donc $S = \emptyset$.

Pour $f(x) \geq 3$ : la courbe est au-dessus de la droite $y = 3$ uniquement sur sa branche de droite, c'est-à-dire pour $x > 1$ (la valeur $1$ est interdite). D'où $S = \left]1\,;\,+\infty\right[$. Vérification algébrique : $f(x) - 3 = \dfrac{x^2 - x + 3 - 3(x-1)}{x-1} = \dfrac{x^2 - 4x + 6}{x - 1}$ ; le numérateur a pour discriminant $16 - 24 = -8 < 0$, il est donc toujours strictement positif, et le signe du quotient est celui de $x - 1$. Confirmé.

\textbf{2.} La droite $y = 5$ coupe la branche de droite en deux points, dont on lit les abscisses $x = 2$ et $x = 4$. Vérification : $f(2) = \dfrac{4-2+3}{1} = 5$ et $f(4) = \dfrac{16-4+3}{3} = 5$. Donc $S = \{2\,;\,4\}$.
\end{exoresolu}

\begin{popfigure}
\begin{center}
\begin{tikzpicture}
\begin{axis}[
  width=0.9\linewidth, height=8.5cm,
  axis lines=middle, xlabel={$x$}, ylabel={$y$},
  xmin=-3.5, xmax=5.5, ymin=-5, ymax=7,
  xtick={-3,-2,-1,1,2,3,4,5}, ytick={-4,-2,2,3,4,6},
  grid=both, grid style={popline!40}
]
\addplot[popcurve, domain=-3.4:0.85, samples=150] {x + 3/(x-1)};
\addplot[popcurve, popgreen, very thick, domain=1.15:5.4, samples=150] {x + 3/(x-1)};
\addplot[dashed, poppink, thick, domain=-3.4:5.4] {3};
\draw[dashed, poporange, thick] (axis cs:1,-5) -- (axis cs:1,7);
\addplot[only marks, mark=*, popgold, mark size=2.5pt] coordinates {(2,5) (4,5)};
\addplot[dashed, popgold, domain=-3.4:5.4] {5};
\node[popgreen, font=\small, anchor=west] at (axis cs:3.2,6.3) {zone $f(x) \geq 3$ : $x > 1$};
\node[poppink, font=\small, anchor=south west] at (axis cs:-3.3,3.1) {$y=3$};
\node[popgold, font=\small, anchor=north west] at (axis cs:-3.3,4.9) {$y=5$};
\end{axis}
\end{tikzpicture}
\end{center}
\legende{Résolution graphique : $y=3$ ne coupe pas $\mathcal{C}_f$ ($f(x)=3$ sans solution) ; la branche verte au-dessus de $y=3$ donne $f(x)\geq 3 \iff x>1$ ; $y=5$ coupe la courbe en $x=2$ et $x=4$.}
\end{popfigure}

\begin{attention}[Toujours exclure les valeurs interdites]
Dans une résolution graphique d'inéquation, une borne d'intervalle solution qui coïncide avec la valeur interdite doit être \emph{ouverte}. Ici, $f(x) \geq 3$ a pour solution $]1\,;\,+\infty[$ et non $[1\,;\,+\infty[$ : $f(1)$ n'existe pas. De même, une intersection située exactement sur l'asymptote verticale est une illusion de lecture : la courbe ne la touche jamais.
\end{attention}

\courssub{Synthèse : plan général d'étude}

\begin{methode}[Plan d'étude complet en 7 étapes de $f(x) = \dfrac{ax^2+bx+c}{dx+e}$]
\begin{enumerate}
\item \textbf{Domaine} : $D_f = \mathbb{R} \setminus \left\{-\dfrac{e}{d}\right\}$.
\item \textbf{Limites aux bornes} : en $-\dfrac{e}{d}$ (à gauche et à droite, en étudiant les signes) et en $\pm\infty$ (par les termes de plus haut degré).
\item \textbf{Asymptotes} : verticale $x = -\dfrac{e}{d}$ ; division euclidienne donnant $f(x) = \alpha x + \beta + \dfrac{\gamma}{dx+e}$, d'où l'asymptote oblique $y = \alpha x + \beta$ et la position relative courbe/asymptote (signe de $\dfrac{\gamma}{dx+e}$).
\item \textbf{Dérivée} : formule du quotient, numérateur trinôme, discriminant, signe.
\item \textbf{Tableau de variations} : complet, avec limites aux bornes, double barre en la valeur interdite, extrema éventuels.
\item \textbf{Éléments de symétrie et points remarquables} : centre $\Omega$ (intersection des asymptotes), intersections avec les axes, tangentes horizontales.
\item \textbf{Tracé} : asymptotes d'abord, extrema et points d'appui, puis les deux branches en respectant le tableau.
\end{enumerate}
\end{methode}

\begin{savaistu}[Des asymptotes obliques dans la vie réelle]
Les courbes à asymptote oblique ne sont pas qu'un exercice de Probatoire ! En économie, le \emph{coût moyen} de production d'une entreprise (par exemple une minoterie à Douala) s'écrit $C_M(x) = \dfrac{ax^2+bx+c}{x}$ : quand la production $x$ devient grande, il se rapproche d'une droite — l'asymptote oblique — qui représente le coût unitaire stabilisé. En informatique, le temps moyen de réponse d'un serveur (chez CAMTEL ou Orange) en fonction du nombre de requêtes suit des modèles analogues : l'asymptote oblique traduit la saturation progressive du système. Comprendre ces courbes, c'est comprendre où un système « s'essouffle ».
\end{savaistu}

\begin{exercice}[Pour t'entraîner]
Soit $g(x) = \dfrac{x^2 + 2x - 2}{x + 1}$.
\begin{enumerate}
\item Déterminer $D_g$ et les limites aux bornes.
\item Montrer que $g(x) = x + 1 - \dfrac{3}{x+1}$ et en déduire les asymptotes.
\item Étudier les variations de $g$ et dresser son tableau de variations.
\item Préciser le centre de symétrie, puis tracer la courbe.
\item Résoudre graphiquement $g(x) = 1$ puis $g(x) \leq 1$.
\end{enumerate}
\end{exercice}

\begin{corrige}[Exercice : fonction $g$]
\textbf{1.} $x + 1 = 0 \iff x = -1$, donc $D_g = \mathbb{R} \setminus \{-1\}$. En $-1$ : le numérateur vaut $1 - 2 - 2 = -3 < 0$ ; pour $x < -1$ le dénominateur est négatif, donc $\displaystyle\lim_{\substack{x \to -1 \\ x < -1}} g(x) = +\infty$ ; pour $x > -1$ il est positif, donc $\displaystyle\lim_{\substack{x \to -1 \\ x > -1}} g(x) = -\infty$. En $\pm\infty$ : $g(x) \sim \dfrac{x^2}{x} = x$, donc $\displaystyle\lim_{x \to -\infty} g(x) = -\infty$ et $\displaystyle\lim_{x \to +\infty} g(x) = +\infty$.

\textbf{2.} $(x+1)(x+1) = x^2 + 2x + 1$, donc $x^2 + 2x - 2 = (x+1)^2 - 3$ et $g(x) = x + 1 - \dfrac{3}{x+1}$. Asymptote verticale : $x = -1$ ; asymptote oblique : $y = x + 1$. Comme $g(x) - (x+1) = -\dfrac{3}{x+1}$, la courbe est en dessous de l'asymptote pour $x > -1$ et au-dessus pour $x < -1$.

\textbf{3.} $g'(x) = 1 + \dfrac{3}{(x+1)^2} = \dfrac{(x+1)^2 + 3}{(x+1)^2} > 0$ sur $D_g$ (le numérateur, somme d'un carré et de $3$, ne s'annule jamais : cas $\Delta < 0$). Donc $g$ est strictement croissante sur $]-\infty\,;\,-1[$ et sur $]-1\,;\,+\infty[$, sans extremum.

\textbf{4.} Centre de symétrie : intersection de $x = -1$ et $y = x+1$, soit $\Omega(-1\,;\,0)$. Vérification : pour tout $h \neq 0$,
\[ g(-1-h) + g(-1+h) = \left(-h + \frac{3}{h}\right) + \left(h - \frac{3}{h}\right) = 0 = 2 \times 0. \]
Intersection avec $(Oy)$ : $g(0) = -2$ ; avec $(Ox)$ : $x^2 + 2x - 2 = 0$ donne $x = -1 - \sqrt{3} \approx -2{,}73$ ou $x = -1 + \sqrt{3} \approx 0{,}73$.

\textbf{5.} $g(x) = 1 \iff x + 1 - \dfrac{3}{x+1} = 1 \iff x(x+1) = 3 \iff x^2 + x - 3 = 0$, soit $x = \dfrac{-1 \pm \sqrt{13}}{2}$, c'est-à-dire $x \approx -2{,}30$ ou $x \approx 1{,}30$ : $S = \left\{\dfrac{-1-\sqrt{13}}{2}\,;\,\dfrac{-1+\sqrt{13}}{2}\right\}$.

Pour $g(x) \leq 1$ : $g(x) - 1 = \dfrac{x^2+x-3}{x+1}$. Le numérateur est positif à l'extérieur de ses racines. Un tableau de signes (ou la lecture de la courbe sous la droite $y = 1$) donne :
\[ S = \left]-\infty\,;\,\frac{-1-\sqrt{13}}{2}\right] \cup \left]-1\,;\,\frac{-1+\sqrt{13}}{2}\right] \]
(la borne $-1$ est exclue : valeur interdite ; les racines sont incluses car l'inéquation est large).
\end{corrige}

\newpage

\coursec{Activité d'intégration}

\courssub{Objectifs de l'activité}

À l'issue de cette activité d'intégration, tu seras capable de :

\begin{itemize}
\item déterminer l'ensemble de définition d'une fonction rationnelle et justifier la cohérence du domaine avec le contexte concret ;
\item déterminer par leurs équations les asymptotes parallèles aux axes (asymptote verticale) et montrer qu'une droite est asymptote oblique à une courbe, en utilisant la division euclidienne du numérateur par le dénominateur ;
\item étudier les variations d'une fonction du type $x \mapsto \dfrac{ax^2+bx+c}{dx+e}$ sur un intervalle et dresser son tableau de variations complet ;
\item mettre en évidence un élément de symétrie d'une courbe (ici un centre de symétrie) et l'exploiter pour construire la représentation graphique ;
\item exploiter la courbe d'une fonction pour résoudre graphiquement puis algébriquement une équation et une inéquation issues d'une contrainte réelle ;
\item interpréter économiquement les résultats mathématiques (coût minimal, seuil de rentabilité, comportement asymptotique).
\end{itemize}

\courssub{Situation}

La coopérative \og Mbog Cacao \fg{}, située dans la région du Centre, regroupe 240 planteurs des villages environnants d'Obala. Chaque année, la coopérative collecte, fermente et sèche le cacao avant de le revendre à un exportateur de Douala. Le gérant, M.~Essomba, a constaté un phénomène paradoxal : lorsque la production augmente trop, le coût moyen de production par tonne \emph{augmente} au lieu de diminuer (surcharge du séchoir, location de camions supplémentaires, heures supplémentaires des employés).

L'économiste de la coopérative a modélisé la situation. Pour une production de $x$ tonnes de cacao marchand, le \textbf{coût moyen de production par tonne}, exprimé en \textbf{milliers de FCFA}, est donné par :
\[ f(x) = \frac{2x^2 + 3x + 9}{x + 1}, \qquad x > 0. \]

L'exportateur de Douala impose une contrainte budgétaire : il ne signe le contrat d'achat que si le coût moyen de production ne dépasse pas \textbf{8 milliers de FCFA par tonne} (soit \num{8000} FCFA), condition nécessaire pour que le prix de vente couvre les frais et dégage une marge.

M.~Essomba pose alors à la coopérative trois questions :
\begin{enumerate}
\item Quelle quantité de cacao faut-il produire pour que le coût moyen par tonne soit \textbf{minimal} ? Quel est ce coût minimal ?
\item Pour quelles productions la contrainte de l'exportateur (coût moyen $\leq 8$ milliers de FCFA) est-elle respectée ?
\item Que prévoit le modèle pour les très grandes productions ? Faut-il accepter toutes les récoltes sans limite ?
\end{enumerate}

Par groupes de quatre, vous êtes les experts-comptables stagiaires de la coopérative. À vous de répondre, courbe à l'appui !

\courssub{Consignes}

\textbf{Partie A : Domaine et asymptotes}
\begin{enumerate}
\item Déterminer l'ensemble de définition $D_f$ de $f$ dans $\mathbb{R}$. Expliquer pourquoi, dans le contexte, on ne s'intéresse qu'à l'intervalle $]0\,;\,+\infty[$.
\item Étudier la limite de $f$ en $-1$ (valeur interdite). En déduire une asymptote verticale à la courbe $\mathcal{C}_f$. Cette asymptote a-t-elle un sens économique ?
\item Effectuer la division euclidienne de $2x^2+3x+9$ par $x+1$. En déduire trois réels $a$, $b$, $c$ tels que $f(x) = ax + b + \dfrac{c}{x+1}$.
\item Montrer que la droite $\Delta$ d'équation $y = 2x+1$ est asymptote oblique à $\mathcal{C}_f$ en $+\infty$. Préciser la position relative de $\mathcal{C}_f$ et de $\Delta$ sur $]0\,;\,+\infty[$.
\end{enumerate}

\textbf{Partie B : Variations et symétrie}
\begin{enumerate}[resume]
\item Calculer $f'(x)$ et montrer que $f'(x) = 2 - \dfrac{8}{(x+1)^2}$.
\item Résoudre $f'(x) = 0$ sur $]0\,;\,+\infty[$, puis dresser le tableau de variations complet de $f$ sur cet intervalle.
\item Montrer que le point $\Omega(-1\,;\,-1)$ est centre de symétrie de la courbe complète de $f$ (on calculera $f(-1+h)+f(-1-h)$).
\end{enumerate}

\textbf{Partie C : Représentation graphique et exploitation}
\begin{enumerate}[resume]
\item Sur papier millimétré, tracer $\mathcal{C}_f$ pour $x \in [0\,;\,10]$ (unités : 1~cm pour 1 tonne en abscisse, 1~cm pour 2 milliers de FCFA en ordonnée), ainsi que l'asymptote $\Delta$.
\item Déterminer graphiquement la quantité qui minimise le coût moyen, puis confirmer par le calcul.
\item Résoudre graphiquement l'inéquation $f(x) \leq 8$, puis la résoudre algébriquement. Conclure sur la contrainte de l'exportateur.
\item Interpréter l'asymptote oblique : que signifie-t-elle pour les très grandes productions ? Rédiger une recommandation de trois lignes à l'attention de M.~Essomba.
\end{enumerate}

Chaque groupe présente sa courbe au tableau et défend ses conclusions devant la classe.

\courssub{Ressources à mobiliser}

\begin{aretenir}[Boîte à outils de l'activité]
\begin{itemize}
\item \cle{Asymptote verticale} : si $\lim\limits_{x \to x_0} f(x) = \pm\infty$, la droite $x = x_0$ est asymptote verticale à $\mathcal{C}_f$.
\item \cle{Asymptote oblique} : si $f(x) = ax+b+\dfrac{c}{x - x_0}$ avec $c \neq 0$, alors $\lim\limits_{x \to \pm\infty} \big[f(x) - (ax+b)\big] = 0$ : la droite $y = ax+b$ est asymptote oblique. Le signe de $\dfrac{c}{x-x_0}$ donne la position relative.
\item \cle{Division euclidienne} de polynômes : pour tout polynôme $P$ et tout binôme $x - x_0$, il existe $Q$ et $R$ avec $P(x) = (x - x_0)Q(x) + R$.
\item \cle{Dérivée} : $\left(\dfrac{1}{x+1}\right)' = -\dfrac{1}{(x+1)^2}$ ; le signe de $f'$ donne les variations de $f$.
\item \cle{Centre de symétrie} : $\Omega(\alpha\,;\,\beta)$ est centre de symétrie de $\mathcal{C}_f$ si $f(\alpha + h) + f(\alpha - h) = 2\beta$ pour tout $h$ admissible.
\item \cle{Résolution graphique} : $f(x) \leq k$ se lit en repérant les abscisses des points de $\mathcal{C}_f$ situés sous la droite $y = k$.
\end{itemize}
\end{aretenir}

\courssub{Démarche et corrigé commenté}

\begin{exoresolu}[Corrigé complet de l'activité]
On considère $f(x) = \dfrac{2x^2+3x+9}{x+1}$ pour $x > 0$, coût moyen en milliers de FCFA par tonne.
\tcblower
\textbf{Partie A. 1.} $f$ est définie si et seulement si $x + 1 \neq 0$, donc $D_f = \mathbb{R} \setminus \{-1\}$. Dans le contexte, $x$ désigne une masse de cacao en tonnes : seules les valeurs $x > 0$ ont un sens économique.

\textbf{2.} Quand $x \to -1^+$, le numérateur tend vers $2 - 3 + 9 = 8 > 0$ et le dénominateur vers $0^+$ : $\lim\limits_{x \to -1^+} f(x) = +\infty$. De même $\lim\limits_{x \to -1^-} f(x) = -\infty$. La droite d'équation $x = -1$ est donc \cle{asymptote verticale}. Elle est hors du domaine économique : aucune interprétation concrète.

\textbf{3.} Division euclidienne : $2x^2 + 3x + 9 = (x+1)(2x+1) + 8$. En effet $(x+1)(2x+1) = 2x^2 + 3x + 1$, et $9 - 1 = 8$. D'où :
\[ f(x) = 2x + 1 + \frac{8}{x+1}. \]

\textbf{4.} On a $f(x) - (2x+1) = \dfrac{8}{x+1}$, et $\lim\limits_{x \to +\infty} \dfrac{8}{x+1} = 0$. La droite $\Delta : y = 2x+1$ est donc \cle{asymptote oblique} à $\mathcal{C}_f$ en $+\infty$. Comme $\dfrac{8}{x+1} > 0$ pour $x > 0$, la courbe est \textbf{au-dessus} de $\Delta$ sur $]0\,;\,+\infty[$.

\textbf{Partie B. 5.} $f$ est dérivable sur $]0\,;\,+\infty[$ et
\[ f'(x) = 2 - \frac{8}{(x+1)^2}. \]

\textbf{6.} $f'(x) = 0 \iff (x+1)^2 = 4 \iff x + 1 = 2$ (car $x+1 > 0$) $\iff x = 1$. De plus $f'(x) < 0$ pour $0 < x < 1$ et $f'(x) > 0$ pour $x > 1$. On calcule : $f(0) = 9$, $f(1) = \dfrac{2+3+9}{2} = 7$, et $\lim\limits_{x \to +\infty} f(x) = +\infty$.

\begin{center}
\begin{tabular}{|c|ccccc|}\hline
$x$ & $0$ & & $1$ & & $+\infty$ \\ \hline
$f'(x)$ & & $-$ & $0$ & $+$ & \\ \hline
$f$ & $9$ & $\searrow$ & $7$ & $\nearrow$ & $+\infty$ \\ \hline
\end{tabular}
\end{center}

\textbf{7.} Pour tout $h \neq 0$ :
\begin{align*}
f(-1+h) + f(-1-h) &= \left[2(-1+h)+1+\frac{8}{h}\right] + \left[2(-1-h)+1-\frac{8}{h}\right] \\
&= -2 + 2h + 1 - 2 - 2h + 1 = -2 = 2 \times (-1).
\end{align*}
Le point $\Omega(-1\,;\,-1)$ est donc \cle{centre de symétrie} de la courbe complète de $f$. (Remarque : $\Omega$ est l'intersection des deux asymptotes $x=-1$ et $y=2x+1$.)

\textbf{Partie C. 8.} Tableau de valeurs pour le tracé :
\begin{center}
\begin{tabular}{|c|cccccc|}\hline
$x$ & $0$ & $0{,}5$ & $1$ & $2$ & $5$ & $10$ \\ \hline
$f(x)$ & $9$ & $7{,}67$ & $7$ & $7{,}67$ & $11{,}33$ & $21{,}73$ \\ \hline
\end{tabular}
\end{center}
On place les points, on trace $\Delta : y = 2x+1$ (passant par $(0\,;\,1)$ et $(4\,;\,9)$), puis la courbe qui descend de $(0\,;\,9)$ vers le minimum $(1\,;\,7)$ et remonte en se rapprochant de $\Delta$.

\textbf{9.} Graphiquement, le point le plus bas de la courbe a pour abscisse $x \approx 1$. Le calcul confirme : $f$ admet un minimum en $x = 1$, et $f(1) = 7$. \textbf{La coopérative doit produire 1 tonne par lot pour un coût moyen minimal de \num{7000} FCFA par tonne.}

\textbf{10.} Graphiquement, la droite $y = 8$ coupe $\mathcal{C}_f$ en deux points d'abscisses $x_1 \approx 0{,}2$ et $x_2 \approx 2{,}3$ ; la courbe est sous la droite entre ces deux valeurs. Algébriquement, pour $x > 0$ :
\begin{align*}
f(x) \leq 8 &\iff 2x^2 + 3x + 9 \leq 8(x+1) \quad (\text{car } x+1 > 0) \\
&\iff 2x^2 - 5x + 1 \leq 0.
\end{align*}
Le discriminant vaut $\Delta = 25 - 8 = 17$ ; les racines sont $x_1 = \dfrac{5 - \sqrt{17}}{4} \approx 0{,}22$ et $x_2 = \dfrac{5 + \sqrt{17}}{4} \approx 2{,}28$. Le trinôme est négatif entre ses racines :
\[ S = \left[\frac{5-\sqrt{17}}{4}\,;\,\frac{5+\sqrt{17}}{4}\right] \approx [0{,}22\,;\,2{,}28]. \]
\textbf{La contrainte de l'exportateur est respectée pour une production comprise entre environ 220 kg et 2,28 tonnes.}

\textbf{11.} Pour les très grandes productions, $f(x) \approx 2x+1$ : le coût moyen \textbf{croît indéfiniment}, d'environ 2 milliers de FCFA par tonne supplémentaire. Recommandation : \og Monsieur Essomba, le modèle montre qu'au-delà de 2,3 tonnes par lot, le coût moyen dépasse le seuil de l'exportateur et croît sans limite. Il faut fractionner la collecte en lots voisins de 1 tonne et investir dans un second séchoir avant d'accepter davantage de récolte. \fg
\end{exoresolu}

\begin{attention}[Pièges rencontrés pendant l'activité]
\begin{itemize}
\item Ne pas oublier que multiplier une inéquation par $x+1$ exige de vérifier son signe : ici $x > 0$ donc $x + 1 > 0$, le sens de l'inégalité est conservé.
\item Le minimum de $f$ n'est \textbf{pas} en $x = 0$ : le bord du domaine n'est pas forcément l'extremum ; seule l'étude du signe de $f'$ décide.
\item L'asymptote verticale $x = -1$ n'a aucun sens économique : distinguer toujours le domaine mathématique du domaine du modèle.
\item Dans la résolution graphique, donner les solutions en \emph{abscisses} (tonnes), pas en ordonnées.
\end{itemize}
\end{attention}

\courssub{Bilan de l'activité}

Cette activité a permis de mobiliser, dans un contexte authentique de la filière cacao camerounaise, l'intégralité de la démarche d'étude d'une fonction rationnelle du type $x \mapsto \dfrac{ax^2+bx+c}{dx+e}$ : recherche du domaine, mise en évidence d'une asymptote verticale, division euclidienne révélant une asymptote oblique, étude des variations par la dérivée, et mise en lumière d'un centre de symétrie situé à l'intersection des asymptotes. Elle a surtout montré la \emph{finalité} de ces outils : le tableau de variations livre la production optimale (1 tonne, coût minimal de \num{7000} FCFA), la courbe permet de résoudre graphiquement l'inéquation budgétaire avant de la confirmer algébriquement, et l'asymptote oblique traduit mathématiquement les \og maladresses de la croissance \fg{} : produire trop, c'est produire cher. Tu as ainsi constaté qu'une courbe bien construite est un véritable instrument de décision économique — compétence directement exigible au Probatoire, où l'exploitation graphique d'une étude de fonction est un classique des problèmes.

\newpage

\coursec{Méthodes et exercices résolus}

\courssub{Méthodes et exercices résolus}

\begin{methode}[Déterminer les asymptotes parallèles aux axes]
Pour déterminer les asymptotes parallèles aux axes à la courbe $\mathcal{C}_f$ d'une fonction $f$ :
\begin{enumerate}
\item \textbf{Asymptote verticale} : repérer les valeurs interdites $a$ de $f$ (valeurs qui annulent le dénominateur). Calculer $\lim\limits_{x\to a^-}f(x)$ et $\lim\limits_{x\to a^+}f(x)$. Si l'une de ces limites est infinie ($+\infty$ ou $-\infty$), alors la droite d'équation $\boxed{x=a}$ est asymptote verticale à $\mathcal{C}_f$.
\item \textbf{Asymptote horizontale} : calculer $\lim\limits_{x\to+\infty}f(x)$ et $\lim\limits_{x\to-\infty}f(x)$. Si l'une de ces limites est un réel fini $b$, alors la droite d'équation $\boxed{y=b}$ est asymptote horizontale à $\mathcal{C}_f$ (en $+\infty$ ou en $-\infty$ selon le cas).
\item \textbf{Rédaction} : annoncer toujours la limite calculée AVANT de conclure sur l'asymptote, en précisant l'équation de la droite.
\end{enumerate}
\textbf{Réflexes} : pour une fonction rationnelle, la limite en l'infini est celle du quotient des termes de plus haut degré ; en une valeur interdite $a$, on étudie le signe du dénominateur au voisinage de $a$ (forme $\dfrac{\text{nombre non nul}}{0^+}$ ou $\dfrac{\text{nombre non nul}}{0^-}$).
\end{methode}

\begin{exoresolu}[Asymptotes parallèles aux axes]
Soit la fonction $f$ définie par $f(x)=\dfrac{2x-3}{x+1}$ et $\mathcal{C}_f$ sa courbe représentative.
\begin{enumerate}
\item Déterminer l'ensemble de définition de $f$.
\item Déterminer par leurs équations les asymptotes parallèles aux axes à la courbe $\mathcal{C}_f$.
\end{enumerate}
\tcblower
\textbf{1. Ensemble de définition.}

$f(x)$ existe si et seulement si $x+1\neq 0$, c'est-à-dire $x\neq -1$. Donc :
\[D_f=\mathbb{R}\setminus\{-1\}=]-\infty;-1[\cup]-1;+\infty[.\]

\textbf{2. Asymptotes.}

\emph{Asymptote verticale.} La valeur interdite est $-1$. Étudions la limite de $f$ en $-1$. Le numérateur vaut $2\times(-1)-3=-5\neq 0$. Le dénominateur $x+1$ tend vers $0$ :
\begin{itemize}
\item si $x\to -1^-$ (c'est-à-dire $x<-1$), alors $x+1<0$, donc $x+1\to 0^-$, d'où $\lim\limits_{x\to -1^-}f(x)=\dfrac{-5}{0^-}=+\infty$ ;
\item si $x\to -1^+$, alors $x+1\to 0^+$, d'où $\lim\limits_{x\to -1^+}f(x)=\dfrac{-5}{0^+}=-\infty$.
\end{itemize}
Comme ces limites sont infinies, la droite d'équation $\boxed{x=-1}$ est asymptote verticale à $\mathcal{C}_f$.

\emph{Asymptote horizontale.} En l'infini, une fonction rationnelle a même limite que le quotient de ses termes de plus haut degré :
\[\lim_{x\to\pm\infty}f(x)=\lim_{x\to\pm\infty}\dfrac{2x}{x}=2.\]
Cette limite est finie, donc la droite d'équation $\boxed{y=2}$ est asymptote horizontale à $\mathcal{C}_f$ en $-\infty$ et en $+\infty$.

\textbf{Conclusion} : $\mathcal{C}_f$ admet l'asymptote verticale d'équation $x=-1$ et l'asymptote horizontale d'équation $y=2$.
\end{exoresolu}

\begin{methode}[Montrer qu'une droite donnée est asymptote oblique]
Pour montrer que la droite $(D)$ d'équation $y=ax+b$ ($a\neq 0$) est asymptote oblique à $\mathcal{C}_f$ en $+\infty$ (ou en $-\infty$) :
\begin{enumerate}
\item Former la différence $f(x)-(ax+b)$ et la réduire au même dénominateur.
\item Calculer $\lim\limits_{x\to+\infty}\big[f(x)-(ax+b)\big]$.
\item Si cette limite vaut $\boxed{0}$, conclure : « la droite $(D)$ d'équation $y=ax+b$ est asymptote oblique à $\mathcal{C}_f$ en $+\infty$ ».
\end{enumerate}
\textbf{Pour trouver $a$ et $b$} quand ils ne sont pas donnés : effectuer la division euclidienne du numérateur par le dénominateur (ou procéder par identification). Pour $f(x)=\dfrac{P(x)}{dx+e}$ avec $P$ de degré $2$, on obtient $f(x)=ax+b+\dfrac{r}{dx+e}$ avec $r$ constant ; alors la droite d'équation $y=ax+b$ est l'asymptote oblique.
\end{methode}

\begin{exoresolu}[Asymptote oblique]
Soit $g$ la fonction définie par $g(x)=\dfrac{x^2-3x+5}{x-1}$ et $\mathcal{C}_g$ sa courbe.
\begin{enumerate}
\item Déterminer trois réels $a$, $b$ et $c$ tels que, pour tout $x\neq 1$ : $g(x)=ax+b+\dfrac{c}{x-1}$.
\item Montrer que la droite $(\Delta)$ d'équation $y=x-2$ est asymptote oblique à $\mathcal{C}_g$.
\end{enumerate}
\tcblower
\textbf{1. Décomposition.} Pour tout $x\neq 1$ :
\[ax+b+\dfrac{c}{x-1}=\dfrac{(ax+b)(x-1)+c}{x-1}=\dfrac{ax^2+(b-a)x+(c-b)}{x-1}.\]
Par identification des coefficients avec $g(x)=\dfrac{x^2-3x+5}{x-1}$, on obtient le système :
\[\begin{cases} a=1 \\ b-a=-3 \\ c-b=5 \end{cases} \iff \begin{cases} a=1 \\ b=-2 \\ c=3 \end{cases}\]
Donc, pour tout $x\neq 1$ : $\boxed{g(x)=x-2+\dfrac{3}{x-1}}$.

\textbf{2. Asymptote oblique.} Formons la différence :
\[g(x)-(x-2)=\dfrac{3}{x-1}.\]
Or $\lim\limits_{x\to+\infty}(x-1)=+\infty$, donc :
\[\lim_{x\to+\infty}\big[g(x)-(x-2)\big]=\lim_{x\to+\infty}\dfrac{3}{x-1}=0.\]
De même $\lim\limits_{x\to-\infty}\big[g(x)-(x-2)\big]=0$.

Comme cette limite est nulle, la droite $(\Delta)$ d'équation $y=x-2$ est asymptote oblique à $\mathcal{C}_g$ en $-\infty$ et en $+\infty$.
\end{exoresolu}

\begin{methode}[Étudier et représenter une fonction homographique]
Soit $f(x)=\dfrac{ax+b}{cx+d}$ avec $c\neq 0$ et $ad-bc\neq 0$. Plan d'étude :
\begin{enumerate}
\item \textbf{Ensemble de définition} : $D_f=\mathbb{R}\setminus\left\{-\dfrac{d}{c}\right\}$.
\item \textbf{Limites et asymptotes} : asymptote verticale d'équation $x=-\dfrac{d}{c}$ (limites infinies) ; asymptote horizontale d'équation $y=\dfrac{a}{c}$ (limite en l'infini).
\item \textbf{Centre de symétrie} : le point $\Omega\left(-\dfrac{d}{c}\,;\,\dfrac{a}{c}\right)$, intersection des deux asymptotes, est centre de symétrie de la courbe. Pour le prouver : vérifier que $f(2x_\Omega-x)+f(x)=2y_\Omega$ pour tout $x$ convenable.
\item \textbf{Dérivée et variations} : $f'(x)=\dfrac{ad-bc}{(cx+d)^2}$. Le signe de $f'$ est celui de $ad-bc$ (le dénominateur est un carré, strictement positif sur $D_f$) : $f$ est strictement monotone sur chaque intervalle de $D_f$.
\item \textbf{Tableau de variations} : la valeur interdite y figure dans une colonne dédoublée ; on ne relie jamais les deux branches par une flèche continue à travers elle.
\item \textbf{Tracé} : placer les asymptotes en pointillés, le centre $\Omega$, quelques points (intersections avec les axes), puis tracer les deux branches de l'hyperbole.
\end{enumerate}
\end{methode}

\begin{exoresolu}[Étude complète d'une fonction homographique]
Soit $f$ la fonction définie par $f(x)=\dfrac{x+3}{x-2}$ et $(\mathcal{C})$ sa courbe dans un repère orthonormé $(O;\vec{i},\vec{j})$.
\begin{enumerate}
\item Déterminer $D_f$ et les limites de $f$ aux bornes de $D_f$.
\item En déduire les asymptotes à $(\mathcal{C})$.
\item Étudier les variations de $f$ et dresser son tableau de variations.
\item Montrer que le point $\Omega(2;1)$ est centre de symétrie de $(\mathcal{C})$.
\item Tracer $(\mathcal{C})$.
\end{enumerate}
\tcblower
\textbf{1. Ensemble de définition et limites.} $x-2=0\iff x=2$, donc $D_f=\mathbb{R}\setminus\{2\}$.

En $2$ : le numérateur vaut $2+3=5>0$.
\begin{itemize}
\item Si $x\to 2^-$, $x-2\to 0^-$, donc $\lim\limits_{x\to 2^-}f(x)=\dfrac{5}{0^-}=-\infty$.
\item Si $x\to 2^+$, $x-2\to 0^+$, donc $\lim\limits_{x\to 2^+}f(x)=\dfrac{5}{0^+}=+\infty$.
\end{itemize}
En l'infini : $\lim\limits_{x\to\pm\infty}f(x)=\lim\limits_{x\to\pm\infty}\dfrac{x}{x}=1$.

\textbf{2. Asymptotes.} Des limites précédentes, on déduit : la droite d'équation $x=2$ est asymptote verticale à $(\mathcal{C})$, et la droite d'équation $y=1$ est asymptote horizontale à $(\mathcal{C})$ en $-\infty$ et en $+\infty$.

\textbf{3. Variations.} $f$ est dérivable sur $D_f$ et, avec $a=1$, $b=3$, $c=1$, $d=-2$ :
\[f'(x)=\dfrac{ad-bc}{(cx+d)^2}=\dfrac{1\times(-2)-3\times 1}{(x-2)^2}=\dfrac{-5}{(x-2)^2}.\]
Pour tout $x\in D_f$, $(x-2)^2>0$, donc $f'(x)<0$ : $f$ est strictement décroissante sur $]-\infty;2[$ et sur $]2;+\infty[$.

\begin{center}
\begin{tabular}{|c|ccc||cccc|}\hline
$x$ & $-\infty$ & & & $2$ & & & $+\infty$ \\ \hline
$f'(x)$ & & $-$ & & & $-$ & & \\ \hline
$f$ & $1$ & $\searrow$ & $-\infty$ & $+\infty$ & $\searrow$ & & $1$ \\ \hline
\end{tabular}
\end{center}

\textbf{4. Centre de symétrie.} Pour tout $x\neq 2$, $2x_\Omega-x=4-x$ appartient aussi à $D_f$ (car $4-x\neq 2\iff x\neq 2$). Calculons :
\begin{align*}
f(4-x)+f(x)&=\dfrac{(4-x)+3}{(4-x)-2}+\dfrac{x+3}{x-2}=\dfrac{7-x}{2-x}+\dfrac{x+3}{x-2}\\
&=\dfrac{-(7-x)+(x+3)}{x-2}=\dfrac{2x-4}{x-2}=\dfrac{2(x-2)}{x-2}=2=2\times 1.
\end{align*}
Donc $\dfrac{f(4-x)+f(x)}{2}=1=y_\Omega$ : le point $\Omega(2;1)$ est bien centre de symétrie de $(\mathcal{C})$.

\textbf{5. Tracé.} Points utiles : $f(0)=-\dfrac{3}{2}$ (intersection avec l'axe des ordonnées) ; $f(x)=0\iff x=-3$ (intersection avec l'axe des abscisses) ; $f(3)=6$ ; $f(1)=-4$.

\begin{popfigure}
\begin{center}
\begin{tikzpicture}[scale=0.75]
\draw[popline] (-5.5,0) -- (7.5,0) node[below left]{\small $x$};
\draw[popline] (0,-5) -- (0,7) node[below left]{\small $y$};
\draw[dashed,poporange,thick] (2,-5) -- (2,7) node[above left]{\small $x=2$};
\draw[dashed,popblue,thick] (-5.5,1) -- (7.5,1) node[above left]{\small $y=1$};
\draw[popink,very thick,domain=-5:1.2,samples=100] plot (\x,{(\x+3)/(\x-2)});
\draw[popink,very thick,domain=2.8:7.2,samples=100] plot (\x,{(\x+3)/(\x-2)});
\fill[popdark] (2,1) circle (2pt) node[below right]{\small $\Omega$};
\fill[popdark] (0,-1.5) circle (1.5pt) node[below left]{\small $-\frac{3}{2}$};
\fill[popdark] (-3,0) circle (1.5pt) node[above left]{\small $-3$};
\foreach \x in {-4,-2,2,4,6} \draw (\x,0.08) -- (\x,-0.08) node[below]{\scriptsize $\x$};
\foreach \y in {-4,-2,2,4,6} \draw (0.08,\y) -- (-0.08,\y) node[left]{\scriptsize $\y$};
\end{tikzpicture}
\end{center}
\legende{Courbe de $f(x)=\dfrac{x+3}{x-2}$ : hyperbole de centre $\Omega(2;1)$.}
\end{popfigure}
\end{exoresolu}

\begin{methode}[Étudier et représenter une fonction polynôme de degré $\leq 3$]
Soit $f(x)=ax^3+bx^2+cx+d$ ($a\neq 0$ pour un vrai degré 3). Plan d'étude :
\begin{enumerate}
\item \textbf{Ensemble de définition} : $D_f=\mathbb{R}$ (un polynôme est défini partout).
\item \textbf{Limites aux bornes} : en $\pm\infty$, $f(x)$ a même limite que son terme de plus haut degré $ax^3$.
\item \textbf{Dérivée} : $f'(x)=3ax^2+2bx+c$ (trinôme du second degré).
\item \textbf{Signe de la dérivée} : calculer le discriminant $\Delta'$ de $f'$.
\begin{itemize}
\item Si $\Delta'\leq 0$ : $f'$ garde un signe constant (celui de $3a$), $f$ est monotone sur $\mathbb{R}$.
\item Si $\Delta'>0$ : deux racines $x_1<x_2$ ; $f'$ s'annule en $x_1$ et $x_2$ et est du signe de $3a$ à l'extérieur des racines.
\end{itemize}
\item \textbf{Tableau de variations} complet : limites aux bornes ET valeurs des extrema $f(x_1)$, $f(x_2)$.
\item \textbf{Tracé} : tableau de valeurs pour quelques points bien choisis, tangentes horizontales aux extrema, et centre de symétrie de la cubique en le point d'abscisse $-\dfrac{b}{3a}$ (point d'inflexion).
\end{enumerate}
\end{methode}

\begin{exoresolu}[Étude d'une fonction polynôme de degré 3]
Soit $f$ la fonction définie sur $\mathbb{R}$ par $f(x)=x^3-3x+2$ et $(\mathcal{C})$ sa courbe représentative.
\begin{enumerate}
\item Calculer les limites de $f$ en $-\infty$ et en $+\infty$.
\item Étudier les variations de $f$ et dresser son tableau de variations.
\item Montrer que le point $I(0;2)$ est centre de symétrie de $(\mathcal{C})$.
\item Tracer $(\mathcal{C})$.
\end{enumerate}
\tcblower
\textbf{1. Limites.} En l'infini, $f(x)$ se comporte comme son terme de plus haut degré $x^3$ :
\[\lim_{x\to-\infty}f(x)=\lim_{x\to-\infty}x^3=-\infty \quad ; \quad \lim_{x\to+\infty}f(x)=\lim_{x\to+\infty}x^3=+\infty.\]

\textbf{2. Variations.} $f$ est dérivable sur $\mathbb{R}$ et :
\[f'(x)=3x^2-3=3(x^2-1)=3(x-1)(x+1).\]
$f'(x)$ est un trinôme de racines $-1$ et $1$, du signe de $a=3>0$ à l'extérieur des racines : $f'(x)>0$ sur $]-\infty;-1[\cup]1;+\infty[$ et $f'(x)<0$ sur $]-1;1[$.

Valeurs remarquables : $f(-1)=-1+3+2=4$ ; $f(1)=1-3+2=0$.

\begin{center}
\begin{tabular}{|c|ccccccc|}\hline
$x$ & $-\infty$ & & $-1$ & & $1$ & & $+\infty$ \\ \hline
$f'(x)$ & & $+$ & $0$ & $-$ & $0$ & $+$ & \\ \hline
$f$ & $-\infty$ & $\nearrow$ & $4$ & $\searrow$ & $0$ & $\nearrow$ & $+\infty$ \\ \hline
\end{tabular}
\end{center}

\textbf{3. Centre de symétrie.} Pour tout réel $x$ :
\[f(-x)+f(x)=\big[(-x)^3-3(-x)+2\big]+\big[x^3-3x+2\big]=-x^3+3x+2+x^3-3x+2=4=2\times 2.\]
Donc $\dfrac{f(0-x)+f(0+x)}{2}=2$ pour tout réel $x$ : le point $I(0;2)$ est centre de symétrie de $(\mathcal{C})$. (On retrouve bien l'abscisse $-\dfrac{b}{3a}=0$.)

\textbf{4. Tracé.} Points : $f(-2)=0$, $f(-1)=4$, $f(0)=2$, $f(1)=0$, $f(2)=4$. Tangentes horizontales en $A(-1;4)$ et $B(1;0)$.

\begin{popfigure}
\begin{center}
\begin{tikzpicture}[scale=0.8]
\draw[popline] (-3.2,0) -- (3.2,0) node[below left]{\small $x$};
\draw[popline] (0,-2.5) -- (0,5) node[below left]{\small $y$};
\draw[popink,very thick,domain=-2.35:2.35,samples=200] plot (\x,{\x*\x*\x-3*\x+2});
\fill[popdark] (0,2) circle (2pt) node[above right]{\small $I$};
\fill[poporange] (-1,4) circle (2pt) node[above]{\small $A(-1;4)$};
\fill[poporange] (1,0) circle (2pt) node[below right]{\small $B(1;0)$};
\fill[popdark] (-2,0) circle (1.5pt) node[below left]{\small $-2$};
\fill[popdark] (2,4) circle (1.5pt);
\foreach \x in {-2,-1,1,2} \draw (\x,0.08) -- (\x,-0.08) node[below]{\scriptsize $\x$};
\foreach \y in {2,4} \draw (0.08,\y) -- (-0.08,\y) node[left]{\scriptsize $\y$};
\end{tikzpicture}
\end{center}
\legende{Courbe de $f(x)=x^3-3x+2$, de centre de symétrie $I(0;2)$.}
\end{popfigure}
\end{exoresolu}

\begin{methode}[Étudier une fonction $x\mapsto\dfrac{ax^2+bx+c}{dx+e}$ ($a\neq 0$, $d\neq 0$)]
\begin{enumerate}
\item \textbf{Ensemble de définition} : $D_f=\mathbb{R}\setminus\left\{-\dfrac{e}{d}\right\}$.
\item \textbf{Limites} : en la valeur interdite, limites infinies (asymptote verticale d'équation $x=-\dfrac{e}{d}$) ; en l'infini, $f(x)$ se comporte comme $\dfrac{ax^2}{dx}=\dfrac{a}{d}x$, donc les limites sont infinies : il n'y a pas d'asymptote horizontale.
\item \textbf{Asymptote oblique} : par division euclidienne (ou identification), écrire $f(x)=\alpha x+\beta+\dfrac{\gamma}{dx+e}$ ; la droite d'équation $y=\alpha x+\beta$ est asymptote oblique car $\lim\limits_{x\to\pm\infty}\dfrac{\gamma}{dx+e}=0$.
\item \textbf{Dérivée} : $f'(x)=\dfrac{N(x)}{(dx+e)^2}$ où $N(x)=(2ax+b)(dx+e)-d(ax^2+bx+c)$ est un trinôme ; étudier son signe via son discriminant (si une racine de $N$ était égale à la valeur interdite, il faudrait l'exclure de l'étude).
\item \textbf{Tableau de variations} avec la valeur interdite en colonne dédoublée, puis tracé avec les asymptotes en pointillés.
\end{enumerate}
\end{methode}

\begin{exoresolu}[Étude de $g(x)=\dfrac{x^2-3x+5}{x-1}$]
On reprend la fonction $g(x)=\dfrac{x^2-3x+5}{x-1}$ de l'exercice sur l'asymptote oblique, de courbe $(\mathcal{C})$.
\begin{enumerate}
\item Déterminer $D_g$ et les asymptotes à $(\mathcal{C})$.
\item Étudier les variations de $g$ et dresser le tableau de variations.
\end{enumerate}
\tcblower
\textbf{1. Ensemble de définition et asymptotes.} $x-1=0\iff x=1$, donc $D_g=\mathbb{R}\setminus\{1\}$.

En $1$ : le numérateur vaut $1-3+5=3>0$. Si $x\to 1^-$, $x-1\to 0^-$ donc $\lim\limits_{x\to 1^-}g(x)=-\infty$ ; si $x\to 1^+$, $x-1\to 0^+$ donc $\lim\limits_{x\to 1^+}g(x)=+\infty$. La droite d'équation $x=1$ est asymptote verticale à $(\mathcal{C})$.

On a montré que $g(x)=x-2+\dfrac{3}{x-1}$ avec $\lim\limits_{x\to\pm\infty}\big[g(x)-(x-2)\big]=0$ : la droite d'équation $y=x-2$ est asymptote oblique à $(\mathcal{C})$ en $-\infty$ et en $+\infty$.

\textbf{2. Variations.} $g$ est dérivable sur $D_g$. Avec $u(x)=x^2-3x+5$ et $v(x)=x-1$ :
\begin{align*}
g'(x)&=\dfrac{(2x-3)(x-1)-(x^2-3x+5)\times 1}{(x-1)^2}=\dfrac{2x^2-5x+3-x^2+3x-5}{(x-1)^2}\\
&=\dfrac{x^2-2x-2}{(x-1)^2}.
\end{align*}
Le dénominateur est strictement positif sur $D_g$, donc $g'(x)$ a le signe du numérateur $x^2-2x-2$. Discriminant : $\Delta=(-2)^2-4\times 1\times(-2)=4+8=12>0$ ; racines :
\[x_1=\dfrac{2-\sqrt{12}}{2}=1-\sqrt{3}\approx -0{,}73 \quad ; \quad x_2=1+\sqrt{3}\approx 2{,}73.\]
Le trinôme est positif à l'extérieur des racines, négatif entre elles (aucune racine n'est égale à la valeur interdite $1$).

Extrema, en utilisant la forme $g(x)=x-2+\dfrac{3}{x-1}$ :
\[g(1-\sqrt{3})=(1-\sqrt{3})-2+\dfrac{3}{-\sqrt{3}}=-1-\sqrt{3}-\sqrt{3}=-1-2\sqrt{3}\approx -4{,}46 ;\]
\[g(1+\sqrt{3})=(1+\sqrt{3})-2+\dfrac{3}{\sqrt{3}}=-1+\sqrt{3}+\sqrt{3}=-1+2\sqrt{3}\approx 2{,}46.\]

\begin{center}
\begin{tabular}{|c|ccc||cccccc|}\hline
$x$ & $-\infty$ & & $1-\sqrt{3}$ & & $1$ & & $1+\sqrt{3}$ & & $+\infty$ \\ \hline
$g'(x)$ & & $+$ & $0$ & $-$ & & $-$ & $0$ & $+$ & \\ \hline
$g$ & $-\infty$ & $\nearrow$ & $-1-2\sqrt{3}$ & $\searrow$ & $-\infty\quad+\infty$ & $\searrow$ & $-1+2\sqrt{3}$ & $\nearrow$ & $+\infty$ \\ \hline
\end{tabular}
\end{center}

\textbf{Conclusion} : $g$ est croissante sur $]-\infty;1-\sqrt{3}]$ et sur $[1+\sqrt{3};+\infty[$, décroissante sur $[1-\sqrt{3};1[$ et sur $]1;1+\sqrt{3}]$.
\end{exoresolu}

\begin{methode}[Résoudre graphiquement une équation ou une inéquation]
À partir de la courbe $\mathcal{C}_f$ tracée dans un repère :
\begin{enumerate}
\item \textbf{Équation $f(x)=k$} : tracer la droite d'équation $y=k$ ; les solutions sont les \emph{abscisses} des points d'intersection de cette droite avec $\mathcal{C}_f$.
\item \textbf{Équation $f(x)=g(x)$} : les solutions sont les abscisses des points d'intersection de $\mathcal{C}_f$ et $\mathcal{C}_g$.
\item \textbf{Inéquation $f(x)\geq k$} : les solutions sont les abscisses des points de $\mathcal{C}_f$ situés \emph{au-dessus} de la droite d'équation $y=k$ (on lit des intervalles sur l'axe des abscisses).
\item \textbf{Inéquation $f(x)\geq g(x)$} : abscisses des points où $\mathcal{C}_f$ est au-dessus de $\mathcal{C}_g$.
\item \textbf{Rédaction} : donner l'ensemble solution sous forme d'intervalle ou de réunion d'intervalles, avec des crochets fermés si l'inégalité est large, ouverts sinon ; exclure toujours les valeurs interdites.
\end{enumerate}
\end{methode}

\begin{exoresolu}[Exploitation graphique d'une courbe]
On reprend la courbe $(\mathcal{C})$ de la fonction $f(x)=\dfrac{x+3}{x-2}$ étudiée plus haut, et on considère la droite $(D)$ d'équation $y=-x+1$.
\begin{enumerate}
\item Résoudre graphiquement l'équation $f(x)=0$.
\item Résoudre graphiquement l'inéquation $f(x)\geq 1$.
\item Déterminer graphiquement le nombre de solutions de l'équation $f(x)=-x+1$, puis confirmer par le calcul.
\end{enumerate}
\tcblower
\textbf{1. Équation $f(x)=0$.} Les solutions sont les abscisses des points où $(\mathcal{C})$ coupe l'axe des abscisses. Graphiquement, $(\mathcal{C})$ coupe l'axe $(Ox)$ en un seul point, d'abscisse $-3$.
\[\boxed{S=\{-3\}}.\]
(Vérification algébrique : $f(x)=0\iff x+3=0\iff x=-3$.)

\textbf{2. Inéquation $f(x)\geq 1$.} On cherche les abscisses des points de $(\mathcal{C})$ situés au-dessus de la droite d'équation $y=1$ (l'asymptote horizontale) ou sur celle-ci. La courbe est strictement au-dessus de cette droite uniquement sur la branche de droite, pour $x>2$ ; la valeur $x=2$ est interdite, et l'égalité $f(x)=1$ est impossible car $\dfrac{x+3}{x-2}=1\iff x+3=x-2$, ce qui est absurde.
\[\boxed{S=]2;+\infty[}.\]

\textbf{3. Équation $f(x)=-x+1$.} Les solutions éventuelles sont les abscisses des points d'intersection de $(\mathcal{C})$ et de $(D)$. La droite $(D)$ passe par les points $(0;1)$ et $(1;0)$ ; en la traçant sur le graphique, on constate qu'elle \textbf{ne rencontre pas} la courbe $(\mathcal{C})$ : l'équation semble n'avoir aucune solution.

Confirmation par le calcul : pour $x\neq 2$,
\[\dfrac{x+3}{x-2}=-x+1\iff x+3=(-x+1)(x-2)\iff x+3=-x^2+3x-2\iff x^2-2x+5=0.\]
Le discriminant vaut $\Delta=(-2)^2-4\times 1\times 5=4-20=-16<0$ : l'équation n'a effectivement aucune solution réelle.
\[\boxed{S=\varnothing}.\]
Cet exemple illustre une règle d'or : une lecture graphique donne une réponse approchée ou un nombre de solutions, et le calcul permet de confirmer lorsque l'énoncé le demande.
\end{exoresolu}

\begin{methode}[Montrer qu'une courbe admet un élément de symétrie]
\begin{enumerate}
\item \textbf{Centre de symétrie $\Omega(a;b)$} : vérifier que pour tout $x$ tel que $a+x\in D_f$, on a aussi $a-x\in D_f$, puis montrer que
\[f(a-x)+f(a+x)=2b \quad \text{c'est-à-dire} \quad \dfrac{f(a-x)+f(a+x)}{2}=b.\]
\item \textbf{Axe de symétrie d'équation $x=a$} : vérifier que $D_f$ est symétrique par rapport à $a$, puis montrer que
\[f(a-x)=f(a+x) \quad \text{pour tout } x \text{ convenable.}\]
\item \textbf{Cas particuliers} : $f$ paire ($f(-x)=f(x)$, avec $D_f$ symétrique par rapport à $0$) $\Rightarrow$ l'axe des ordonnées est axe de symétrie ; $f$ impaire ($f(-x)=-f(x)$) $\Rightarrow$ l'origine $O$ est centre de symétrie.
\item \textbf{Candidats naturels} : pour une fonction homographique, le centre est l'intersection des deux asymptotes ; pour une cubique $ax^3+bx^2+cx+d$, le centre est le point d'abscisse $-\dfrac{b}{3a}$ (point d'inflexion).
\end{enumerate}
\end{methode}

\begin{exoresolu}[Éléments de symétrie]
\begin{enumerate}
\item Soit $f(x)=x^2-4x+5$. Montrer que la droite d'équation $x=2$ est axe de symétrie de la courbe de $f$.
\item Soit $h(x)=\dfrac{2x+1}{x-3}$. Montrer que le point $\Omega(3;2)$ est centre de symétrie de la courbe de $h$.
\end{enumerate}
\tcblower
\textbf{1. Axe de symétrie.} $D_f=\mathbb{R}$ est symétrique par rapport à $2$. Pour tout réel $x$ :
\begin{align*}
f(2-x)&=(2-x)^2-4(2-x)+5=4-4x+x^2-8+4x+5=x^2+1,\\
f(2+x)&=(2+x)^2-4(2+x)+5=4+4x+x^2-8-4x+5=x^2+1.
\end{align*}
Donc $f(2-x)=f(2+x)$ pour tout réel $x$ : la droite d'équation $\boxed{x=2}$ est axe de symétrie de la courbe de $f$. (C'est la parabole de sommet $S(2;1)$.)

\textbf{2. Centre de symétrie.} $D_h=\mathbb{R}\setminus\{3\}$ : si $3+x\in D_h$ alors $x\neq 0$, et $3-x\neq 3$, donc $3-x\in D_h$. Pour tout $x\neq 0$ :
\begin{align*}
h(3-x)+h(3+x)&=\dfrac{2(3-x)+1}{(3-x)-3}+\dfrac{2(3+x)+1}{(3+x)-3}=\dfrac{7-2x}{-x}+\dfrac{7+2x}{x}\\
&=\dfrac{-(7-2x)+(7+2x)}{x}=\dfrac{4x}{x}=4=2\times 2.
\end{align*}
Donc $\dfrac{h(3-x)+h(3+x)}{2}=2$ pour tout $x\neq 0$ : le point $\boxed{\Omega(3;2)}$ est centre de symétrie de la courbe de $h$. (On remarque que $\Omega$ est bien l'intersection des asymptotes d'équations $x=3$ et $y=2$.)
\end{exoresolu}

\begin{aretenir}[L'essentiel de la leçon]
\begin{itemize}
\item \textbf{Asymptote verticale} $x=a$ : $\lim\limits_{x\to a}f(x)=\pm\infty$ (souvent en une valeur interdite). \textbf{Asymptote horizontale} $y=b$ : $\lim\limits_{x\to\pm\infty}f(x)=b$ (réel fini). \textbf{Asymptote oblique} $y=ax+b$ : $\lim\limits_{x\to\pm\infty}[f(x)-(ax+b)]=0$.
\item \textbf{Homographique} $f(x)=\dfrac{ax+b}{cx+d}$ : hyperbole de centre $\Omega\left(-\dfrac{d}{c};\dfrac{a}{c}\right)$, dérivée $f'(x)=\dfrac{ad-bc}{(cx+d)^2}$ de signe constant.
\item \textbf{Polynôme de degré 3} : dérivée trinôme ; le signe de $f'$ se lit avec le discriminant ; la courbe admet un centre de symétrie d'abscisse $-\dfrac{b}{3a}$.
\item \textbf{Fonction $\dfrac{ax^2+bx+c}{dx+e}$} : asymptote verticale + asymptote oblique obtenue par écriture $f(x)=\alpha x+\beta+\dfrac{\gamma}{dx+e}$.
\item \textbf{Lecture graphique} : solutions d'équations = abscisses d'intersections ; solutions d'inéquations = intervalles d'abscisses.
\end{itemize}
\end{aretenir}

\begin{attention}[Les 5 erreurs qui coûtent des points au Probatoire]
\begin{enumerate}
\item \textbf{Conclure sur une asymptote sans calculer la limite.} Écrire « $x=2$ est asymptote verticale » sans avoir montré que $\lim\limits_{x\to 2}f(x)=\pm\infty$ est une réponse non justifiée : la limite calculée doit toujours précéder la conclusion.
\item \textbf{Oublier la valeur interdite dans le tableau de variations.} Pour une fonction homographique ou rationnelle, la valeur interdite doit apparaître dans le tableau (colonne dédoublée) ; on ne trace jamais une flèche « continue » à travers elle, et on ne dit jamais « $f$ est décroissante sur $\mathbb{R}$ » mais « décroissante sur chacun des intervalles $]-\infty;2[$ et $]2;+\infty[$ ».
\item \textbf{Se tromper dans la formule de la dérivée de l'homographique.} $f'(x)=\dfrac{ad-bc}{(cx+d)^2}$ : attention aux signes de $b$, $c$, $d$ dans l'identification (pour $f(x)=\dfrac{x+3}{x-2}$, on a $d=-2$ et non $2$). Une erreur de signe fausse tout le tableau de variations.
\item \textbf{Confondre les conditions d'existence d'une asymptote oblique.} Dire « $y=ax+b$ est asymptote » parce que $f(x)-(ax+b)$ « devient petit » ne suffit pas : il faut écrire explicitement $\lim\limits_{x\to\pm\infty}[f(x)-(ax+b)]=0$. De plus, une asymptote oblique n'existe que si la limite de $f$ en l'infini est infinie.
\item \textbf{Mal rédiger une résolution graphique.} Les solutions d'une équation $f(x)=k$ sont des \emph{abscisses} (des valeurs de $x$), pas des points ; une inéquation se résout par un \emph{intervalle} lu sur l'axe des abscisses, en excluant les valeurs interdites et en respectant les crochets (ouverts pour $<$, fermés pour $\leq$).
\end{enumerate}
\end{attention}

\newpage

\coursec{Exercices}

\courssub{Je vérifie mes connaissances}

\begin{exercice}[QCM : une seule bonne réponse]
Pour chaque question, choisir la bonne réponse (aucune justification n'est demandée ici).
\begin{enumerate}
\item La courbe de $f(x)=\dfrac{2x+1}{x-3}$ admet pour asymptote verticale la droite d'équation :
\begin{itemize}
\item[a)] $x=-\dfrac{1}{2}$ \qquad b) $x=3$ \qquad c) $y=2$ \qquad d) $y=3$
\end{itemize}
\item Si $\displaystyle\lim_{x\to +\infty}f(x)=-4$, alors la courbe de $f$ admet :
\begin{itemize}
\item[a)] une asymptote verticale d'équation $x=-4$
\item[b)] une asymptote horizontale d'équation $y=-4$
\item[c)] une asymptote oblique d'équation $y=-4x$
\item[d)] aucune asymptote
\end{itemize}
\item La droite d'équation $y=ax+b$ est asymptote oblique à la courbe de $f$ en $+\infty$ si :
\begin{itemize}
\item[a)] $\displaystyle\lim_{x\to +\infty}f(x)=ax+b$
\item[b)] $\displaystyle\lim_{x\to +\infty}\bigl[f(x)-(ax+b)\bigr]=0$
\item[c)] $\displaystyle\lim_{x\to +\infty}\dfrac{f(x)}{x}=0$
\item[d)] $f(x)=ax+b$ pour tout $x$ assez grand
\end{itemize}
\item Le point $\Omega(\alpha\,;\,\beta)$ est centre de symétrie de la courbe de $f$ si et seulement si, pour tout $h$ tel que $\alpha-h$ et $\alpha+h$ appartiennent au domaine :
\begin{itemize}
\item[a)] $f(\alpha+h)=f(\alpha-h)$
\item[b)] $f(\alpha+h)+f(\alpha-h)=2\beta$
\item[c)] $f(\alpha+h)-f(\alpha-h)=2\beta$
\item[d)] $f(\alpha+h)=2\beta-f(\alpha)+f(h)$
\end{itemize}
\end{enumerate}
\end{exercice}

\begin{exercice}[Vrai ou faux, en justifiant]
Répondre par \textbf{vrai} ou \textbf{faux}, en justifiant chaque réponse par une démonstration courte ou un contre-exemple.
\begin{enumerate}
\item « Si $f$ n'est pas définie en $a$, alors sa courbe admet une asymptote verticale d'équation $x=a$. »
\item « Toute fonction polynôme de degré $3$ admet un centre de symétrie. »
\item « Une courbe peut couper son asymptote horizontale. »
\item « La courbe de la fonction $x\longmapsto \dfrac{1}{x}$ admet la droite d'équation $y=x$ pour axe de symétrie. »
\end{enumerate}
\end{exercice}

\begin{exercice}[Lecture de limites et asymptotes]
On donne une fonction $f$ dont on connaît les limites suivantes :
\[ \lim_{\substack{x\to 1\\ x<1}}f(x)=-\infty,\qquad \lim_{\substack{x\to 1\\ x>1}}f(x)=+\infty,\qquad \lim_{x\to +\infty}f(x)=2,\qquad \lim_{x\to -\infty}f(x)=2. \]
\begin{enumerate}
\item Donner les équations de toutes les asymptotes à la courbe de $f$ que l'on peut déduire de ces informations.
\item La courbe de $f$ peut-elle admettre d'autres asymptotes que celles trouvées ? Justifier.
\item Tracer, sur un repère à main levée, une allure possible de la courbe de $f$ compatible avec ces données.
\end{enumerate}
\end{exercice}

\begin{exercice}[Calculs directs]
\begin{enumerate}
\item Déterminer le domaine de définition de chacune des fonctions suivantes :
\[ f(x)=\frac{3x-2}{x+1}\;;\qquad g(x)=x^3-3x^2+2\;;\qquad h(x)=\frac{x^2+2x+2}{x+1}. \]
\item Calculer les limites de $f$ en $-1$ (à gauche et à droite) et en $+\infty$.
\item Montrer que, pour tout réel $x\neq -1$ : $\;h(x)=x+1+\dfrac{1}{x+1}$.
\item En déduire $\displaystyle\lim_{x\to +\infty}\bigl[h(x)-(x+1)\bigr]$ et interpréter graphiquement le résultat.
\end{enumerate}
\end{exercice}

\courssub{Je m'entraîne}

\begin{exercice}[Asymptotes d'une fonction rationnelle]
Soit $f$ la fonction définie par $f(x)=\dfrac{2x^2-3x+1}{x-2}$ et $\mathscr{C}_f$ sa courbe représentative.
\begin{enumerate}
\item Déterminer le domaine de définition de $f$.
\item Déterminer les limites de $f$ aux bornes de son domaine et en déduire l'équation de l'asymptote verticale à $\mathscr{C}_f$.
\item Déterminer trois réels $a$, $b$ et $c$ tels que, pour tout $x\neq 2$ :
\[ f(x)=ax+b+\frac{c}{x-2}. \]
\item Montrer que la droite $\mathscr{D}$ d'équation $y=2x+1$ est asymptote oblique à $\mathscr{C}_f$ en $+\infty$ et en $-\infty$.
\item Étudier la position relative de $\mathscr{C}_f$ et de $\mathscr{D}$.
\end{enumerate}
\end{exercice}

\begin{exercice}[Axe de symétrie d'une parabole]
Soit $g$ la fonction définie sur $\mathbb{R}$ par $g(x)=x^2-4x+5$ et $\mathscr{P}$ sa courbe dans un repère orthonormé.
\begin{enumerate}
\item Démontrer que la droite d'équation $x=2$ est axe de symétrie de $\mathscr{P}$ (on pourra comparer $g(2+h)$ et $g(2-h)$).
\item Dresser le tableau de variations de $g$ et tracer $\mathscr{P}$.
\item Résoudre graphiquement l'inéquation $g(x)\leq 2$, puis retrouver le résultat par le calcul.
\end{enumerate}
\end{exercice}

\begin{exercice}[Centre de symétrie d'une cubique]
Soit $f$ la fonction définie sur $\mathbb{R}$ par $f(x)=x^3-3x^2+2$ et $\mathscr{C}$ sa courbe représentative.
\begin{enumerate}
\item Calculer $f(1+h)+f(1-h)$ pour tout réel $h$.
\item En déduire que le point $\Omega(1\,;\,0)$ est centre de symétrie de $\mathscr{C}$.
\item Calculer $f'(x)$, étudier son signe et dresser le tableau de variations de $f$.
\item Déterminer une équation de la tangente à $\mathscr{C}$ au point $\Omega$. Que constate-t-on ?
\end{enumerate}
\end{exercice}

\begin{exercice}[Exploitation graphique]
La courbe $\mathscr{C}$ ci-dessous représente une fonction $f$ définie sur $[-4\,;\,6]$. La droite d'équation $y=1$ est asymptote à $\mathscr{C}$ en $+\infty$ côté droit du graphique, et la droite d'équation $x=-1$ est asymptote verticale.
\begin{enumerate}
\item Reproduire soigneusement la courbe sur ton cahier (repère orthonormé, unité $\SI{1}{cm}$).
\item Résoudre graphiquement les équations et inéquations suivantes :
\begin{itemize}
\item[a)] $f(x)=0$ ; \qquad b) $f(x)=2$ ; \qquad c) $f(x)>1$ ; \qquad d) $f(x)\leq 0$.
\end{itemize}
\item Déterminer graphiquement le nombre de solutions de l'équation $f(x)=m$ selon les valeurs du réel $m$.
\end{enumerate}
\end{exercice}

\courssub{Je me prépare au Probatoire}

\begin{exercice}[Étude complète d'une fonction homographique]
\emph{(D'après Probatoire C -- D, session 2019.)}

On considère la fonction $f$ définie par $f(x)=\dfrac{3x-2}{x+1}$ et on note $\mathscr{C}_f$ sa courbe représentative dans le plan muni d'un repère orthonormé $(O\,;\,\vec{i},\vec{j})$ d'unité $\SI{1}{cm}$.

\textbf{Partie A : étude de la fonction}
\begin{enumerate}
\item Déterminer le domaine de définition $D_f$ de $f$.
\item Calculer les limites de $f$ aux bornes de $D_f$ et en déduire les équations des asymptotes à $\mathscr{C}_f$.
\item Montrer que, pour tout $x\in D_f$ : $\;f'(x)=\dfrac{5}{(x+1)^2}$.
\item Dresser le tableau de variations de $f$.
\item Montrer que le point $\Omega(-1\,;\,3)$ est centre de symétrie de $\mathscr{C}_f$.
\item Déterminer les coordonnées des points d'intersection de $\mathscr{C}_f$ avec les axes du repère.
\end{enumerate}

\textbf{Partie B : construction et exploitation}
\begin{enumerate}
\item Tracer les asymptotes, puis la courbe $\mathscr{C}_f$.
\item Tracer dans le même repère la droite $\mathscr{D}$ d'équation $y=x+4$.
\item Résoudre graphiquement l'équation $f(x)=x+4$, puis vérifier le résultat par le calcul.
\item Un opérateur de téléphonie mobile facture un service selon le modèle $f$ (en milliers de francs CFA) pour $x$ gigaoctets consommés ($x\geq 0$). À l'aide du graphique, déterminer à partir de quelle consommation la facture dépasse $\num{2500}$ FCFA.
\end{enumerate}
\end{exercice}

\begin{exercice}[Cubique et discussion suivant un paramètre]
\emph{(D'après Probatoire D, session 2021.)}

Soit $f$ la fonction définie sur $\mathbb{R}$ par $f(x)=x^3-3x^2+2$ et $\mathscr{C}$ sa courbe représentative dans un repère orthonormé $(O\,;\,\vec{i},\vec{j})$ d'unité $\SI{2}{cm}$.

\textbf{Partie A : étude et symétrie}
\begin{enumerate}
\item Calculer les limites de $f$ en $-\infty$ et en $+\infty$.
\item Calculer $f'(x)$, étudier son signe et dresser le tableau de variations de $f$.
\item Montrer que le point $\Omega(1\,;\,0)$ est centre de symétrie de $\mathscr{C}$.
\item Déterminer une équation de la tangente $T$ à $\mathscr{C}$ au point $\Omega$.
\item Montrer que l'équation $f(x)=0$ admet une unique solution $\alpha$ dans l'intervalle $[2\,;\,3]$, puis donner un encadrement de $\alpha$ à $10^{-1}$ près.
\end{enumerate}

\textbf{Partie B : tracé et discussion graphique}
\begin{enumerate}
\item Tracer $T$ et $\mathscr{C}$ (on placera les extrema et on utilisera la symétrie).
\item Soit $m$ un réel. Déterminer graphiquement, suivant les valeurs de $m$, le nombre de solutions de l'équation $x^3-3x^2+2=m$.
\item Une coopérative agricole de l'Ouest modélise son bénéfice mensuel (en centaines de milliers de FCFA) par $f(x)$, où $x$ est le nombre de tonnes de maïs vendues ($0\leq x\leq 3$). Déterminer graphiquement les quantités pour lesquelles le bénéfice est positif.
\end{enumerate}
\end{exercice}

\begin{exercice}[Fonction rationnelle et asymptote oblique]
\emph{(Type Probatoire C -- TI, session 2022.)}

Une entreprise de transformation de cacao à Douala estime que le coût moyen de production, en milliers de FCFA par tonne, de $x$ tonnes de cacao ($x>0$) est donné par
\[ f(x)=\frac{x^2-3x+6}{x-2}\quad \text{pour } x>2. \]
On considère la fonction $f$ définie par cette formule sur $\mathbb{R}\setminus\{2\}$ et $\mathscr{C}$ sa courbe dans un repère orthogonal $(O\,;\,\vec{i},\vec{j})$ (unités : $\SI{1}{cm}$ en abscisse, $\SI{0,5}{cm}$ en ordonnée).

\textbf{Partie A : étude de la fonction}
\begin{enumerate}
\item Calculer les limites de $f$ aux bornes de son domaine de définition. En déduire l'équation de l'asymptote verticale à $\mathscr{C}$.
\item Montrer que, pour tout $x\neq 2$ : $\;f(x)=x-1+\dfrac{4}{x-2}$.
\item En déduire que la droite $\Delta$ d'équation $y=x-1$ est asymptote oblique à $\mathscr{C}$ en $+\infty$ et en $-\infty$.
\item Étudier, suivant les valeurs de $x$, la position relative de $\mathscr{C}$ et de $\Delta$.
\item Montrer que $f'(x)=\dfrac{x^2-4x}{(x-2)^2}$, puis dresser le tableau de variations de $f$.
\item Montrer que le point $\Omega(2\,;\,1)$ est centre de symétrie de $\mathscr{C}$.
\end{enumerate}

\textbf{Partie B : construction et exploitation}
\begin{enumerate}
\item Tracer les asymptotes, puis la courbe $\mathscr{C}$.
\item Résoudre graphiquement l'équation $f(x)=x$ et donner un encadrement à $10^{-1}$ près de chaque solution.
\item Retrouver algébriquement les solutions exactes de l'équation $f(x)=x$.
\item L'entreprise vend chaque tonne à un prix unitaire modélisé par la droite d'équation $y=x$. À l'aide des questions précédentes, indiquer pour quelles quantités produites ($x>2$) le coût moyen est inférieur au prix de vente.
\end{enumerate}
\end{exercice}

\newpage

\coursec{Corrigés des exercices}

\courssub{Je vérifie mes connaissances}

\begin{corrige}[Exercice 1 — QCM]
\begin{enumerate}
\item \textbf{Réponse b) :} $x=3$. La fonction $f$ n'est pas définie en $3$ et $\displaystyle\lim_{\substack{x\to 3\\ x>3}}f(x)=+\infty$ (numérateur proche de $7>0$, dénominateur proche de $0^+$), donc la droite d'équation $x=3$ est asymptote verticale. La réponse c) $y=2$ est l'asymptote \emph{horizontale}, pas verticale.
\item \textbf{Réponse b) :} une asymptote horizontale d'équation $y=-4$. C'est la définition même : si $\displaystyle\lim_{x\to +\infty}f(x)=\ell$ (fini), la droite $y=\ell$ est asymptote horizontale en $+\infty$.
\item \textbf{Réponse b) :} $\displaystyle\lim_{x\to +\infty}\bigl[f(x)-(ax+b)\bigr]=0$. C'est la caractérisation de l'asymptote oblique : la distance verticale entre la courbe et la droite tend vers $0$.
\item \textbf{Réponse b) :} $f(\alpha+h)+f(\alpha-h)=2\beta$, c'est-à-dire $\dfrac{f(\alpha+h)+f(\alpha-h)}{2}=\beta$ : le point $\Omega(\alpha\,;\,\beta)$ est le milieu des points de la courbe d'abscisses $\alpha-h$ et $\alpha+h$.
\end{enumerate}
\end{corrige}

\begin{corrige}[Exercice 2 — Vrai ou faux]
\begin{enumerate}
\item \textbf{Faux.} Il faut de plus que la limite de $f$ en $a$ (à gauche ou à droite) soit infinie. \emph{Contre-exemple :} la fonction $f$ définie par $f(x)=\dfrac{x^2-1}{x-1}$ n'est pas définie en $1$, mais pour $x\neq 1$, $f(x)=\dfrac{(x-1)(x+1)}{x-1}=x+1$, donc $\displaystyle\lim_{x\to 1}f(x)=2$ (limite finie) : la droite $x=1$ n'est \emph{pas} asymptote, la courbe présente seulement un « trou » au point $(1\,;\,2)$.
\item \textbf{Vrai.} Soit $f(x)=ax^3+bx^2+cx+d$ ($a\neq 0$). Posons $\alpha=-\dfrac{b}{3a}$ et $\beta=f(\alpha)$. Un calcul direct (développer $f(\alpha+h)$ et $f(\alpha-h)$) montre que les termes en $h$ et en $h^3$ s'annulent dans la somme, de sorte que $f(\alpha+h)+f(\alpha-h)=2\beta$ pour tout réel $h$. Donc $\Omega(\alpha\,;\,\beta)$ est centre de symétrie de la courbe. (Voir l'exercice 7 pour un exemple complet mené pas à pas.)
\item \textbf{Vrai.} L'asymptote horizontale décrit le comportement \emph{à l'infini} uniquement ; rien n'interdit à la courbe de la croiser pour des valeurs finies de $x$. \emph{Contre-exemple :} soit $f(x)=\dfrac{x^2+x-1}{x^2}=1+\dfrac{1}{x}-\dfrac{1}{x^2}$, définie sur $\mathbb{R}^*$. On a $\displaystyle\lim_{x\to\pm\infty}f(x)=1$, donc la droite $y=1$ est asymptote horizontale ; pourtant $f(1)=1+1-1=1$, donc la courbe coupe son asymptote au point $(1\,;\,1)$.
\item \textbf{Vrai.} La fonction inverse est \emph{involutive} : si $y=\dfrac{1}{x}$, alors $x=\dfrac{1}{y}$. Échanger les coordonnées d'un point correspond exactement à la symétrie par rapport à la droite $y=x$ (première bissectrice) ; la courbe est donc invariante par cette symétrie. On peut aussi vérifier que le symétrique de $M\left(x\,;\,\tfrac1x\right)$ par rapport à la première bissectrice est $M'\left(\tfrac1x\,;\,x\right)$, qui appartient bien à la courbe puisque $x=\dfrac{1}{1/x}$.
\end{enumerate}
\end{corrige}

\begin{corrige}[Exercice 3 — Lecture de limites et asymptotes]
\begin{enumerate}
\item Des limites infinies en $1$ (à gauche \emph{et} à droite), on déduit l'\textbf{asymptote verticale} d'équation $\boxed{x=1}$. Des limites finies égales à $2$ en $-\infty$ et en $+\infty$, on déduit l'\textbf{asymptote horizontale} d'équation $\boxed{y=2}$, valable des deux côtés.
\item \textbf{Oui, c'est possible.} Les informations données ne concernent que les voisinages de $1$, de $-\infty$ et de $+\infty$. Si la fonction n'était pas définie en un autre point $a$ avec une limite infinie en $a$, sa courbe admettrait une asymptote verticale supplémentaire $x=a$. Les données ne permettent ni de l'affirmer ni de l'exclure.
\item Allure possible : une branche à gauche de $x=1$, partant de l'asymptote $y=2$ en $-\infty$ et plongeant vers $-\infty$ en $1^-$ ; une branche à droite, descendant de $+\infty$ en $1^+$ et se rapprochant de $y=2$ en $+\infty$.
\begin{popfigure}
\begin{center}
\begin{tikzpicture}[scale=0.8]
\draw[popline,->] (-4.5,0) -- (5.5,0) node[below left]{$x$};
\draw[popline,->] (0,-3) -- (0,4.5) node[below left]{$y$};
\draw[poporange,dashed,thick] (1,-3) -- (1,4.5) node[right]{$x=1$};
\draw[popgreen,dashed,thick] (-4.5,2) -- (5.5,2) node[above left]{$y=2$};
\draw[popblue,very thick,domain=-4.2:0.55,samples=80] plot (\x,{2+1/(\x-1)});
\draw[popblue,very thick,domain=1.35:5.2,samples=80] plot (\x,{2+1/(\x-1)});
\node[below left] at (0,0) {$O$};
\end{tikzpicture}
\end{center}
\legende{Une allure compatible avec les quatre limites données.}
\end{popfigure}
\end{enumerate}
\end{corrige}

\begin{corrige}[Exercice 4 — Calculs directs]
\begin{enumerate}
\item \textbf{Domaines de définition.}
\begin{itemize}
\item $f(x)=\dfrac{3x-2}{x+1}$ est définie si et seulement si $x+1\neq 0$, donc $D_f=\mathbb{R}\setminus\{-1\}$.
\item $g(x)=x^3-3x^2+2$ est un polynôme : $D_g=\mathbb{R}$.
\item $h(x)=\dfrac{x^2+2x+2}{x+1}$ : $D_h=\mathbb{R}\setminus\{-1\}$.
\end{itemize}
\item \textbf{Limites de $f$.} En $-1$, le numérateur vaut $3\times(-1)-2=-5<0$.
\begin{itemize}
\item Si $x\to -1$ avec $x<-1$, alors $x+1<0$ (proche de $0^-$) : $\displaystyle\lim_{\substack{x\to -1\\ x<-1}}f(x)=+\infty$.
\item Si $x\to -1$ avec $x>-1$, alors $x+1>0$ : $\displaystyle\lim_{\substack{x\to -1\\ x>-1}}f(x)=-\infty$.
\end{itemize}
En $+\infty$, limite du quotient des termes de plus haut degré : $\displaystyle\lim_{x\to +\infty}f(x)=\lim_{x\to +\infty}\dfrac{3x}{x}=3$.
\item Pour tout $x\neq -1$ :
\[ (x+1)+\frac{1}{x+1}=\frac{(x+1)^2+1}{x+1}=\frac{x^2+2x+1+1}{x+1}=\frac{x^2+2x+2}{x+1}=h(x). \]
\item D'après la question précédente, $h(x)-(x+1)=\dfrac{1}{x+1}$, donc
\[ \lim_{x\to +\infty}\bigl[h(x)-(x+1)\bigr]=\lim_{x\to +\infty}\frac{1}{x+1}=0. \]
\textbf{Interprétation graphique :} la droite d'équation $y=x+1$ est \cle{asymptote oblique} à la courbe de $h$ en $+\infty$ (et aussi en $-\infty$, par le même calcul).
\end{enumerate}
\end{corrige}

\courssub{Je m'entraîne}

\begin{corrige}[Exercice 5 — Asymptotes d'une fonction rationnelle]
\begin{enumerate}
\item $f$ est définie si et seulement si $x-2\neq 0$ : $D_f=\mathbb{R}\setminus\{2\}$.
\item \textbf{Limites aux bornes.}
\begin{itemize}
\item En $+\infty$ : $\displaystyle\lim_{x\to +\infty}f(x)=\lim_{x\to +\infty}\frac{2x^2}{x}=\lim_{x\to +\infty}2x=+\infty$. De même $\displaystyle\lim_{x\to -\infty}f(x)=-\infty$.
\item En $2$ : le numérateur vaut $2\times 4-3\times 2+1=3>0$. Si $x\to 2$ avec $x<2$, $x-2<0$ : $\displaystyle\lim_{\substack{x\to 2\\ x<2}}f(x)=-\infty$ ; si $x>2$ : $\displaystyle\lim_{\substack{x\to 2\\ x>2}}f(x)=+\infty$.
\end{itemize}
La limite étant infinie en $2$, la droite d'équation $\boxed{x=2}$ est \cle{asymptote verticale} à $\mathscr{C}_f$.
\item On procède par division euclidienne de $2x^2-3x+1$ par $x-2$ :
\[ (x-2)(2x+1)=2x^2-3x-2 \quad\Longrightarrow\quad 2x^2-3x+1=(x-2)(2x+1)+3. \]
Donc, pour tout $x\neq 2$ : $f(x)=2x+1+\dfrac{3}{x-2}$, c'est-à-dire $\boxed{a=2,\ b=1,\ c=3}$.
\item Pour tout $x\neq 2$ : $f(x)-(2x+1)=\dfrac{3}{x-2}$. Or $\displaystyle\lim_{x\to +\infty}\frac{3}{x-2}=0$ et $\displaystyle\lim_{x\to -\infty}\frac{3}{x-2}=0$. Donc
\[ \lim_{x\to +\infty}\bigl[f(x)-(2x+1)\bigr]=0 \quad\text{et}\quad \lim_{x\to -\infty}\bigl[f(x)-(2x+1)\bigr]=0, \]
ce qui prouve que la droite $\mathscr{D}$ d'équation $y=2x+1$ est \cle{asymptote oblique} à $\mathscr{C}_f$ en $+\infty$ et en $-\infty$.
\item \textbf{Position relative.} Le signe de $f(x)-(2x+1)=\dfrac{3}{x-2}$ est celui de $x-2$ :
\begin{itemize}
\item si $x>2$, $f(x)-(2x+1)>0$ : $\mathscr{C}_f$ est \textbf{au-dessus} de $\mathscr{D}$ ;
\item si $x<2$, $f(x)-(2x+1)<0$ : $\mathscr{C}_f$ est \textbf{en dessous} de $\mathscr{D}$.
\end{itemize}
\end{enumerate}
\end{corrige}

\begin{corrige}[Exercice 6 — Axe de symétrie d'une parabole]
\begin{enumerate}
\item Pour tout réel $h$ :
\begin{align*}
g(2+h)&=(2+h)^2-4(2+h)+5=4+4h+h^2-8-4h+5=h^2+1,\\
g(2-h)&=(2-h)^2-4(2-h)+5=4-4h+h^2-8+4h+5=h^2+1.
\end{align*}
Donc $g(2+h)=g(2-h)$ pour tout réel $h$ : la droite d'équation $\boxed{x=2}$ est \cle{axe de symétrie} de $\mathscr{P}$.
\item On a la forme canonique $g(x)=(x-2)^2+1$. Comme $(x-2)^2\geq 0$, $g$ admet un minimum en $x=2$, égal à $g(2)=1$. De plus $g'(x)=2x-4$, négative pour $x<2$, positive pour $x>2$.
\begin{center}
\begin{tabular}{|c|ccccc|}\hline
$x$ & $-\infty$ & & $2$ & & $+\infty$ \\ \hline
$g'(x)$ & & $-$ & $0$ & $+$ & \\ \hline
$g$ & $+\infty$ & $\searrow$ & $1$ & $\nearrow$ & $+\infty$ \\ \hline
\end{tabular}
\end{center}
\begin{popfigure}
\begin{center}
\begin{tikzpicture}[scale=0.75]
\draw[popline,->] (-1.5,0) -- (5.5,0) node[below left]{$x$};
\draw[popline,->] (0,-0.5) -- (0,6) node[below left]{$y$};
\draw[poporange,dashed] (2,-0.5) -- (2,6) node[above]{$x=2$};
\draw[popblue,very thick,domain=-0.2:4.2,samples=100] plot (\x,{(\x-2)*(\x-2)+1});
\draw[popgreen,dashed] (-1.5,2) -- (5.5,2) node[above left]{$y=2$};
\fill[popink] (2,1) circle (2pt) node[below right]{$S(2\,;\,1)$};
\fill[popink] (1,2) circle (2pt); \fill[popink] (3,2) circle (2pt);
\foreach \x in {1,2,3,4} \draw (\x,0.08) -- (\x,-0.08) node[below]{\small $\x$};
\foreach \y in {1,2} \draw (0.08,\y) -- (-0.08,\y) node[left]{\small $\y$};
\node[below left] at (0,0) {\small $O$};
\end{tikzpicture}
\end{center}
\legende{La parabole $\mathscr{P}$, son axe de symétrie et la droite $y=2$.}
\end{popfigure}
\item \textbf{Graphiquement :} $g(x)\leq 2$ lorsque $\mathscr{P}$ est sous la droite $y=2$, qui coupe $\mathscr{P}$ aux points d'abscisses $1$ et $3$ : $S=[1\,;\,3]$.
\textbf{Par le calcul :} $g(x)\leq 2 \iff (x-2)^2+1\leq 2 \iff (x-2)^2\leq 1 \iff -1\leq x-2\leq 1 \iff 1\leq x\leq 3$. On retrouve bien $S=[1\,;\,3]$.
\end{enumerate}
\end{corrige}

\begin{corrige}[Exercice 7 — Centre de symétrie d'une cubique]
\begin{enumerate}
\item Pour tout réel $h$ :
\begin{align*}
f(1+h)&=(1+h)^3-3(1+h)^2+2=(1+3h+3h^2+h^3)-3(1+2h+h^2)+2=h^3-3h,\\
f(1-h)&=(1-h)^3-3(1-h)^2+2=(1-3h+3h^2-h^3)-3(1-2h+h^2)+2=-h^3+3h.
\end{align*}
Donc $f(1+h)+f(1-h)=(h^3-3h)+(-h^3+3h)=0$.
\item On a montré que pour tout réel $h$ : $f(1+h)+f(1-h)=0=2\times 0$. C'est la condition caractéristique avec $\alpha=1$ et $\beta=0$ : le point $\boxed{\Omega(1\,;\,0)}$ est \cle{centre de symétrie} de $\mathscr{C}$.
\item $f'(x)=3x^2-6x=3x(x-2)$. Ce trinôme est positif à l'extérieur des racines $0$ et $2$, négatif entre les racines. Valeurs : $f(0)=2$ et $f(2)=8-12+2=-2$.
\begin{center}
\begin{tabular}{|c|ccccccc|}\hline
$x$ & $-\infty$ & & $0$ & & $2$ & & $+\infty$ \\ \hline
$f'(x)$ & & $+$ & $0$ & $-$ & $0$ & $+$ & \\ \hline
$f$ & $-\infty$ & $\nearrow$ & $2$ & $\searrow$ & $-2$ & $\nearrow$ & $+\infty$ \\ \hline
\end{tabular}
\end{center}
\item $f'(1)=3-6=-3$ et $f(1)=0$. La tangente en $\Omega$ a pour équation :
\[ y=f'(1)(x-1)+f(1)=-3(x-1) \quad\text{soit}\quad \boxed{y=-3x+3}. \]
\textbf{Constat :} étudions la position de la courbe par rapport à cette tangente :
\[ f(x)-(-3x+3)=x^3-3x^2+2+3x-3=x^3-3x^2+3x-1=(x-1)^3, \]
qui est négatif pour $x<1$ et positif pour $x>1$ : la tangente au centre de symétrie \emph{traverse} la courbe. On dit que $\Omega$ est un \cle{point d'inflexion}. C'est un fait général : le centre de symétrie de la courbe d'une fonction polynôme de degré $3$ est toujours son point d'inflexion.
\end{enumerate}
\end{corrige}

\begin{corrige}[Exercice 8 — Exploitation graphique]
\emph{Les réponses ci-dessous correspondent à la courbe fournie ; les valeurs lues sont des valeurs approchées, justifiées par des pointillés de lecture sur le graphique reproduit.}
\begin{enumerate}
\item Reproduction : on trace le repère orthonormé (unité $\SI{1}{cm}$), l'asymptote verticale $x=-1$ et l'asymptote horizontale $y=1$ en pointillés, puis on reporte les points remarquables de la courbe avant de la tracer à main levée.
\item \textbf{Lectures graphiques} (on trace les droites horizontales $y=0$, $y=1$, $y=2$ et on projette les intersections sur l'axe des abscisses) :
\begin{itemize}
\item[a)] $f(x)=0$ : la courbe coupe l'axe des abscisses en deux points : $S=\{-3\;;\;0\}$.
\item[b)] $f(x)=2$ : la droite $y=2$ coupe la courbe en trois points : $S=\{-2\;;\;1\;;\;3\}$.
\item[c)] $f(x)>1$ : la courbe est strictement au-dessus de la droite $y=1$ pour $S=\left]-2{,}5\,;\,-1\right[\cup\left]0{,}5\,;\,6\right]$.
\item[d)] $f(x)\leq 0$ : la courbe est sous l'axe des abscisses (ou le touche) pour $S=[-4\,;\,-3]\cup\left]-1\,;\,0\right]$. \textbf{Attention :} $-1$ est exclu car $f$ n'y est pas définie (asymptote verticale).
\end{itemize}
\item \textbf{Discussion suivant $m$.} On compte les points d'intersection de $\mathscr{C}$ avec la droite horizontale $y=m$ :
\begin{center}
\begin{tabularx}{\linewidth}{|l|X|}\hline
\rowcolor{popblueL} Valeurs de $m$ & Nombre de solutions \\ \hline
$m>3$ & $1$ solution \\ \hline
$m=3$ & $2$ solutions \\ \hline
$1<m<3$ & $3$ solutions \\ \hline
$-2\leq m\leq 1$ & $2$ solutions \\ \hline
$m<-2$ & $1$ solution \\ \hline
\end{tabularx}
\end{center}
(La branche gauche, croissante de $-2$ vers $+\infty$, fournit une solution dès que $m\geq -2$ ; la branche droite, qui monte de $-\infty$ au maximum $3$ puis redescend vers $1$, fournit deux solutions si $1<m<3$, une si $m\leq 1$ ou $m=3$, aucune si $m>3$.)
\end{enumerate}
\end{corrige}

\courssub{Je me prépare au Probatoire}

\begin{corrige}[Exercice 9 — Étude complète d'une fonction homographique]
\textbf{Partie A}
\begin{enumerate}
\item $f$ est définie si et seulement si $x+1\neq 0$ : $\boxed{D_f=\mathbb{R}\setminus\{-1\}}$.
\item \textbf{Limites.} En $+\infty$ et $-\infty$ : $\displaystyle\lim_{x\to\pm\infty}f(x)=\lim_{x\to\pm\infty}\frac{3x}{x}=3$.
En $-1$ : le numérateur vaut $3\times(-1)-2=-5<0$.
\begin{itemize}
\item Si $x\to -1$ avec $x<-1$, $x+1<0$ : $\displaystyle\lim_{\substack{x\to -1\\ x<-1}}f(x)=+\infty$.
\item Si $x\to -1$ avec $x>-1$, $x+1>0$ : $\displaystyle\lim_{\substack{x\to -1\\ x>-1}}f(x)=-\infty$.
\end{itemize}
\textbf{Asymptotes :} la droite $\boxed{x=-1}$ est asymptote verticale (limites infinies en $-1$) et la droite $\boxed{y=3}$ est asymptote horizontale en $-\infty$ et en $+\infty$.
\item $f$ est de la forme $\dfrac{u}{v}$ avec $u(x)=3x-2$, $u'(x)=3$, $v(x)=x+1$, $v'(x)=1$. Pour tout $x\in D_f$ :
\[ f'(x)=\frac{u'v-uv'}{v^2}=\frac{3(x+1)-(3x-2)\times 1}{(x+1)^2}=\frac{3x+3-3x+2}{(x+1)^2}=\frac{5}{(x+1)^2}. \]
\item Pour tout $x\in D_f$, $(x+1)^2>0$ donc $f'(x)=\dfrac{5}{(x+1)^2}>0$ : $f$ est \textbf{strictement croissante sur chacun des intervalles} $\left]-\infty\,;\,-1\right[$ et $\left]-1\,;\,+\infty\right[$.
\begin{center}
\begin{tabular}{|c|ccccc|}\hline
$x$ & $-\infty$ & & $-1$ & & $+\infty$ \\ \hline
$f'(x)$ & & $+$ & & $+$ & \\ \hline
$f$ & $3$ & $\nearrow$ & $+\infty$\ \textbf{||}\ $-\infty$ & $\nearrow$ & $3$ \\ \hline
\end{tabular}
\end{center}
\begin{attention}[Point de vigilance]
Il est interdit d'écrire « $f$ est croissante sur $D_f$ » : la croissance est établie \emph{sur chaque intervalle} séparément, car $f$ n'est pas définie en $-1$. La double barre dans le tableau matérialise la valeur interdite.
\end{attention}
\item Pour tout $h\neq 0$, les réels $-1-h$ et $-1+h$ appartiennent à $D_f$ et :
\[ f(-1+h)=\frac{3(-1+h)-2}{h}=\frac{3h-5}{h}=3-\frac{5}{h}\;;\qquad f(-1-h)=\frac{3(-1-h)-2}{-h}=\frac{-3h-5}{-h}=3+\frac{5}{h}. \]
Donc $f(-1+h)+f(-1-h)=6=2\times 3$ : le point $\boxed{\Omega(-1\,;\,3)}$ est \cle{centre de symétrie} de $\mathscr{C}_f$. (On remarque que $\Omega$ est l'intersection des deux asymptotes.)
\item \textbf{Intersections avec les axes.}
\begin{itemize}
\item Axe des ordonnées : $f(0)=\dfrac{-2}{1}=-2$ : point $A(0\,;\,-2)$.
\item Axe des abscisses : $f(x)=0\iff 3x-2=0\iff x=\dfrac{2}{3}$ : point $B\left(\dfrac{2}{3}\,;\,0\right)$.
\end{itemize}
\end{enumerate}

\textbf{Partie B}
\begin{enumerate}
\item Construction : on trace en pointillés les asymptotes $x=-1$ et $y=3$, on place $\Omega$, $A$ et $B$, puis on trace les deux branches d'hyperbole, symétriques par rapport à $\Omega$.
\begin{popfigure}
\begin{center}
\begin{tikzpicture}[scale=0.62]
\draw[popline,->] (-6.5,0) -- (5.5,0) node[below left]{$x$};
\draw[popline,->] (0,-5) -- (0,7) node[below left]{$y$};
\draw[poporange,dashed,thick] (-1,-5) -- (-1,7) node[above]{$x=-1$};
\draw[poporange,dashed,thick] (-6.5,3) -- (5.5,3) node[above left]{$y=3$};
\draw[popgreen,thick,domain=-6.5:5.2,samples=60] plot (\x,{\x+4}) node[above left]{$\mathscr{D}$};
\draw[popblue,very thick,domain=-6.5:-1.85,samples=80] plot (\x,{(3*\x-2)/(\x+1)});
\draw[popblue,very thick,domain=-0.35:5.2,samples=80] plot (\x,{(3*\x-2)/(\x+1)});
\fill[popink] (-1,3) circle (2pt) node[above right]{$\Omega$};
\fill[popink] (0,-2) circle (2pt) node[below left]{$A$};
\fill[popink] (0.667,0) circle (2pt) node[above right]{$B$};
\node[below left] at (0,0) {\small $O$};
\foreach \x in {-5,-4,-3,-2,1,2,3,4,5} \draw (\x,0.07) -- (\x,-0.07) node[below]{\scriptsize $\x$};
\foreach \y in {-4,-2,2,4,6} \draw (0.07,\y) -- (-0.07,\y) node[left]{\scriptsize $\y$};
\end{tikzpicture}
\end{center}
\legende{La courbe $\mathscr{C}_f$, ses asymptotes et la droite $\mathscr{D}$ : $y=x+4$.}
\end{popfigure}
\item La droite $\mathscr{D}$ passe par $(0\,;\,4)$ et $(-4\,;\,0)$. \textbf{Graphiquement}, $\mathscr{D}$ ne rencontre pas $\mathscr{C}_f$ : l'équation $f(x)=x+4$ n'a \textbf{aucune solution}.
\textbf{Vérification algébrique :} pour $x\neq -1$,
\[ f(x)=x+4 \iff 3x-2=(x+4)(x+1) \iff 3x-2=x^2+5x+4 \iff x^2+2x+6=0. \]
Discriminant : $\Delta=4-24=-20<0$ (ou directement $x^2+2x+6=(x+1)^2+5>0$ pour tout $x$) : pas de solution réelle. $S=\varnothing$. Le calcul confirme la lecture graphique.
\item La facture dépasse $\num{2500}$ FCFA, c'est-à-dire $2{,}5$ milliers de FCFA, lorsque $f(x)>2{,}5$. Par le calcul, pour $x\neq -1$ :
\[ f(x)=2{,}5 \iff 3x-2=2{,}5(x+1) \iff 0{,}5x=4{,}5 \iff x=9. \]
Comme $f$ est strictement croissante sur $\left]-1\,;\,+\infty\right[$, on a $f(x)>2{,}5 \iff x>9$ sur cet intervalle. La facture dépasse donc $\num{2500}$ FCFA \textbf{à partir de $9$ gigaoctets consommés} (strictement au-delà de $9$ Go).
\end{enumerate}

\textbf{Barème indicatif (sur 10) :} A1 : $0{,}5$ ; A2 : $1{,}5$ ; A3 : $1$ ; A4 : $1$ ; A5 : $1$ ; A6 : $1$ ; B1 : $1{,}5$ ; B2 : $0{,}5$ ; B3 : $1$ ; B4 : $1$.

\textbf{Points de vigilance :} citer la condition $x\neq -1$ avant chaque calcul de limite ou d'équivalence ; justifier l'asymptote par la \emph{limite} et non par la seule valeur interdite ; pour le centre de symétrie, écrire explicitement la relation $f(\alpha+h)+f(\alpha-h)=2\beta$ ; ne pas conclure « croissante sur $D_f$ ».
\end{corrige}

\begin{corrige}[Exercice 10 — Cubique et discussion suivant un paramètre]
\textbf{Partie A}
\begin{enumerate}
\item $f$ est un polynôme : sa limite en l'infini est celle de son terme de plus haut degré $x^3$ :
\[ \lim_{x\to -\infty}f(x)=-\infty \qquad\text{et}\qquad \lim_{x\to +\infty}f(x)=+\infty. \]
\item $f'(x)=3x^2-6x=3x(x-2)$ : positif sur $\left]-\infty\,;\,0\right[\cup\left]2\,;\,+\infty\right[$, négatif sur $\left]0\,;\,2\right[$. Avec $f(0)=2$ et $f(2)=-2$ :
\begin{center}
\begin{tabular}{|c|ccccccc|}\hline
$x$ & $-\infty$ & & $0$ & & $2$ & & $+\infty$ \\ \hline
$f'(x)$ & & $+$ & $0$ & $-$ & $0$ & $+$ & \\ \hline
$f$ & $-\infty$ & $\nearrow$ & $2$ & $\searrow$ & $-2$ & $\nearrow$ & $+\infty$ \\ \hline
\end{tabular}
\end{center}
\item Voir l'exercice 7 : pour tout réel $h$, $f(1+h)=h^3-3h$ et $f(1-h)=-h^3+3h$, donc $f(1+h)+f(1-h)=0=2\times 0$ : $\boxed{\Omega(1\,;\,0)}$ est centre de symétrie de $\mathscr{C}$.
\item $f'(1)=-3$ et $f(1)=0$ : $T : y=-3(x-1)+0$, soit $\boxed{T : y=-3x+3}$.
\item Sur $[2\,;\,3]$, $f$ est \textbf{continue} (polynôme) et \textbf{strictement croissante} (car $f'>0$ sur $\left]2\,;\,3\right[$). De plus $f(2)=-2<0$ et $f(3)=27-27+2=2>0$, donc $0\in\left[f(2)\,;\,f(3)\right]$. D'après le \textbf{théorème des valeurs intermédiaires} (cas des fonctions strictement monotones), l'équation $f(x)=0$ admet une \textbf{unique} solution $\alpha$ dans $[2\,;\,3]$.
\textbf{Encadrement :} $f(2{,}7)=19{,}683-21{,}87+2=-0{,}187<0$ et $f(2{,}8)=21{,}952-23{,}52+2=0{,}432>0$. Donc $\boxed{2{,}7<\alpha<2{,}8}$.
\end{enumerate}

\textbf{Partie B}
\begin{enumerate}
\item On place les extrema $(0\,;\,2)$ et $(2\,;\,-2)$, le centre $\Omega(1\,;\,0)$, on trace $T$ (qui passe par $\Omega$ et $(0\,;\,3)$), puis $\mathscr{C}$ en exploitant la symétrie de centre $\Omega$ : le symétrique de $(0\,;\,2)$ est $(2\,;\,-2)$, celui de $(-1\,;\,-2)$ est $(3\,;\,2)$.
\begin{popfigure}
\begin{center}
\begin{tikzpicture}[scale=0.9]
\draw[popline,->] (-2.2,0) -- (3.8,0) node[below left]{$x$};
\draw[popline,->] (0,-3) -- (0,3.8) node[below left]{$y$};
\draw[popgreen,thick,domain=-0.4:2.1,samples=20] plot (\x,{-3*\x+3}) node[right]{$T$};
\draw[popblue,very thick,domain=-1.35:3.35,samples=120] plot (\x,{\x*\x*\x-3*\x*\x+2});
\fill[popink] (1,0) circle (2pt) node[below left]{$\Omega$};
\fill[popink] (0,2) circle (2pt) node[above right]{\scriptsize $(0\,;\,2)$};
\fill[popink] (2,-2) circle (2pt) node[below right]{\scriptsize $(2\,;\,-2)$};
\node[below left] at (0,0) {\small $O$};
\foreach \x in {-1,1,2,3} \draw (\x,0.06) -- (\x,-0.06) node[below]{\scriptsize $\x$};
\foreach \y in {-2,2} \draw (0.06,\y) -- (-0.06,\y) node[left]{\scriptsize $\y$};
\end{tikzpicture}
\end{center}
\legende{La cubique $\mathscr{C}$, sa tangente au centre de symétrie $\Omega$.}
\end{popfigure}
\item L'équation $f(x)=m$ revient à chercher les intersections de $\mathscr{C}$ avec la droite horizontale $y=m$. D'après le tableau de variations (maximum local $2$, minimum local $-2$) :
\begin{center}
\begin{tabularx}{\linewidth}{|l|X|}\hline
\rowcolor{popblueL} Valeurs de $m$ & Nombre de solutions \\ \hline
$m<-2$ ou $m>2$ & $1$ solution \\ \hline
$m=-2$ ou $m=2$ & $2$ solutions (dont une double) \\ \hline
$-2<m<2$ & $3$ solutions \\ \hline
\end{tabularx}
\end{center}
\item Le bénéfice est positif lorsque $f(x)>0$, c'est-à-dire lorsque $\mathscr{C}$ est au-dessus de l'axe des abscisses. Factorisons : $f(1)=0$, donc $x-1$ se met en facteur et l'on obtient $f(x)=(x-1)(x^2-2x-2)$ ; le trinôme $x^2-2x-2$ a pour discriminant $\Delta=4+8=12$ et pour racines $1-\sqrt3$ et $1+\sqrt3\approx 2{,}73$. Sur $[0\,;\,3]$, la courbe est strictement au-dessus de l'axe pour
\[ x\in[0\,;\,1[\;\cup\;\left]1+\sqrt3\,;\,3\right]. \]
\textbf{Conclusion concrète :} le bénéfice est positif pour une vente strictement inférieure à $1$ tonne, ou strictement supérieure à environ $2{,}73$ tonnes (jusqu'à $3$ tonnes). Entre $1$ et $2{,}73$ tonnes, la coopérative est en déficit.
\end{enumerate}

\textbf{Barème indicatif (sur 10) :} A1 : $0{,}5$ ; A2 : $1{,}5$ ; A3 : $1$ ; A4 : $0{,}5$ ; A5 : $1{,}5$ ; B1 : $1{,}5$ ; B2 : $1{,}5$ ; B3 : $1{,}5$ ; présentation : $0{,}5$.

\textbf{Points de vigilance :} pour le TVI, citer les trois hypothèses (continuité, stricte monotonie, $0$ entre les images) ; pour la discussion en $m$, s'appuyer sur les extrema locaux et préciser les cas d'égalité ; ne pas oublier la contrainte $0\leq x\leq 3$ dans la question concrète.
\end{corrige}

\begin{corrige}[Exercice 11 — Fonction rationnelle et asymptote oblique]
\textbf{Partie A}
\begin{enumerate}
\item \textbf{Limites aux bornes de $D_f=\mathbb{R}\setminus\{2\}$.}
\begin{itemize}
\item En $\pm\infty$ : $\displaystyle\lim_{x\to\pm\infty}f(x)=\lim_{x\to\pm\infty}\frac{x^2}{x}=\lim_{x\to\pm\infty}x$, donc $\displaystyle\lim_{x\to +\infty}f(x)=+\infty$ et $\displaystyle\lim_{x\to -\infty}f(x)=-\infty$.
\item En $2$ : le numérateur vaut $4-6+6=4>0$. Si $x\to 2$ avec $x<2$ : $x-2<0$, donc $\displaystyle\lim_{\substack{x\to 2\\ x<2}}f(x)=-\infty$ ; si $x>2$ : $\displaystyle\lim_{\substack{x\to 2\\ x>2}}f(x)=+\infty$.
\end{itemize}
La limite étant infinie en $2$, la droite d'équation $\boxed{x=2}$ est \cle{asymptote verticale} à $\mathscr{C}$.
\item Pour tout $x\neq 2$ :
\[ (x-1)+\frac{4}{x-2}=\frac{(x-1)(x-2)+4}{x-2}=\frac{x^2-3x+2+4}{x-2}=\frac{x^2-3x+6}{x-2}=f(x). \]
Donc $\boxed{f(x)=x-1+\dfrac{4}{x-2}}$.
\item D'après A2 : $f(x)-(x-1)=\dfrac{4}{x-2}$, et $\displaystyle\lim_{x\to +\infty}\frac{4}{x-2}=0$, de même en $-\infty$. Donc
\[ \lim_{x\to \pm\infty}\bigl[f(x)-(x-1)\bigr]=0, \]
ce qui prouve que la droite $\Delta$ d'équation $\boxed{y=x-1}$ est \cle{asymptote oblique} à $\mathscr{C}$ en $+\infty$ et en $-\infty$.
\item \textbf{Position relative.} Le signe de $f(x)-(x-1)=\dfrac{4}{x-2}$ est celui de $x-2$ :
\begin{itemize}
\item si $x>2$ : $\mathscr{C}$ est strictement \textbf{au-dessus} de $\Delta$ ;
\item si $x<2$ : $\mathscr{C}$ est strictement \textbf{en dessous} de $\Delta$.
\end{itemize}
\item En dérivant la forme de A2 : $f'(x)=1-\dfrac{4}{(x-2)^2}=\dfrac{(x-2)^2-4}{(x-2)^2}=\dfrac{x^2-4x+4-4}{(x-2)^2}=\dfrac{x^2-4x}{(x-2)^2}$.
Signe de $f'(x)$ : celui de $x(x-4)$ (le dénominateur est strictement positif sur $D_f$), donc positif sur $\left]-\infty\,;\,0\right[\cup\left]4\,;\,+\infty\right[$, négatif sur $\left]0\,;\,2\right[\cup\left]2\,;\,4\right[$.
Valeurs : $f(0)=\dfrac{6}{-2}=-3$ et $f(4)=\dfrac{16-12+6}{2}=5$.
\begin{center}
\begin{tabular}{|c|ccccccccc|}\hline
$x$ & $-\infty$ & & $0$ & & $2$ & & $4$ & & $+\infty$ \\ \hline
$f'(x)$ & & $+$ & $0$ & $-$ & & $-$ & $0$ & $+$ & \\ \hline
$f$ & $-\infty$ & $\nearrow$ & $-3$ & $\searrow$ & $-\infty$\ \textbf{||}\ $+\infty$ & $\searrow$ & $5$ & $\nearrow$ & $+\infty$ \\ \hline
\end{tabular}
\end{center}
\item Pour tout $h\neq 0$ :
\[ f(2+h)=(2+h)-1+\frac{4}{h}=1+h+\frac{4}{h}\;;\qquad f(2-h)=(2-h)-1-\frac{4}{h}=1-h-\frac{4}{h}. \]
Donc $f(2+h)+f(2-h)=2=2\times 1$ : le point $\boxed{\Omega(2\,;\,1)}$ est \cle{centre de symétrie} de $\mathscr{C}$. (C'est l'intersection des asymptotes $x=2$ et $y=x-1$.)
\end{enumerate}

\textbf{Partie B}
\begin{enumerate}
\item On trace en pointillés les asymptotes $x=2$ et $y=x-1$, on place les extrema $(0\,;\,-3)$ et $(4\,;\,5)$, symétriques par rapport à $\Omega$, puis les deux branches de $\mathscr{C}$.
\begin{popfigure}
\begin{center}
\begin{tikzpicture}[x=0.55cm,y=0.32cm]
\draw[popline,->] (-5.5,0) -- (8.5,0) node[below left]{$x$};
\draw[popline,->] (0,-9) -- (0,10) node[below left]{$y$};
\draw[poporange,dashed,thick] (2,-9) -- (2,10) node[above]{$x=2$};
\draw[poporange,dashed,thick,domain=-5.5:8.2] plot (\x,{\x-1}) node[above left]{$\Delta$};
\draw[popgreen,thick,domain=-5.5:8.2] plot (\x,{\x}) node[below left]{$y=x$};
\draw[popblue,very thick,domain=-5.5:1.5,samples=90] plot (\x,{\x-1+4/(\x-2)});
\draw[popblue,very thick,domain=2.45:8.2,samples=90] plot (\x,{\x-1+4/(\x-2)});
\fill[popink] (2,1) circle (2.5pt) node[above right]{$\Omega$};
\fill[popink] (0,-3) circle (2.5pt) node[below left]{\scriptsize $(0\,;\,-3)$};
\fill[popink] (4,5) circle (2.5pt) node[above left]{\scriptsize $(4\,;\,5)$};
\fill[poppink] (6,6) circle (2.5pt) node[below right]{\scriptsize $6$};
\node[below left] at (0,0) {\small $O$};
\end{tikzpicture}
\end{center}
\legende{La courbe $\mathscr{C}$, ses asymptotes et la droite $y=x$.}
\end{popfigure}
\item \textbf{Graphiquement}, la droite $y=x$ ne rencontre $\mathscr{C}$ qu'en un seul point, d'abscisse voisine de $6$ (branche de droite ; la branche de gauche reste sous $\Delta$, elle-même sous $y=x$). L'équation $f(x)=x$ admet une unique solution, voisine de $\boxed{6}$.
\item \textbf{Algébriquement}, pour $x\neq 2$ :
\[ f(x)=x \iff \frac{x^2-3x+6}{x-2}=x \iff x^2-3x+6=x^2-2x \iff -x+6=0 \iff x=6. \]
L'unique solution exacte est $\boxed{x=6}$, ce qui confirme la lecture graphique.
\item Le coût moyen est inférieur au prix de vente lorsque $f(x)<x$ avec $x>2$. D'après B3, $f(x)-x=\dfrac{6-x}{x-2}$ ; pour $x>2$, le dénominateur est positif, donc cette quantité est négative si et seulement si $x>6$.
\textbf{Conclusion :} l'entreprise réalise une marge positive (coût moyen inférieur au prix de vente) lorsqu'elle produit \textbf{plus de $6$ tonnes} de cacao.
\end{enumerate}

\textbf{Barème indicatif (sur 12) :} A1 : $1{,}5$ ; A2 : $1$ ; A3 : $1$ ; A4 : $1$ ; A5 : $2$ ; A6 : $1$ ; B1 : $1{,}5$ ; B2 : $1$ ; B3 : $1$ ; B4 : $1$.

\textbf{Points de vigilance :} pour l'asymptote oblique, il faut impérativement écrire la limite de la \emph{différence} $f(x)-(ax+b)$ et conclure par la phrase de définition ; la position relative se déduit du \emph{signe} de cette même différence ; dans le tableau de variations, ne pas oublier la double barre en $x=2$ et les deux limites infinies de part et d'autre ; dans la question concrète, respecter la contrainte du contexte ($x>2$).
\end{corrige}

\newpage

\coursec{Fiche bilan}

\begin{popfigure}
\begin{tikzpicture}[mindmap, grow cyclic, every node/.style={concept, text=white, font=\small\bfseries},
root concept/.append style={concept color=popdark, font=\bfseries},
level 1/.append style={level distance=3.4cm, sibling angle=60},
level 2/.append style={level distance=2.1cm, sibling angle=45, font=\scriptsize}]
\node[root concept]{Étude et représentation graphique d'une fonction}
child[concept color=popblue]{ node {Asymptotes $\parallel$ axes}
	child[concept color=popblue!70]{ node {verticale : $x=a$ si $\lim_{a}f=\pm\infty$} }
	child[concept color=popblue!70]{ node {horizontale : $y=b$ si $\lim_{\pm\infty}f=b$} }
}
child[concept color=poporange]{ node {Asymptote oblique}
	child[concept color=poporange!70]{ node {$y=ax+b$ si $\lim_{\pm\infty}[f(x)-(ax+b)]=0$} }
}
child[concept color=popgreen]{ node {Symétries}
	child[concept color=popgreen!70]{ node {centre $\Omega(a;b)$ : $f(a+h)+f(a-h)=2b$} }
	child[concept color=popgreen!70]{ node {axe $x=a$ : $f(a+h)=f(a-h)$} }
}
child[concept color=poppurple]{ node {Homographique $x\mapsto\frac{ax+b}{cx+d}$}
	child[concept color=poppurple!70]{ node {hyperbole, centre $\Omega\left(-\frac{d}{c};\frac{a}{c}\right)$} }
}
child[concept color=poppink]{ node {Polynômes degré $\leq 3$}
	child[concept color=poppink!70]{ node {dérivée, TV, tangentes, points remarquables} }
}
child[concept color=popgold]{ node {$x\mapsto\frac{ax^2+bx+c}{dx+e}$}
	child[concept color=popgold!70]{ node {asymptote oblique par division} }
	child[concept color=popgold!70]{ node {lecture graphique : équations, inéquations} }
};
\end{tikzpicture}
\legende{Carte mentale de la leçon : les six compétences clés autour de l'étude d'une fonction.}
\end{popfigure}

\begin{bilan}
\textbf{1. Définitions clés}
\begin{itemize}
\item \cle{Asymptote verticale} : la droite $x=a$ est asymptote à $\mathcal{C}_f$ si $\displaystyle\lim_{x\to a}f(x)=\pm\infty$ (souvent une valeur interdite).
\item \cle{Asymptote horizontale} : la droite $y=b$ est asymptote si $\displaystyle\lim_{x\to\pm\infty}f(x)=b$ ($b$ réel).
\item \cle{Asymptote oblique} : la droite $y=ax+b$ ($a\neq 0$) est asymptote si $\displaystyle\lim_{x\to\pm\infty}\big[f(x)-(ax+b)\big]=0$.
\item \cle{Centre de symétrie} $\Omega(a;b)$ : pour tout $h$ tel que $a\pm h\in D_f$, $f(a+h)+f(a-h)=2b$.
\item \cle{Axe de symétrie} $x=a$ : pour tout $h$, $f(a+h)=f(a-h)$ (avec $D_f$ symétrique par rapport à $a$).
\end{itemize}

\textbf{2. Formules et résultats à connaître par cœur}
\begin{center}
\begin{tabularx}{\linewidth}{|l|X|X|}\hline
\rowcolor{popblueL} \textbf{Fonction} & \textbf{Éléments remarquables} & \textbf{Courbe} \\ \hline
$f(x)=\dfrac{ax+b}{cx+d}$ & $D_f=\mathbb{R}\setminus\left\{-\frac{d}{c}\right\}$ ; asymptotes $x=-\frac{d}{c}$ et $y=\frac{a}{c}$ ; centre $\Omega\left(-\frac{d}{c}\,;\,\frac{a}{c}\right)$ & Hyperbole \\ \hline
$f(x)=ax^3+bx^2+cx+d$ & $f'(x)=3ax^2+2bx+c$ ; signe de $f'$ via $\Delta$ ; centre de symétrie au point d'inflexion $x_0=-\frac{b}{3a}$ & Cubique \\ \hline
$f(x)=\dfrac{ax^2+bx+c}{dx+e}$ & Division : $f(x)=\alpha x+\beta+\dfrac{\gamma}{dx+e}$ ; asymptote oblique $y=\alpha x+\beta$ ; asymptote verticale $x=-\frac{e}{d}$ & Courbe avec asymptote oblique \\ \hline
\end{tabularx}
\end{center}

\textbf{3. Méthodes express}
\begin{itemize}
\item \cle{Plan d'étude type} : (1) ensemble de définition ; (2) limites aux bornes et asymptotes ; (3) dérivée et signe ; (4) tableau de variations complet ; (5) éléments de symétrie ; (6) points remarquables (intersections avec les axes, tangente) ; (7) tracé soigné.
\item \cle{Prouver une asymptote oblique} $y=ax+b$ : calculer $f(x)-(ax+b)$, montrer que la limite vaut $0$ en $\pm\infty$.
\item \cle{Prouver un centre de symétrie} $\Omega(a;b)$ : vérifier $f(a+h)+f(a-h)=2b$, ou poser $X=x-a$, $Y=y-b$ et montrer que la nouvelle fonction est impaire.
\item \cle{Lecture graphique} : $f(x)=k$ se résout par les abscisses des intersections de $\mathcal{C}_f$ avec la droite $y=k$ ; $f(x)\geq k$ : abscisses des points de $\mathcal{C}_f$ situés au-dessus de cette droite.
\end{itemize}
\end{bilan}

\begin{aretenir}[Checklist avant le Probatoire]
\begin{itemize}
\item[$\square$] Je sais déterminer l'ensemble de définition d'une fonction rationnelle (valeurs interdites).
\item[$\square$] Je sais repérer une asymptote verticale à partir d'une limite infinie en une valeur interdite.
\item[$\square$] Je sais justifier une asymptote horizontale par une limite finie en $\pm\infty$.
\item[$\square$] Je sais démontrer qu'une droite $y=ax+b$ est asymptote oblique en calculant $\lim\big[f(x)-(ax+b)\big]=0$.
\item[$\square$] Je sais effectuer la division de $ax^2+bx+c$ par $dx+e$ pour obtenir l'asymptote oblique.
\item[$\square$] Je connais par cœur le centre de symétrie $\Omega\left(-\frac{d}{c}\,;\,\frac{a}{c}\right)$ de l'hyperbole homographique.
\item[$\square$] Je sais prouver un centre de symétrie par la relation $f(a+h)+f(a-h)=2b$.
\item[$\square$] Je sais prouver un axe de symétrie par la relation $f(a+h)=f(a-h)$.
\item[$\square$] Je sais dresser un tableau de variations complet : signe de $f'$, flèches, limites et valeurs aux extrema.
\item[$\square$] Je sais placer asymptotes, points remarquables et tangentes avant de tracer la courbe.
\item[$\square$] Je sais résoudre graphiquement $f(x)=k$, $f(x)\geq k$ et comparer $f$ à sa droite asymptote (position relative).
\item[$\square$] Je rédige proprement : limites écrites, conditions vérifiées, conclusions en français.
\end{itemize}
\end{aretenir}

\findecours

\end{document}
